\subsection{The weighted virtual matrix space}%
\label{subsec:weighted-virtual-matrix-space}

\begin{theorem}[Weighted virtual inner product]%
    \label{thm:weighted-virtual-inner-product}
    \lean{Matrix.rhoWeighted_inner}
    \leanok
    Let $\rho\in M_D(\C)$ be positive definite.  The $\rho$-weighted inner
    product on $M_D(\C)$ is
    \begin{align}
        \langle X,Y\rangle_\rho
         & =\tr\!\left(\rho X^\dagger Y\right).
    \end{align}
    This is the auxiliary inner product in
    \cite[Equation~(5.6)]{Fannes1992Finitely}.
\end{theorem}

\begin{theorem}[Weighted norm formula]%
    \label{thm:weighted-virtual-norm-formula}
    \lean{Matrix.rhoWeighted_norm_sq}
    \leanok
    \uses{thm:weighted-virtual-inner-product}
    For every $X\in M_D(\C)$,
    \begin{align}
        \lVert X\rVert_\rho^2
         & =\operatorname{Re}\tr\!\left(\rho X^\dagger X\right).
    \end{align}
\end{theorem}
\begin{proof}
    \leanok
    \uses{thm:weighted-virtual-inner-product}
    By Theorem~\ref{thm:weighted-virtual-inner-product} with $Y=X$,
    \begin{align}
        \lVert X\rVert_\rho^2
         & =\operatorname{Re}\langle X,X\rangle_\rho             \\
         & =\operatorname{Re}\tr\!\left(\rho X^\dagger X\right).
    \end{align}
\end{proof}

\begin{definition}[Euclidean realization of the weighted matrix space]%
    \label{def:weighted-virtual-euclidean-realization}
    \lean{Matrix.rhoWeightedEquivEuclidean}
    \lean{Matrix.rhoWeightedEquivEuclidean_apply}
    \lean{Matrix.rhoWeightedEquivEuclidean_symm_apply}
    \leanok
    \uses{thm:weighted-virtual-inner-product}
    The map
    \begin{align}
        \Phi_\rho\colon M_D(\C) & \longrightarrow\C^{D^2},
                                &
        \Phi_\rho(X)            & =\operatorname{vec}\!\left(X\sqrt\rho\right)
    \end{align}
    is a complex linear isometric isomorphism.  Its inverse is
    \begin{align}
        \Phi_\rho^{-1}(z)
         & =\operatorname{unvec}(z)\,(\sqrt\rho)^{-1}.
    \end{align}
    Thus the weighted matrix space has complex Hilbert dimension $D^2$.
    This Euclidean realization is not asserted as a result of Fannes,
    Nachtergaele, and Werner.
\end{definition}

\subsection{Nested ground projections and martingale differences}%
\label{subsec:nested-ground-projection-differences}

% Source: References/cond-mat_9410110/main.tex:1060--1073, equation En.
% Source: References/cond-mat_9410110/main.tex:1195--1205,
% Theorem 2.1(i), equations resolutionpsi and resolutionnormpsi.

\begin{definition}[Nested ground projections and their martingale differences]%
    \label{def:nested-ground-projection-differences}
    \lean{FrustrationFree.NestedGroundProjections}
    \lean{FrustrationFree.NestedGroundProjections.martingaleDifference}
    \leanok
    Let $\mathcal H$ be a complex Hilbert space.  A \emph{nested ground-projection
        family} is a sequence $(G_n)_{n\geq0}$ of orthogonal projections such that
    \begin{align}
        m\leq n
        \quad & \Longrightarrow\quad
        \operatorname{ran}(G_n)\subseteq\operatorname{ran}(G_m).
    \end{align}
    Its $n$th martingale difference is
    \begin{align}
        E_n & =G_n-G_{n+1}.
    \end{align}
\end{definition}

\begin{theorem}[Orthogonal martingale-difference resolution]%
    \label{thm:orthogonal-martingale-difference-resolution}
    \lean{FrustrationFree.NestedGroundProjections.martingaleDifference_isSymmetricProjection}
    \lean{FrustrationFree.NestedGroundProjections.martingaleDifference_comp}
    \lean{FrustrationFree.NestedGroundProjections.sum_martingaleDifference_eq_id}
    \leanok
    \uses{def:nested-ground-projection-differences}
    Each $E_n$ is an orthogonal projection, and
    \begin{align}
        E_mE_n & =\delta_{m,n}E_m.
    \end{align}
    With Nachtergaele's endpoint conventions
    $G_0=\mathbf 1$ and $G_{N+1}=0$, the differences form a complete family:
    \begin{align}
        \sum_{n=0}^{N}E_n & =\mathbf 1.
    \end{align}
    These are the assertions following the display labelled $E_n$ in
    \cite{Nachtergaele1996Spectral}.
\end{theorem}

\begin{proof}
    \leanok
    \uses{def:nested-ground-projection-differences}
    If $m\leq n$, nestedness and orthogonality give
    $G_mG_n=G_nG_m=G_n$.  Consequently,
    \begin{align}
        E_n^2
               & =(G_n-G_{n+1})^2
        =G_n-G_{n+1}=E_n,          \\
        E_mE_n & =0\qquad(m\ne n).
    \end{align}
    Self-adjointness follows from
    $E_n^\dagger=G_n^\dagger-G_{n+1}^\dagger=E_n$.
    The complete sum telescopes with both endpoints visible:
    \begin{align}
        \sum_{n=0}^{N}E_n
         & =\sum_{n=0}^{N}(G_n-G_{n+1})
        =G_0-G_{N+1}=\mathbf 1.
    \end{align}
\end{proof}

\begin{theorem}[Truncated martingale-difference reconstruction]%
    \label{thm:truncated-martingale-difference-reconstruction}
    \lean{FrustrationFree.NestedGroundProjections.sum_martingaleDifference_apply_of_mem_orthogonal}
    \leanok
    \uses{thm:orthogonal-martingale-difference-resolution}
    Suppose $G_0=\mathbf 1$.  For every $v\perp\operatorname{ran}(G_N)$,
    \begin{align}
        v & =\sum_{n=0}^{N-1}E_nv.
    \end{align}
    This is the identity \textup{resolutionpsi} in the proof of
    \cite[Theorem~2.1(i)]{Nachtergaele1996Spectral}.
\end{theorem}

\begin{proof}
    \leanok
    \uses{thm:orthogonal-martingale-difference-resolution}
    Since $G_N$ is the orthogonal projection onto its range,
    $v\perp\operatorname{ran}(G_N)$ implies $G_Nv=0$.  Hence
    \begin{align}
        \sum_{n=0}^{N-1}E_nv
         & =(G_0-G_N)v=v.
    \end{align}
\end{proof}

\begin{theorem}[Pythagorean identity for martingale differences]%
    \label{thm:martingale-difference-pythagorean-identity}
    \lean{FrustrationFree.NestedGroundProjections.norm_sq_sum_martingaleDifference_finset}
    \lean{FrustrationFree.NestedGroundProjections.norm_sq_sum_martingaleDifference}
    \leanok
    \uses{thm:orthogonal-martingale-difference-resolution}
    For every $v\in\mathcal H$ and finite $S\subset\N$,
    \begin{align}
        \left\lVert\sum_{n\in S}E_nv\right\rVert^2
         & =\sum_{n\in S}\lVert E_nv\rVert^2.
    \end{align}
    In particular, for $S=\{0,\ldots,N-1\}$,
    \begin{align}
        \left\lVert\sum_{n=0}^{N-1}E_nv\right\rVert^2
         & =\sum_{n=0}^{N-1}\lVert E_nv\rVert^2.
    \end{align}
\end{theorem}

\begin{proof}
    \leanok
    \uses{thm:orthogonal-martingale-difference-resolution}
    The summands are pairwise orthogonal because
    $\langle E_mv,E_nv\rangle=\langle v,E_mE_nv\rangle=0$ for $m\ne n$.
    The identity follows by the finite Pythagorean theorem.
\end{proof}

\begin{theorem}[Norm-square decomposition on the final orthogonal complement]%
    \label{thm:martingale-difference-final-complement-norm-decomposition}
    \lean{FrustrationFree.NestedGroundProjections.norm_sq_eq_sum_martingaleDifference_of_mem_orthogonal}
    \leanok
    \uses{thm:truncated-martingale-difference-reconstruction, thm:martingale-difference-pythagorean-identity}
    Suppose $G_0=\mathbf 1$.  For every $v\perp\operatorname{ran}(G_N)$,
    \begin{align}
        \lVert v\rVert^2
         & =\sum_{n=0}^{N-1}\lVert E_nv\rVert^2.
    \end{align}
    This is the identity \textup{resolutionnormpsi} in the proof of
    \cite[Theorem~2.1(i)]{Nachtergaele1996Spectral}.
\end{theorem}

\begin{proof}
    \leanok
    \uses{thm:truncated-martingale-difference-reconstruction, thm:martingale-difference-pythagorean-identity}
    Substitute the truncated reconstruction
    $v=\sum_{n=0}^{N-1}E_nv$ into the Pythagorean identity.
\end{proof}

% Source: References/cond-mat_9410110/main.tex:1206--1220,
% Theorem 2.1(i), equations Enpsi and Enpsi2.

\begin{theorem}[Outside-window martingale cancellation]%
    \label{thm:outside-window-martingale-cancellation}
    \lean{FrustrationFree.NestedGroundProjections.martingaleDifference_comp_localProjection_comp_eq_zero}
    \lean{FrustrationFree.NestedGroundProjections.inner_sum_martingaleDifference_localProjection_eq_inner_sum_Icc}
    \leanok
    \uses{thm:orthogonal-martingale-difference-resolution}
    Let $Q_n$ be a local ground projection and suppose that
    \begin{align}
        E_mQ_n & =Q_nE_m
               &         & \text{whenever }m<n-l\text{ or }n<m.
    \end{align}
    Then every index outside $[n-l,n]$ cancels:
    \begin{align}
        E_mQ_nE_n & =0
                  &    & \text{whenever }m<n-l\text{ or }n<m.
    \end{align}
    If $n<N$, the complete resolution term therefore reduces to
    \begin{align}
        \left\langle v,
        \sum_{m=0}^{N-1}E_mQ_nE_nv\right\rangle
         & =\left\langle
        \sum_{m=n-l}^{n}E_mv,Q_nE_nv\right\rangle.
    \end{align}
    This is the reduction from \textup{Enpsi} to the first line of
    \textup{Enpsi2} in \cite[Theorem~2.1(i)]{Nachtergaele1996Spectral}.
\end{theorem}

\begin{proof}
    \leanok
    \uses{thm:orthogonal-martingale-difference-resolution}
    For an outside index, commutation and orthogonality give
    \begin{align}
        E_mQ_nE_n & =Q_nE_mE_n=0,
    \end{align}
    since either inequality implies $m\ne n$.  Hence only
    $n-l\leq m\leq n$ remains in the finite sum.  Finally, symmetry of each
    $E_m$ gives
    \begin{align}
        \langle v,E_mQ_nE_nv\rangle
         & =\langle E_mv,Q_nE_nv\rangle,
    \end{align}
    and summing over the active interval proves the displayed identity.
\end{proof}

% Source: References/cond-mat_9410110/main.tex:1060--1073, equation En.
% Source: References/cond-mat_9410110/main.tex:1206--1220,
% Theorem 2.1(i), equations Enpsi and Enpsi2.

\begin{definition}[Fixed-volume prefix and interval ground projections]%
    \label{def:fixed-volume-prefix-ground-projections}
    \lean{MPSTensor.openPrefixParentHamiltonianES,
        MPSTensor.openPrefixGroundProjectionES,
        MPSTensor.openIntervalGroundProjectionES}
    \leanok
    \uses{rem:euclidean_local_projector_ingredients}
    Fix an $N$-site Hilbert space and let $h_i$ denote the nonwrapping
    interaction of length $L$ starting at $i$.  For every $m\geq0$, set
    \begin{align}
        H_m^{(N)} & =\sum_{i+L\leq m}h_i,
                  & G_m^{(N)}             & =P_{\ker H_m^{(N)}}.
    \end{align}
    All operators act on the same $N$-site space.  The local interval ground
    projection in the proof of Theorem~2.1(i) is
    \begin{align}
        Q_n & =P_{\ker H_{[n-l,n+1)}},
            & H_{[n-l,n+1)}            & =\sum_{n-l\leq i,\ i+L\leq n+1}h_i.
    \end{align}
\end{definition}

\begin{lemma}[Positivity of the fixed-volume Hamiltonians]%
    \label{lem:fixed-volume-hamiltonians-positive}
    \lean{MPSTensor.openPrefixParentHamiltonianES_isPositive,
        MPSTensor.openSuffixParentHamiltonianES_isPositive}
    \leanok
    \uses{def:fixed-volume-prefix-ground-projections,
        lem:transported_es_positive}
    The prefix and interval Hamiltonians are positive and symmetric:
    \begin{align}
        H_m^{(N)}
         & =\sum_{i+L\leq m}h_i\geq0,
         & \bigl(H_m^{(N)}\bigr)^\dagger
         & =H_m^{(N)},                              \\
        H_{[n-l,n+1)}
         & =\sum_{n-l\leq i,\ i+L\leq n+1}h_i\geq0,
         & H_{[n-l,n+1)}^\dagger
         & =H_{[n-l,n+1)}.
    \end{align}
\end{lemma}

\begin{proof}
    \leanok
    \uses{def:fixed-volume-prefix-ground-projections,
        lem:transported_es_positive}
    Each $h_i$ is an orthogonal projection and is therefore positive.  Finite
    sums of positive operators are positive, and every positive operator is
    symmetric.
\end{proof}

\begin{theorem}[Nested fixed-volume prefix ground projections]%
    \label{thm:fixed-volume-prefix-ground-projections-nested}
    \lean{MPSTensor.localTermES_eq_zero_of_openPrefixParentHamiltonianES_eq_zero,
        MPSTensor.ker_openPrefixParentHamiltonianES_antitone,
        MPSTensor.fixedAmbientNestedGroundProjectionsES}
    \leanok
    \uses{def:fixed-volume-prefix-ground-projections,
        def:nested-ground-projection-differences}
    If $m\leq n$, then
    \begin{align}
        \ker H_n^{(N)} & \subseteq\ker H_m^{(N)},
                       & \operatorname{ran}G_n^{(N)}
                       & \subseteq\operatorname{ran}G_m^{(N)}.
    \end{align}
    Hence $(G_m^{(N)})_{m\geq0}$ is a nested ground-projection family.
\end{theorem}

\begin{proof}
    \leanok
    \uses{def:fixed-volume-prefix-ground-projections,
        lem:fixed-volume-hamiltonians-positive}
    If $H_n^{(N)}v=0$, then orthogonality of the local projections gives
    \begin{align}
        0=\langle v,H_n^{(N)}v\rangle
         & =\sum_{i+L\leq n}\lVert h_iv\rVert^2.
    \end{align}
    Every summand is nonnegative, so $h_iv=0$ whenever $i+L\leq n$.
    When $m\leq n$, every summand of $H_m^{(N)}$ occurs in $H_n^{(N)}$,
    and hence $H_m^{(N)}v=0$.  This proves the kernel inclusion, and the
    range inclusion follows from
    $\operatorname{ran}G_m^{(N)}=\ker H_m^{(N)}$.
\end{proof}

\begin{theorem}[Physical outside-window commutation and cancellation]%
    \label{thm:physical-outside-window-martingale-cancellation}
    \lean{LinearMap.IsSymmetric.kernelProjection_commute_of_commute,
        Submodule.starProjection_commute_of_le,
        MPSTensor.cyclicWindowsDisjoint_of_nonwrapping_ordered,
        MPSTensor.openPrefixParentHamiltonianES_commute_openSuffixParentHamiltonianES,
        MPSTensor.openPrefixGroundProjectionES_commute_openIntervalGroundProjectionES_of_le,
        MPSTensor.ker_openPrefixParentHamiltonianES_le_ker_openSuffixParentHamiltonianES,
        MPSTensor.openPrefixGroundProjectionES_commute_openIntervalGroundProjectionES_of_lt,
        MPSTensor.fixedAmbient_martingaleDifference_commute_openIntervalGroundProjectionES,
        MPSTensor.inner_sum_fixedAmbient_martingaleDifference_openIntervalGroundProjectionES_eq_Icc}
    \leanok
    \uses{def:fixed-volume-prefix-ground-projections,
        lem:fixed-volume-hamiltonians-positive,
        thm:fixed-volume-prefix-ground-projections-nested,
        thm:outside-window-martingale-cancellation}
    Assume $L\leq N$ and let $E_m=G_m^{(N)}-G_{m+1}^{(N)}$.  Then
    \begin{align}
        E_mQ_n & =Q_nE_m
               &         & \text{if }m<n-l\text{ or }n<m.
    \end{align}
    Consequently, for $n<N$ and every vector $v$,
    \begin{align}
        \left\langle v,
        \sum_{m=0}^{N-1}E_mQ_nE_nv\right\rangle
         & =\left\langle\sum_{m=n-l}^{n}E_mv,Q_nE_nv\right\rangle.
    \end{align}
\end{theorem}

\begin{proof}
    \leanok
    \uses{def:fixed-volume-prefix-ground-projections,
        lem:fixed-volume-hamiltonians-positive,
        thm:fixed-volume-prefix-ground-projections-nested,
        thm:outside-window-martingale-cancellation}
    If $m<n-l$, the supports of every summand of $H_m^{(N)}$ and every
    summand of $H_{[n-l,n+1)}$ are disjoint.  Their local interactions commute,
    hence the two symmetric Hamiltonians commute.  More generally, suppose
    that $T$ is symmetric and that an arbitrary linear operator $S$ commutes
    with $T$.  Then $\ker S$ is invariant under $T$, and symmetry of $T$ makes
    $(\ker S)^\perp$ invariant as well.  Thus $P_{\ker S}$ commutes with $T$.
    Since $P_{\ker S}$ is symmetric, applying the same argument to $\ker T$
    shows
    \begin{align}
        P_{\ker S}P_{\ker T} & =P_{\ker T}P_{\ker S}.
    \end{align}
    This proves the first range of indices for both $G_m^{(N)}$ and
    $G_{m+1}^{(N)}$.

    If $n<m$, every local interaction in $H_{[n-l,n+1)}$ also occurs in
    $H_m^{(N)}$.  Therefore
    \begin{align}
        \ker H_m^{(N)} & \subseteq\ker H_{[n-l,n+1)},
    \end{align}
    and the corresponding nested orthogonal projections commute.  Subtracting
    the two prefix identities gives $E_mQ_n=Q_nE_m$ in both index ranges.  The
    abstract outside-window cancellation theorem now restricts the full sum to
    $m\in[n-l,n]$.
\end{proof}

% Source: References/cond-mat_9410110/main.tex, condition C3, lines 1083--1094;
% proof of Theorem 2.1(i), lines 1178--1259.
\begin{lemma}[Right-spectator extension]
    \label{lem:right-spectator-extension}
    \lean{ContinuousLinearMap.rightFiber,
        ContinuousLinearMap.rightFiberwiseMap,
        ContinuousLinearMap.rightFiberwiseMap_add,
        ContinuousLinearMap.rightFiberwiseMap_sub,
        ContinuousLinearMap.rightFiberwiseMap_sum,
        ContinuousLinearMap.rightFiberwiseMap_comp,
        ContinuousLinearMap.norm_conj_linearIsometryEquiv,
        ContinuousLinearMap.norm_rightFiberwiseMap_le,
        ContinuousLinearMap.norm_rightFiberwiseMap}
    \leanok
    Let $T$ and $U$ be operators on a finite-dimensional active Hilbert space,
    and let $S$ be a finite set of right-spectator configurations.  The operator
    $T^{\oplus S}$, defined by applying $T$ independently at every $s\in S$,
    satisfies
    \begin{align}
        (T^{\oplus S})(U^{\oplus S}) & =(TU)^{\oplus S}.
    \end{align}
    It also satisfies
    \begin{align}
        \lVert T^{\oplus S}\rVert & \leq\lVert T\rVert.
    \end{align}
    If $S$ is nonempty, then the second relation is an equality.
\end{lemma}

\begin{proof}
    \leanok
    The squared norm splits as a sum over spectator fibers:
    \begin{align}
        \lVert T^{\oplus S}v\rVert^2
         & =\sum_{s\in S}\lVert Tv_s\rVert^2
        \leq\lVert T\rVert^2\sum_{s\in S}\lVert v_s\rVert^2.
    \end{align}
    For nonempty $S$, a vector supported on one spectator fiber gives the
    reverse inequality between the two operator norms.  Addition, subtraction,
    and composition are checked independently on each fiber.
\end{proof}

\begin{lemma}[Right-spectator coordinates and kernel projections]
    \label{lem:right-spectator-kernel-projection}
    \lean{MPSTensor.rightSpectatorConfigEquiv,
        MPSTensor.rightSpectatorConfigLinearIsometryEquiv,
        ContinuousLinearMap.mem_ker_rightFiberwiseMap_iff,
        ContinuousLinearMap.isSymmetricProjection_rightFiberwiseMap,
        ContinuousLinearMap.ker_starProjection_rightFiberwiseMap,
        MPSTensor.ker_starProjection_conj_linearIsometryEquiv}
    \leanok
    For $n,r\in\N$, restriction to the first $n$ sites and the final $r$ sites
    gives a canonical equivalence
    \begin{align}
        [d]^{n+r} & \simeq [d]^n\times[d]^r
    \end{align}
    and the corresponding Euclidean linear isometry.  A vector belongs to
    $\ker(T^{\oplus S})$ exactly when each spectator fiber belongs to $\ker T$.
    If $P_{\ker T}$ is the orthogonal projection onto $\ker T$, then
    \begin{align}
        P_{\ker(T^{\oplus S})} & =(P_{\ker T})^{\oplus S}.
    \end{align}
\end{lemma}

\begin{proof}
    \leanok
    The configuration equivalence is the inverse of concatenation.  The kernel
    characterization follows by evaluating each coordinate.  Fiberwise
    extension preserves symmetry and idempotence, and its range consists
    exactly of the vectors whose fibers lie in $\ker T$.  Equality of the two
    symmetric projections follows from equality of their ranges.
\end{proof}

\begin{theorem}[Nonwrapping interactions with right spectators]
    \label{thm:local-term-right-spectator-conjugacy}
    \lean{MPSTensor.localTermES_conj_rightSpectatorConfigLinearIsometryEquiv}
    \leanok
    \uses{lem:right-spectator-extension}
    Let $i+L\leq n$.  Under the canonical identification
    \begin{align}
        \ell^2([d]^{n+r}) & \simeq
        \ell^2([d]^n\times[d]^r),
    \end{align}
    the nonwrapping length-$L$ interaction beginning at $i$ is
    \begin{align}
        h^{(n+r)}_{i,L} & =(h^{(n)}_{i,L})^{\oplus[d]^r}.
    \end{align}
\end{theorem}

\begin{proof}
    \leanok
    The nonwrapping hypothesis identifies cyclic restriction with contiguous
    restriction on both chain lengths.  Concatenation leaves the first $n$
    coordinates unchanged and fixes the final $r$ spectator coordinates.
    Hence window extraction and restriction agree on every spectator fiber,
    and applying the same local parent interaction gives the identity.
\end{proof}

\begin{theorem}[Open Hamiltonian sums with right spectators]
    \label{thm:open-hamiltonian-right-spectator-conjugacy}
    \lean{MPSTensor.openPrefixParentHamiltonianES_conj_rightSpectatorConfigLinearIsometryEquiv,
        MPSTensor.openSuffixParentHamiltonianES_conj_rightSpectatorConfigLinearIsometryEquiv}
    \leanok
    \uses{thm:local-term-right-spectator-conjugacy}
    Let $L>0$ and $p\leq n$.  Under the canonical identification
    \begin{align}
        \ell^2([d]^{n+r}) & \simeq
        \ell^2([d]^n\times[d]^r),
    \end{align}
    both the prefix Hamiltonian on $[0,p)$ and every suffix-interval
    Hamiltonian ending at $p$ are direct sums over the right-spectator
    configurations of their operators in the $n$-site volume.
\end{theorem}

\begin{proof}
    \leanok
    The condition $L>0$ makes the active starts in the two finite sums exactly
    the starts before $p$.  Concatenation gives a bijection between the
    corresponding start sets in the two volumes.  Apply
    Theorem~\ref{thm:local-term-right-spectator-conjugacy} to each summand and
    then sum over this bijection.
\end{proof}

\begin{theorem}[Fixed-volume C3 ground projections with right spectators]
    \label{thm:fixed-volume-c3-projector-conjugacy}
    \lean{MPSTensor.openPrefixParentHamiltonianES_self_eq_openParentHamiltonianES,
        MPSTensor.ker_openPrefixParentHamiltonianES_eq_openChainLeftGroundSpaceES,
        MPSTensor.openSuffixParentHamiltonianES_full_window_eq_localTermES,
        MPSTensor.ker_openSuffixParentHamiltonianES_eq_openChainTailGroundSpaceES,
        MPSTensor.ker_openPrefixParentHamiltonianES_self_eq_groundSpaceES,
        MPSTensor.openPrefixGroundProjectionES_conj_rightSpectatorConfigLinearIsometryEquiv,
        MPSTensor.openPrefixWholeGroundProjectionES_conj_rightSpectatorConfigLinearIsometryEquiv,
        MPSTensor.openIntervalGroundProjectionES_conj_rightSpectatorConfigLinearIsometryEquiv}
    \leanok
    \uses{lem:right-spectator-kernel-projection,
        thm:open-hamiltonian-right-spectator-conjugacy}
    Let $K,l,n,r\in\mathbb{N}$.  Suppose that $A$ has nonzero bond dimension,
    is block injective at length $L_0>0$, and satisfies $L_0+1\leq n$.  The
    kernel of the range-$(L_0+1)$ prefix Hamiltonian on the first $n$ sites of
    the $(n+1)$-site volume is the left open-chain ground space.  When that
    prefix instead fills its own $n$-site volume, its kernel is the whole
    $n$-site ground space.  Independently,
    the kernel of the full final length-$(l+1)$ suffix window beginning at $K$
    is the tail open-chain ground space.  Let $U_{\mathrm{pre}}$ and
    $U_{\mathrm{tail}}$ be the canonical right-spectator isometries for the
    active-volume decompositions $(n+1+r)=(n+1)+r$ and
    $(K+l+1+r)=(K+l+1)+r$, respectively.  Their orthogonal projections satisfy
    \begin{align}
        U_{\mathrm{pre}} G_{\lbrack 0,n)}^{(n+1+r)}U_{\mathrm{pre}}^{-1}
         & =(G_{\lbrack 0,n)}^{(n+1)})^{\oplus{[d]}^r},       \\
        U_{\mathrm{tail}} G_{\lbrack K,K+l+1)}^{(K+l+1+r)}U_{\mathrm{tail}}^{-1}
         & =(G_{\lbrack K,K+l+1)}^{(K+l+1)})^{\oplus{[d]}^r}.
    \end{align}
\end{theorem}

\begin{proof}
    \leanok
    The prefix kernel is read fiberwise.  Block injectivity identifies every
    fiber kernel with the $n$-site MPS ground space, which is precisely the
    left-window membership condition after fixing the last physical index.
    When the active prefix fills its smaller volume, the same argument
    identifies its kernel projection with the whole open-chain ground
    projection.  The full suffix interval contains one interaction; its cyclic
    restriction
    is the ordinary restriction obtained by fixing the $K$-site prefix, giving
    the tail-window membership condition.  Orthogonal kernel projections are
    preserved by the canonical isometry, and the fiberwise kernel-projection
    identity gives the two displayed formulas.
\end{proof}

\begin{theorem}[C3 boundary indices in the fixed final volume]
    \label{thm:fixed-volume-c3-boundary-indices}
    \lean{MPSTensor.openPrefixParentHamiltonianES_eq_zero_of_lt,
        MPSTensor.openPrefixGroundProjectionES_eq_one_of_lt,
        MPSTensor.fixedAmbient_martingaleDifference_eq_zero_of_lt,
        MPSTensor.openPrefixParentHamiltonianES_eq_openSuffixParentHamiltonianES_at_endpoint,
        MPSTensor.openIntervalGroundProjectionES_at_endpoint,
        MPSTensor.openIntervalGroundProjectionES_comp_martingaleDifference_at_endpoint}
    \leanok
    \uses{def:fixed-volume-prefix-ground-projections}
    Suppose that $0<l$ and $l+1\leq N$.  Take interaction length $L=l+1$ and
    define $E_n=G_n^{(N)}-G_{n+1}^{(N)}$.  Then the two boundary ranges in the
    fixed $N$-site space are explicit:
    \begin{align}
        n<l & \Longrightarrow E_n=0.
    \end{align}
    At the remaining boundary index,
    \begin{align}
        Q_lE_l & =0.
    \end{align}
    At $n=l$, both $Q_l$ and $G_{l+1}^{(N)}$ are the kernel projection of the
    unique interaction on $[0,l+1)$.  Thus
    \begin{align}
        E_l & =\Id-G_{l+1}^{(N)}.
    \end{align}
    Hence
    \begin{align}
        Q_lE_l & =G_{l+1}^{(N)}(\Id-G_{l+1}^{(N)})=0.
    \end{align}
\end{theorem}

\begin{proof}
    \leanok
    \uses{def:fixed-volume-prefix-ground-projections}
    If $n<l$, neither the $n$-site prefix nor the $(n+1)$-site prefix contains a full
    length-$(l+1)$ interaction, so both ground projections are the identity.
    At $n=l$, the next prefix and the interval $[0,l+1)$ contain exactly the
    same single interaction.  Idempotence of its kernel projection proves the
    second identity.
\end{proof}

\subsection{Spectator-indexed boundary maps}\label{subsec:spectator_boundary_maps}

The tail and left boundary maps embed virtual-space data into a common
$(K+L)$-site Hilbert space, expressing overlapping ground spaces as ranges of
linear maps on the same ambient space to compare their orthogonal projectors.

\begin{definition}[Prefix restriction]\label{def:prefix_restrict}
    \lean{MPSTensor.prefixRestrictₗ, MPSTensor.prefixRestrictₗ_apply}
    \leanok
    \uses{def:tail_restrict_parent}
    For $K,L\in\N$ and a fixed suffix configuration $\tau\colon\Fin L\to\Fin d$, the
    \emph{prefix restriction} is the $\C$-linear map
    $P_\tau\colon\C^{\Fin d^{K+L}}\to\C^{\Fin d^K}$ defined by
    $(P_\tau\psi)(\sigma)=\psi(\sigma,\tau)$ for $\sigma\colon\Fin K\to\Fin d$.
    This is the left-sided analogue of the suffix restriction of
    Definition~\ref{def:tail_restrict_parent}.
\end{definition}

\begin{lemma}[Prefix restriction of a ground-space vector]\label{lem:prefix_restrict_ground_space_map}
    \lean{MPSTensor.prefixRestrictₗ_groundSpaceMap}
    \leanok
    \uses{def:prefix_restrict, def:ground_space_map, lem:eval_word_append}
    For every boundary matrix $X\in\MN{D}$ and every suffix $\tau$,
    \begin{align}
        P_\tau\bigl(\Gamma_{K+L}(X)\bigr)
         & =\Gamma_K\bigl(A^\tau X\bigr).
    \end{align}
\end{lemma}

\begin{proof}
    \leanok
    \uses{lem:eval_word_append}
    Expand $\Gamma_{K+L}(X)(\sigma,\tau)=\tr(A^{\sigma,\tau}X)$,
    factor $A^{\sigma,\tau}=A^\sigma A^\tau$ by
    Lemma~\ref{lem:eval_word_append}, and reassociate:
    $\tr((A^\sigma A^\tau)X)=\tr(A^\sigma(A^\tau X))
        =\Gamma_K(A^\tau X)(\sigma)$.
\end{proof}

\begin{definition}[Tail boundary map]\label{def:tail_boundary_map}
    \lean{MPSTensor.tailBoundaryMap, MPSTensor.tailBoundaryMap_apply}
    \leanok
    \uses{def:eval_word}
    For $K,L\in\N$, the \emph{tail boundary map} is the $\C$-linear map
    \begin{align}
        \Gamma^{\mathrm{tail}}_{K,L}\colon
        \bigl((\Fin K\to\Fin d)\to\MN{D}\bigr)
        \to\C^{\Fin d^{K+L}}
    \end{align}
    defined for every family $Y\colon(\Fin K\to\Fin d)\to\MN{D}$ by
    \begin{align}
        \Gamma^{\mathrm{tail}}_{K,L}(Y)(u,\tau)
         & =\tr\bigl(A^\tau\,Y_u\bigr),
    \end{align}
    where $u\colon\Fin K\to\Fin d$ is the prefix and
    $\tau\colon\Fin L\to\Fin d$ is the suffix.
    When $Y_u=X A^u$ the output equals the standard $\Gamma_{K+L}(X)$.
\end{definition}

\begin{definition}[Left boundary map]\label{def:left_boundary_map}
    \lean{MPSTensor.leftBoundaryMap, MPSTensor.leftBoundaryMap_apply}
    \leanok
    \uses{def:eval_word}
    For $K,L\in\N$, the \emph{left boundary map} is the $\C$-linear map
    \begin{align}
        \Gamma^{\mathrm{left}}_{K,L}\colon
        \bigl((\Fin L\to\Fin d)\to\MN{D}\bigr)
        \to\C^{\Fin d^{K+L}}
    \end{align}
    defined for every family $Z\colon(\Fin L\to\Fin d)\to\MN{D}$ by
    \begin{align}
        \Gamma^{\mathrm{left}}_{K,L}(Z)(u,\tau)
         & =\tr\bigl(A^{u}\,Z_{\tau}\bigr),
    \end{align}
    where $u\colon\Fin K\to\Fin d$ is the prefix and
    $\tau\colon\Fin L\to\Fin d$ is the suffix.
    When $Z_\tau=A^\tau X$ the output equals the standard $\Gamma_{K+L}(X)$.
\end{definition}

\begin{theorem}[Tail boundary map on a prefix--suffix pair]\label{thm:tail_boundary_map_append}
    \lean{MPSTensor.tailBoundaryMap_append}
    \leanok
    \uses{def:tail_boundary_map, def:ground_space_map}
    For every prefix $u\colon\Fin K\to\Fin d$ and suffix $\tau\colon\Fin L\to\Fin d$,
    \begin{align}
        \Gamma^{\mathrm{tail}}_{K,L}(Y)(u,\tau)
         & =\Gamma_L(Y_u)(\tau).
    \end{align}
\end{theorem}

\begin{proof}
    \leanok
    Unfold the definition of $\Gamma^{\mathrm{tail}}_{K,L}$ on the pair
    $(u,\tau)=\Fin.\!append\;u\;\tau$, use the decomposition
    $(\Fin.\!append\;u\;\tau)\circ\Fin.\!natAdd\;K=\tau$ and
    $(\Fin.\!append\;u\;\tau)\circ\Fin.\!castAdd\;L=u$, and apply
    $\Gamma_L(Y_u)(\tau)=\tr(A^\tau Y_u)$.
\end{proof}

\begin{theorem}[Suffix restriction of the tail boundary map]\label{thm:tail_restrict_tail_boundary_map}
    \lean{MPSTensor.tailRestrictₗ_tailBoundaryMap}
    \leanok
    \uses{thm:tail_boundary_map_append, def:tail_restrict_parent}
    Restricting the tail boundary map to the suffix for a fixed prefix
    $u\in\Fin d^K$ recovers $\Gamma_L(Y_u)$:
    \begin{align}
        R^{\mathrm{tail}}_u\bigl(\Gamma^{\mathrm{tail}}_{K,L}(Y)\bigr)
         & =\Gamma_L(Y_u).
    \end{align}
\end{theorem}

\begin{proof}
    \leanok
    \uses{thm:tail_boundary_map_append}
    This is Theorem~\ref{thm:tail_boundary_map_append} applied pointwise:
    for each $\tau$, $(R^{\mathrm{tail}}_u\psi)(\tau)=\psi(u,\tau)$.
\end{proof}

\begin{theorem}[Tail boundary map factorization]\label{thm:tail_boundary_map_factorization}
    \lean{MPSTensor.tailBoundaryMap_factorization}
    \leanok
    \uses{def:tail_boundary_map, def:ground_space_map, lem:eval_word_append}
    The standard length-$(K+L)$ boundary map factors through the tail boundary map:
    \begin{align}
        \Gamma_{K+L}(X)
         & =\Gamma^{\mathrm{tail}}_{K,L}\bigl(u\mapsto X A^{u}\bigr)
    \end{align}
    for every $X\in\MN{D}$.
\end{theorem}

\begin{proof}
    \leanok
    \uses{lem:eval_word_append}
    Write a configuration $\sigma\in\Fin d^{K+L}$ as $\sigma=(u,\tau)$ with
    $u\in\Fin d^K$ and $\tau\in\Fin d^L$.  Substitute $Y_u=X A^u$ into the
    definition of $\Gamma^{\mathrm{tail}}_{K,L}$, apply the three-factor
    trace cycle $\tr(A^\tau(X A^u))=\tr(A^u(A^\tau X))$, reassociate
    $\tr(A^u(A^\tau X))=\tr((A^u A^\tau)X)$, factor $A^u A^\tau=A^{u,\tau}$
    by multiplicativity of word evaluation (Lemma~\ref{lem:eval_word_append}),
    and identify the result with $\Gamma_{K+L}(X)(u,\tau)=\tr(A^{u,\tau}X)$.
\end{proof}

\begin{theorem}[Tail boundary map range]\label{thm:tail_boundary_map_range_eq}
    \lean{MPSTensor.tailBoundaryMap_range_eq}
    \leanok
    \uses{def:tail_boundary_map, def:ground_space_parent, def:tail_restrict_parent}
    The range of the tail boundary map is exactly the tail ground space:
    \begin{align}
        \range\bigl(\Gamma^{\mathrm{tail}}_{K,L}\bigr)
         & =\bigl\{\psi\in\C^{\Fin d^{K+L}}
        \mid \forall u\in\Fin d^K,\;
        R^{\mathrm{tail}}_u\psi\in G_L(A)\bigr\}.
    \end{align}
\end{theorem}

\begin{proof}
    \leanok
    \uses{thm:tail_restrict_tail_boundary_map, def:tail_restrict_parent, def:ground_space_parent}
    For the forward inclusion, if $\psi=\Gamma^{\mathrm{tail}}_{K,L}(Y)$ then
    $R^{\mathrm{tail}}_u\psi=\Gamma_L(Y_u)$ for every prefix $u$, so the suffix
    restriction lies in $G_L(A)$.
    Conversely, given $\psi$ with $R^{\mathrm{tail}}_u\psi\in G_L(A)$ for every $u$,
    choose for each $u$ a matrix $Y_u\in\MN{D}$ with
    $\Gamma_L(Y_u)=R^{\mathrm{tail}}_u\psi$.  Then
    $\Gamma^{\mathrm{tail}}_{K,L}(Y)(u,\tau)=\tr(A^\tau Y_u)
        =\Gamma_L(Y_u)(\tau)=(R^{\mathrm{tail}}_u\psi)(\tau)=\psi(u,\tau)$,
    so $\psi$ lies in the range.
\end{proof}

\begin{theorem}[Left boundary map on a prefix--suffix pair]\label{thm:left_boundary_map_append}
    \lean{MPSTensor.leftBoundaryMap_append}
    \leanok
    \uses{def:left_boundary_map, def:ground_space_map}
    For every prefix $u\colon\Fin K\to\Fin d$ and suffix $\tau\colon\Fin L\to\Fin d$,
    \begin{align}
        \Gamma^{\mathrm{left}}_{K,L}(Z)(u,\tau)
         & =\Gamma_K(Z_\tau)(u).
    \end{align}
\end{theorem}

\begin{proof}
    \leanok
    Unfold the definition of $\Gamma^{\mathrm{left}}_{K,L}$ on the pair
    $(u,\tau)$, use $(\Fin.\!append\;u\;\tau)\circ\Fin.\!castAdd\;L=u$ and
    $(\Fin.\!append\;u\;\tau)\circ\Fin.\!natAdd\;K=\tau$, and apply
    $\Gamma_K(Z_\tau)(u)=\tr(A^u Z_\tau)$.
\end{proof}

\begin{theorem}[Prefix restriction of the left boundary map]\label{thm:prefix_restrict_left_boundary_map}
    \lean{MPSTensor.prefixRestrictₗ_leftBoundaryMap}
    \leanok
    \uses{thm:left_boundary_map_append, def:prefix_restrict}
    Restricting the left boundary map to the prefix for a fixed suffix
    $\tau\in\Fin d^L$ recovers $\Gamma_K(Z_\tau)$:
    \begin{align}
        P_\tau\bigl(\Gamma^{\mathrm{left}}_{K,L}(Z)\bigr)
         & =\Gamma_K(Z_\tau).
    \end{align}
\end{theorem}

\begin{proof}
    \leanok
    \uses{thm:left_boundary_map_append}
    Pointwise application of Theorem~\ref{thm:left_boundary_map_append}:
    for each $\sigma$, $(P_\tau\psi)(\sigma)=\psi(\sigma,\tau)$.
\end{proof}

\begin{theorem}[Left boundary map factorization]\label{thm:left_boundary_map_factorization}
    \lean{MPSTensor.leftBoundaryMap_factorization}
    \leanok
    \uses{def:left_boundary_map, def:ground_space_map, lem:eval_word_append}
    The standard length-$(K+L)$ boundary map factors through the left boundary map:
    \begin{align}
        \Gamma_{K+L}(X)
         & =\Gamma^{\mathrm{left}}_{K,L}\bigl(\tau\mapsto A^{\tau}X\bigr)
    \end{align}
    for every $X\in\MN{D}$.
\end{theorem}

\begin{proof}
    \leanok
    \uses{lem:eval_word_append}
    Substitute $Z_\tau=A^\tau X$ into the definition of
    $\Gamma^{\mathrm{left}}_{K,L}$, giving
    $\tr(A^u(A^\tau X))=\tr((A^u A^\tau)X)$.  Factor $A^u A^\tau=A^{u,\tau}$ by
    Lemma~\ref{lem:eval_word_append} and identify the result with
    $\Gamma_{K+L}(X)(u,\tau)$.
\end{proof}

\begin{theorem}[Left boundary map range]\label{thm:left_boundary_map_range_eq}
    \lean{MPSTensor.leftBoundaryMap_range_eq}
    \leanok
    \uses{def:left_boundary_map, def:ground_space_parent, def:prefix_restrict}
    The range of the left boundary map is exactly the left ground space:
    \begin{align}
        \range\bigl(\Gamma^{\mathrm{left}}_{K,L}\bigr)
         & =\bigl\{\psi\in\C^{\Fin d^{K+L}}
        \mid \forall \tau\in\Fin d^L,\;
        P_\tau\psi\in G_K(A)\bigr\}.
    \end{align}
\end{theorem}

\begin{proof}
    \leanok
    \uses{thm:prefix_restrict_left_boundary_map, def:prefix_restrict, def:ground_space_parent}
    The proof is symmetric to Theorem~\ref{thm:tail_boundary_map_range_eq}.
    If $\psi=\Gamma^{\mathrm{left}}_{K,L}(Z)$ then
    $P_\tau\psi=\Gamma_K(Z_\tau)$, so the prefix restriction lies in $G_K(A)$.
    Conversely, given $\psi$ with $P_\tau\psi\in G_K(A)$ for every $\tau$, choose
    $Z_\tau$ with $\Gamma_K(Z_\tau)=P_\tau\psi$; then
    $\Gamma^{\mathrm{left}}_{K,L}(Z)(u,\tau)=\tr(A^u Z_\tau)
        =\Gamma_K(Z_\tau)(u)=(P_\tau\psi)(u)=\psi(u,\tau)$.
\end{proof}

\begin{lemma}[Prefix restriction at $L=1$]\label{lem:prefix_restrict_one_eq_restrict_last}
    \lean{MPSTensor.prefixRestrictₗ_one_eq_restrictLast}
    \leanok
    \uses{def:prefix_restrict, def:restrict_last}
    For $L=1$, the prefix restriction $P_\tau$ coincides with the last-site
    restriction: $P_\tau\psi=R^{\mathrm{last}}(\psi,\tau(0))$ for every
    $\tau\colon\Fin 1\to\Fin d$.
\end{lemma}

\begin{proof}
    \leanok
    Unfolding both sides: for $\sigma\in\Fin d^K$,
    $(P_\tau\psi)(\sigma)=\psi(\sigma,\tau)$ while
    $R^{\mathrm{last}}(\psi,j)(\sigma)=\psi(\sigma_1,\dots,\sigma_K,j)$ with
    $j=\tau(0)$.  These are equal because $\Fin.\!append\;\sigma\;\tau$ puts
    $\tau$ after $\sigma$, matching the $\Fin.\!snoc$ pattern of $R^{\mathrm{last}}$.
\end{proof}

\begin{theorem}[Tail boundary map range: Euclidean-space transport]\label{thm:tail_boundary_map_range_es_transport}
    \lean{MPSTensor.tailBoundaryMap_range_map_symm_eq_openChainTailGroundSpaceES}
    \leanok
    \uses{thm:tail_boundary_map_range_eq, def:open_chain_ground_subspaces_es}
    For every $K,L\in\N$, transport of the tail boundary map range by the
    inverse canonical $\ell^2$ isomorphism
    identifies it with the open-chain tail ground submodule:
    \begin{align}
        \iota_{K+L}\bigl(\range(\Gamma^{\mathrm{tail}}_{K,L})\bigr)
         & =\mathcal V_{K,L}(A).
    \end{align}
    Concretely, a vector $v$ in the $(K+L)$-site Euclidean space belongs to
    $\mathcal V_{K,L}(A)$ exactly when the corresponding function satisfies the
    tail ground condition (Definition~\ref{def:in_tail_ground}).
\end{theorem}

\begin{proof}
    \leanok
    \uses{thm:tail_boundary_map_range_eq}
    By Theorem~\ref{thm:tail_boundary_map_range_eq}, the range equals the
    intersection over prefixes $u$ of the ground-space comaps
    $(R^{\mathrm{tail}}_u)^{-1}(G_L(A))$.  Transporting this intersection
    through the canonical $\ell^2$ identification and unfolding the definition
    of $\mathcal V_{K,L}(A)$ (Definition~\ref{def:open_chain_ground_subspaces_es})
    gives the stated equality.
\end{proof}

\begin{theorem}[Left boundary map range: Euclidean-space transport ($L=1$)]%
    \label{thm:left_boundary_map_range_es_transport}
    \lean{MPSTensor.leftBoundaryMap_range_map_symm_eq_openChainLeftGroundSpaceES}
    \leanok
    \uses{thm:left_boundary_map_range_eq, def:open_chain_ground_subspaces_es}
    At $L=1$, transport of the left boundary map range by the inverse canonical
    $\ell^2$ isomorphism identifies
    it with the open-chain left ground submodule:
    \begin{align}
        \iota_{K+1}\bigl(\range(\Gamma^{\mathrm{left}}_{K,1})\bigr)
         & =\mathcal U_K(A).
    \end{align}
    Concretely, a vector $v$ in the $(K+1)$-site Euclidean space belongs to
    $\mathcal U_K(A)$ exactly when the corresponding function satisfies the
    left ground condition (every fixed-final-site restriction lies in $G_K(A)$;
    Definition~\ref{def:in_left_ground}).
\end{theorem}

\begin{proof}
    \leanok
    \uses{thm:left_boundary_map_range_eq, lem:prefix_restrict_one_eq_restrict_last, def:restrict_last}
    By Theorem~\ref{thm:left_boundary_map_range_eq} with $L=1$, the range equals
    the intersection over suffixes $\tau\colon\Fin 1\to\Fin d$ of the comaps
    $P_\tau^{-1}(G_K(A))$.  For $L=1$, each suffix $\tau$ is determined by a
    single index $j\in\Fin d$, and Lemma~\ref{lem:prefix_restrict_one_eq_restrict_last}
    identifies $P_\tau$ with $R^{\mathrm{last}}(\cdot,j)$, the last-site restriction of
    Definition~\ref{def:restrict_last}.  Transporting the
    resulting intersection $\bigcap_{j\in\Fin d}(R^{\mathrm{last}}(\cdot,j))^{-1}(G_K(A))$
    through the canonical $\ell^2$ identification and
    unfolding Definition~\ref{def:open_chain_ground_subspaces_es} gives
    $\mathcal U_K(A)$.
\end{proof}

\begin{definition}[Euclidean spectator boundary maps]
    \label{def:spectator_boundary_maps_es}
    \lean{MPSTensor.tailBoundaryMapES, MPSTensor.leftBoundaryMapES}
    \leanok
    \uses{def:tail_boundary_map, def:left_boundary_map}
    Transport the tail and left boundary maps through the canonical Euclidean
    identifications on their virtual matrix families and physical coefficient
    spaces.  Denote the resulting continuous maps by
    $\Gamma^{\mathrm{tail}}_{K,L}$ and
    $\Gamma^{\mathrm{left}}_{K,L}$.
\end{definition}

\begin{theorem}[Euclidean spectator boundary ranges and factorizations]%
    \label{thm:spectator_boundary_hilbert_factorization}
    \lean{MPSTensor.range_tailBoundaryMapES}
    \lean{MPSTensor.range_leftBoundaryMapES_one}
    \lean{MPSTensor.tailBoundaryMapES_comp_tailVirtualMapES}
    \lean{MPSTensor.leftBoundaryMapES_comp_leftVirtualMapES}
    \leanok
    \uses{def:spectator_boundary_maps_es, thm:tail_boundary_map_range_es_transport, thm:left_boundary_map_range_es_transport, thm:tail_boundary_map_factorization, thm:left_boundary_map_factorization}
    Equip each virtual matrix fiber with the unweighted Hilbert--Schmidt inner
    product and flatten the finite family of fibers into a Euclidean space.
    The resulting continuous tail and left boundary maps have ranges
    $\mathcal V_{K,L}(A)$ and, at suffix length one, $\mathcal U_K(A)$.
    If $J^{\mathrm{tail}}_K(X)_u=XA^u$ and
    $J^{\mathrm{left}}_L(X)_\tau=A^\tau X$, then
    \begin{align}
        \Gamma^{\mathrm{tail}}_{K,L}J^{\mathrm{tail}}_K & =\Gamma_{K+L}, &
        \Gamma^{\mathrm{left}}_{K,L}J^{\mathrm{left}}_L & =\Gamma_{K+L}.
    \end{align}
\end{theorem}
\begin{proof}
    \leanok
    \uses{def:spectator_boundary_maps_es, thm:tail_boundary_map_range_es_transport, thm:left_boundary_map_range_es_transport, thm:tail_boundary_map_factorization, thm:left_boundary_map_factorization}
    Transport the algebraic range descriptions through the Euclidean
    identification.  The two factorization identities follow by evaluating a
    virtual matrix \(X\) fiberwise and applying the corresponding algebraic
    tail or left factorization.
\end{proof}

\begin{theorem}[Injectivity of the left spectator boundary map]
    \label{thm:left_boundary_map_es_injective}
    \lean{MPSTensor.leftBoundaryMapES_injective_of_groundSpaceMapES_injective}
    \leanok
    \uses{def:ground_space_map_es, def:spectator_boundary_maps_es}
    If the length-$K$ ground-space map is injective, then the left boundary map
    with any spectator length $L$ is injective.
\end{theorem}

\begin{proof}
    \leanok
    \uses{def:ground_space_map_es, def:spectator_boundary_maps_es, thm:ground_space_map_es_frobenius, def:prefix_restrict}
    Recover every spectator fiber by restricting to its fixed suffix and use
    injectivity of the length-$K$ ground-space map.
\end{proof}

\begin{theorem}[Direct-sum spectator Gram operators]%
    \label{thm:spectator_boundary_direct_sum_gram}
    \lean{MPSTensor.tailBoundaryMapES_adjoint_comp_self_eq_fiberwise_groundSpaceGram}
    \lean{MPSTensor.leftBoundaryMapES_adjoint_comp_self_eq_fiberwise_groundSpaceGram}
    \lean{MPSTensor.inverseGram_tailBoundaryMapES_eq_fiberwise_inverseGram}
    \lean{MPSTensor.inverseGram_leftBoundaryMapES_eq_fiberwise_inverseGram}
    \leanok
    The tail Gram is the direct sum over prefix configurations of copies of
    $\mathcal K_L$, and the left Gram is the direct sum over suffix
    configurations of copies of $\mathcal K_K$.  Whenever the
    corresponding local boundary map is injective, the inverse Gram is the
    direct sum of the local inverse Grams.  These identities include empty
    spectator index types and assert no norm estimate.
\end{theorem}
\begin{proof}
    \leanok
    Reindex the physical inner product by prefix--suffix configurations and
    evaluate each fiber.  This gives the direct-sum Gram identities.  Composing
    fiberwise with the local inverse Gram on either side gives the stated
    inverse formulas.
\end{proof}

\begin{theorem}[Norm of a fiberwise boundary operator]%
    \label{thm:boundary_fiberwise_map_norm}
    \lean{MPSTensor.norm_boundaryFiberwiseMap_le}
    \leanok
    Let $S$ be a finite spectator index set, let $V$ be the Euclidean virtual
    matrix space, and let $G:V\to V$ be linear.  If $G^{\oplus S}$ acts on a
    boundary family $x=(x_s)_{s\in S}$ by
    $(G^{\oplus S}x)_s=Gx_s$, then
    \begin{align}
        \lVert G^{\oplus S}\rVert\leq \lVert G\rVert .
    \end{align}
    This estimate also holds when $S$ is empty.
\end{theorem}
\begin{proof}
    \leanok
    Square both sides of the pointwise estimate
    $\lVert Gx_s\rVert\leq\lVert G\rVert\lVert x_s\rVert$, sum over $s\in S$,
    and use the Euclidean direct-sum norm identity.
\end{proof}

\begin{theorem}[Exact C3 projector residual identity]%
    \label{thm:c3_spectator_projector_residual}
    \lean{MPSTensor.c3_injectiveRangeProjector_residual_eq}
    \leanok
    \uses{thm:spectator_boundary_hilbert_factorization}
    Fix $K,L_0\in\mathbb N$.  Suppose that the tail boundary map for
    $(K,L_0+1)$, the left boundary map for $(K+L_0,1)$, and
    $\Gamma_{K+L_0+1}$ are injective.  The difference between the product of
    the first two inverse-Gram range projectors and the third inverse-Gram range
    projector is exactly the common-factorization residual obtained from the
    tail and left virtual maps.  Their ranges are, respectively,
    $\mathcal V_{K,L_0+1}(A)$, $\mathcal U_{K+L_0}(A)$, and
    $\mathcal G_{K+L_0+1}(A)$.  This identity asserts no numerical contraction
    estimate.
\end{theorem}
\begin{proof}
    \leanok
    \uses{thm:spectator_boundary_hilbert_factorization}
    Apply the common-factorization projector identity to the two spectator
    boundary maps, the full boundary map, and their tail and left virtual maps.
    These range equalities identify the three physical ranges.
\end{proof}

\begin{theorem}[Arbitrary-increment mixed Gram of the overlapping spectator boundary maps]%
    \label{thm:spectator_boundary_mixed_gram}
    \lean{MPSTensor.inner_tailBoundaryMapES_adjoint_leftBoundaryMapES_q}
    \lean{MPSTensor.inner_tailBoundaryMapES_adjoint_leftBoundaryMapES}
    \leanok
    \uses{def:tail_boundary_map, def:left_boundary_map, def:ground_space_gram}
    Let $q\in\mathbb N$, let $X=(X_u)_{u\in\Fin d^K}$, and let
    $Z=(Z_j)_{j\in\Fin d^q}$ be virtual boundary-matrix families.  After the
    canonical identification of $K+(l+q)$ with $(K+l)+q$, the mixed Gram for
    the tail window of length $l+q$ and the left window with a $q$-site
    increment is
    \begin{align}
        \left\langle X,
        (\Gamma^{\mathrm{tail}}_{K,l+q})^\dagger
        \Gamma^{\mathrm{left}}_{K+l,q}Z\right\rangle
         & =\sum_{u\in\Fin d^K}\sum_{j\in\Fin d^q}
        \left\langle
        U(A^jX_u),\,
        \mathcal K_l\,U(Z_jA^u)
        \right\rangle,
    \end{align}
    where $U$ is Frobenius vectorization and
    $\mathcal K_l=\Gamma_l^\dagger\Gamma_l$.  The multiplication order is
    fixed by
    $\tr(A^{\tau,j}X_u)=\tr(A^\tau(A^jX_u))$ and
    $\tr(A^{u,\tau}Z_j)=\tr(A^\tau(Z_jA^u))$.
    The one-site identity is the specialization $q=1$.  This identity expresses
    the algebraic form of condition (C3$'$) in
    \cite[Equation~(2.5)]{Nachtergaele1996Spectral}.
\end{theorem}
\begin{proof}
    \leanok
    \uses{def:tail_boundary_map, def:left_boundary_map, def:ground_space_gram, lem:eval_word_append, thm:ground_space_map_es_frobenius}
    Split a physical configuration into a prefix $u$, an overlap word $\tau$
    of length $l$, and an increment configuration $j$ of length $q$.  The
    adjoint relation turns the mixed Gram into the physical inner product of
    the two boundary vectors.  For each $(u,\tau,j)$, word multiplicativity
    and trace cyclicity give
    \begin{align}
        \tr(A^{\tau,j}X_u) & =\tr(A^\tau(A^jX_u)), &
        \tr(A^{u,\tau}Z_j) & =\tr(A^\tau(Z_jA^u)).
    \end{align}
    For fixed $(u,j)$, summing the resulting inner products over $\tau$ is
    exactly the boundary-map Gram pairing
    $\langle U(A^jX_u),\mathcal K_l U(Z_jA^u)\rangle$.  Summing over $u$ and
    $j$ proves the arbitrary-increment identity; setting $q=1$ gives the
    original formula.
\end{proof}

\begin{theorem}[Centered mixed Gram at the limiting metric]%
    \label{thm:spectator_boundary_mixed_gram_centered}
    \lean{MPSTensor.inner_tailBoundaryMapES_adjoint_leftBoundaryMapES_centered}
    \leanok
    \uses{thm:spectator_boundary_mixed_gram, thm:fixed_point_gram_metric_exact}
    Specialize Theorem~\ref{thm:spectator_boundary_mixed_gram} to $q=1$.  Let
    $X=(X_u)_{u\in\Fin d^K}$ and
    $Z=(Z_j)_{j\in\Fin d^1}$ be the resulting boundary families.  If
    $\rho\in\MN{D}$ has nonzero trace and
    $\mathcal K_\infty=\mathscr R(P_\rho)$, then
    \begin{align}
         & \left\langle X,
        (\Gamma^{\mathrm{tail}}_{K,l+1})^\dagger
        \Gamma^{\mathrm{left}}_{K+l,1}Z\right\rangle
        -\sum_{u,j}\left\langle
        U(A^jX_u),\mathcal K_\infty U(Z_jA^u)\right\rangle \\
         & \qquad=\sum_{u,j}\left\langle
        U(A^jX_u),(\mathcal K_l-\mathcal K_\infty)U(Z_jA^u)
        \right\rangle .
    \end{align}
    Here
    $\mathcal K_\infty U(Z_jA^u)
        =U((\tr\rho)^{-1}Z_jA^u\rho)$,
    so the limiting multiplication is on the right and is generally not the
    identity.
\end{theorem}
\begin{proof}
    \leanok
    \uses{thm:spectator_boundary_mixed_gram, thm:fixed_point_gram_metric_exact}
    Substitute the exact mixed-Gram formula and the right-multiplication action
    of $\mathcal K_\infty$, then distribute subtraction through the two finite
    sums and the second argument of the inner product.
\end{proof}

\begin{theorem}[Limiting metric inverse and virtual-map adjoints]%
    \label{thm:limiting_metric_inverse_virtual_adjoints}
    \lean{Matrix.inverse_gramReshuffle_fixedPointProj_frobeniusEquivEuclidean_apply}
    \lean{MPSTensor.tailVirtualMapES_adjoint_apply}
    \lean{MPSTensor.leftVirtualMapES_adjoint_apply}
    \leanok
    \uses{def:gram_reshuffle, thm:spectator_boundary_hilbert_factorization}
    For arbitrary virtual dimension, the adjoints of the tail and left virtual
    maps are
    \begin{align}
        (J^{\mathrm{tail}}_K)^\dagger(X_u)_u
         & =U\!\left(\sum_u X_u(A^u)^\dagger\right),          &
        (J^{\mathrm{left}}_L)^\dagger(Z_\tau)_\tau
         & =U\!\left(\sum_\tau(A^\tau)^\dagger Z_\tau\right).
    \end{align}
    If, in addition, $D>0$ and $\rho\in\MN{D}$ is positive definite, set
    $\mathcal K_\infty=\mathscr R(P_\rho)$ and
    $S_\infty=\mathcal K_\infty^{-1}$.  Then
    \begin{align}
        \mathcal K_\infty U(X)
         & =U\!\left((\tr\rho)^{-1}X\rho\right), &
        S_\infty U(X)
         & =U\!\left((\tr\rho)X\rho^{-1}\right).
    \end{align}
\end{theorem}
\begin{proof}
    \leanok
    \uses{thm:fixed_point_gram_metric_exact, thm:fixed_point_gram_metric_is_unit}
    For the tail adjoint, test against an arbitrary matrix $Y$ and use the
    Hilbert--Schmidt trace pairing fiberwise:
    \begin{align}
        \sum_u\langle U(YA^u),U(X_u)\rangle
        =\left\langle U(Y),
        U\!\left(\sum_uX_u(A^u)^\dagger\right)\right\rangle.
    \end{align}
    The left formula follows in the same way from
    $\langle U(A^\tau Y),U(Z_\tau)\rangle
        =\langle U(Y),U((A^\tau)^\dagger Z_\tau)\rangle$.

    Under the additional positive-definiteness hypothesis, the forward action
    of $\mathcal K_\infty$ is the exact limiting-metric formula and the
    operator is invertible.  Applying the forward formula to
    $U((\tr\rho)X\rho^{-1})$ cancels both scalar factors and the two right
    matrix factors, which gives the displayed inverse action.
\end{proof}

\begin{theorem}[Exact limiting-metric telescope for the C3 residual]%
    \label{thm:c3_projector_gram_error_telescope}
    \lean{MPSTensor.c3_virtualResidual_eq_centered_sub_gramErrors}
    \lean{MPSTensor.inner_limitingMixedGramIntertwining}
    \lean{MPSTensor.inner_c3_centeredVirtualResidual_eq_gramError}
    \lean{MPSTensor.c3_injectiveRangeProjector_residual_eq_centered_sub_gramErrors}
    \leanok
    \uses{thm:c3_spectator_projector_residual, def:ground_space_gram, def:gram_reshuffle}
    Let $\rho\in\MN{D}$ be positive definite and set
    $\mathcal K_\infty=\mathscr R(P_\rho)$ and
    $S_\infty=\mathcal K_\infty^{-1}$.  Write $J_{\rm t}$ and $J_{\rm l}$
    for the tail and left virtual maps, and write $S_n=\mathcal K_n^{-1}$
    whenever the length-$n$ boundary map is injective.  The virtual residual
    in the exact tail-after-left C3 projector factorization has the telescoping
    decomposition
    \begin{align}
         & M_l-\mathcal K_{l+1}^{\oplus K}J_{\rm t}S_{K+l+1}
        J_{\rm l}^\dagger\mathcal K_{K+l}^{\oplus 1}                \\
         & =\left(M_l-\mathcal K_\infty^{\oplus K}J_{\rm t}S_\infty
        J_{\rm l}^\dagger\mathcal K_\infty^{\oplus 1}\right)        \\
         & \quad-(\mathcal K_{l+1}^{\oplus K}
        -\mathcal K_\infty^{\oplus K})J_{\rm t}S_{K+l+1}
        J_{\rm l}^\dagger\mathcal K_{K+l}^{\oplus 1}                \\
         & \quad-\mathcal K_\infty^{\oplus K}J_{\rm t}
        (S_{K+l+1}-S_\infty)J_{\rm l}^\dagger
        \mathcal K_{K+l}^{\oplus 1}                                 \\
         & \quad-\mathcal K_\infty^{\oplus K}J_{\rm t}S_\infty
        J_{\rm l}^\dagger(\mathcal K_{K+l}^{\oplus 1}
        -\mathcal K_\infty^{\oplus 1}),
    \end{align}
    where
    $M_l=(\Gamma^{\rm tail}_{K,l+1})^\dagger
        \Gamma^{\rm left}_{K+l,1}$ and the superscript $\oplus$ denotes the
    corresponding fiberwise spectator operator.  Substituting this equality
    into the exact inverse-Gram factorization gives the same decomposition of
    the physical projector residual, sandwiched between the tail boundary map,
    the two local fiberwise inverse Grams, and the adjoint left boundary map.

    The limiting product in the first parenthesis has, without an additional
    hypothesis, the explicit limiting mixed pairing
    \begin{align}
         & \left\langle X,
        \mathcal K_\infty^{\oplus K}J_{\rm t}S_\infty
        J_{\rm l}^\dagger\mathcal K_\infty^{\oplus 1}Z
        \right\rangle                    \\
         & \qquad=\sum_{u,j}\left\langle
        U(A^jX_u),\mathcal K_\infty U(Z_jA^u)\right\rangle .
    \end{align}
    Consequently, pairing the first parenthesis against boundary families is
    exactly the same finite sum with
    $\mathcal K_l-\mathcal K_\infty$ in place of $\mathcal K_l$.

    These identities are reconstructed operator algebra.  They assert neither
    the numerical FNW/Nachtergaele contraction nor any identification of a
    physical ground projector with $P_\rho$.
\end{theorem}
\begin{proof}
    \leanok
    \uses{thm:c3_spectator_projector_residual, thm:spectator_boundary_direct_sum_gram, thm:spectator_boundary_mixed_gram_centered, thm:fixed_point_gram_metric_exact, thm:limiting_metric_inverse_virtual_adjoints}
    Add and subtract the product in which all three finite Gram factors are
    replaced by their limiting values, and telescope the resulting product
    from left to right.  Substitute the direct-sum Gram and inverse-Gram
    formulas into the exact projector factorization.

    It remains to identify the leading limiting product.  In Frobenius
    coordinates the limiting metric and its inverse act by
    \begin{align}
        \mathcal K_\infty U(X)
         & =U\!\left((\tr\rho)^{-1}X\rho\right), &
        S_\infty U(X)
         & =U\!\left((\tr\rho)X\rho^{-1}\right).
    \end{align}
    The adjoint left virtual map sends a family $W=(W_j)_j$ to
    $U(\sum_j(A^j)^\dagger W_j)$.  For
    $W_j=(\tr\rho)^{-1}Z_j\rho$, the right multiplication by $\rho$ and
    the factor $(\tr\rho)^{-1}$ supplied by the left limiting spectator
    metric cancel exactly through $S_\infty$.  Applying the tail virtual map and the
    tail limiting spectator metric leaves, in the fiber indexed by $u$,
    \begin{align}
        U\!\left((\tr\rho)^{-1}
        \sum_j (A^j)^\dagger Z_jA^u\rho\right).
    \end{align}
    Its inner product with $X_u$, followed by summation over $u$, is the
    displayed limiting mixed pairing by the Frobenius trace identity.  No
    transfer fixed-point equation is needed for this zeroth-order
    cancellation.  The centered mixed-Gram identity then turns the leading
    term into the pairing with $\mathcal K_l-\mathcal K_\infty$.
\end{proof}

\subsection{Martingale gap consequences}\label{subsec:spectral_gap_martingale}

\begin{lemma}[Weighted inner-product estimate]
    \label{lem:nachtergaele_weighted_inner_product}
    \lean{FrustrationFree.re_inner_le_weighted_norm_sq}
    \leanok
    Let $x$ and $y$ belong to a complex inner-product space. For every $c>0$,
    \begin{align}
        \re\langle x,y\rangle
         & \leq \frac{1}{2c}\lVert x\rVert^2
        +\frac{c}{2}\lVert y\rVert^2.
        \label{eq:nachtergaele_weighted_inner_product}
    \end{align}
\end{lemma}

\begin{proof}
    \leanok
    Cauchy--Schwarz bounds the left-hand side by
    $\lVert x\rVert\lVert y\rVert$. Apply the weighted arithmetic--geometric
    mean inequality to $\lVert y\rVert$ and $\lVert x\rVert$.
\end{proof}

\begin{lemma}[C3 projection-composition estimate]
    \label{lem:nachtergaele_c3_projection_composition}
    \lean{FrustrationFree.norm_sq_apply_projection_le_of_norm_comp_le}
    \leanok
    Let $P$ be an orthogonal projection on a complex inner-product space, and let
    $Q$ be a bounded linear operator from that space to another. If
    \begin{align}
        \lVert QP\rVert & \leq \epsilon,
    \end{align}
    then, for every $v$,
    \begin{align}
        \lVert QPv\rVert^2
         & \leq \epsilon^2\lVert Pv\rVert^2.
        \label{eq:nachtergaele_c3_projection_composition}
    \end{align}
    This includes $\epsilon=0$.
\end{lemma}

\begin{proof}
    \leanok
    Idempotency gives $QPv=QP(Pv)$. Apply the operator-norm estimate for
    $QP$ to $Pv$ and square the resulting nonnegative inequality. The
    hypothesis itself implies $\epsilon\geq0$.
\end{proof}

\begin{theorem}[Moving-window double counting]
    \label{thm:moving_window_sum_le}
    \lean{FrustrationFree.movingWindow_sum_le}
    \lean{FrustrationFree.movingWindow_sum_range_le}
    \lean{FrustrationFree.movingWindow_sum_Ico_le}
    \leanok
    Let $(a_m)_{m\geq0}$ be a nonnegative real sequence. For all $l,N\in\N$,
    \begin{align}
        \sum_{n=0}^{N-1}\sum_{m=\max\{0,n-l\}}^{n}a_m
         & \leq (l+1)\sum_{m=0}^{N-1}a_m.
        \label{eq:moving_window_sum_le}
    \end{align}
    The same bound holds after omitting the terms with $n<l$. More generally,
    for a threshold $n_0\in\N$ only the coefficients indexed by
    $m\geq\max\{0,n_0-l\}$ occur, so
    \begin{align}
        \sum_{n=n_0}^{N-1}\sum_{m=\max\{0,n-l\}}^{n}a_m
         & \leq (l+1)\sum_{m=\max\{0,n_0-l\}}^{N-1}a_m.
        \label{eq:moving_window_sum_Ico_le}
    \end{align}
\end{theorem}

\begin{proof}
    \leanok
    Write $k=n-m$. Each pair in the double sum has $0\leq k\leq l$.
    For fixed $k$, the map $n\mapsto n-k$ is injective and has image in
    $\{0,\ldots,N-1\}$. Nonnegativity bounds its contribution by
    $\sum_{m=0}^{N-1}a_m$. Summing over the $l+1$ values of $k$ proves
    (\ref{eq:moving_window_sum_le}). This is the counting step in the proof of
    Nachtergaele's Theorem~2.1(i), between Equation~(Enpsi2) and the following
    summed estimate: for $a_m=\lVert E_m\psi\rVert^2$, each term occurs in at
    most $l+1$ consecutive windows
    \cite[Proof of Theorem~2.1(i)]{Nachtergaele1996Spectral}. For
    (\ref{eq:moving_window_sum_Ico_le}), replace $a_m$ by the sequence that
    vanishes below $\max\{0,n_0-l\}$ and agrees with $a_m$ above it. Every
    window with $n\geq n_0$ is unchanged, and the previous bound applied to
    the replacement sequence gives the claim.
\end{proof}

% Source: References/cond-mat_9410110/main.tex, conditions C1--C3 at
% lines 1030--1094, Theorem 2.1(i) at lines 1119--1130, and its proof at
% lines 1195--1259. The statement below isolates the estimate above a chosen
% threshold and explicitly records how its C1 range relates to source C1.
\begin{theorem}[Threshold C1--C3 energy estimate]
    \label{thm:nachtergaele_c1_c3_threshold_energy_estimate}
    \lean{FrustrationFree.NestedGroundProjections.energy_lower_bound_of_nachtergaele_c1_c3_threshold}
    \leanok
    Let $(G_n)_{n\geq0}$ be nested orthogonal projections on a
    finite-dimensional complex Hilbert space $\mathcal H$, put
    $E_n=G_n-G_{n+1}$, and assume $G_0=\mathbf 1$. Fix $l,N,n_0\in\N$ with
    $n_0\leq N$, and let $Q_n$ be
    orthogonal projections for $n_0\leq n<N$. Suppose $E_mQ_n=Q_nE_m$
    whenever $m<n-l$ or $n<m$. Let $H_n$ be the local Hamiltonians and $H$
    the full Hamiltonian. Assume the quadratic-form version of C1 on the same
    range, for every $x\in\mathcal H$,
    \begin{align}
        0
         & \leq\sum_{n=n_0}^{N-1}\re\langle H_nx,x\rangle
        \leq d_{l+1}\re\langle Hx,x\rangle,
        \label{eq:nachtergaele_c1_threshold_form}
    \end{align}
    the quadratic-form version of C2 for every $x\in\mathcal H$ and
    $n_0\leq n<N$,
    \begin{align}
        \gamma_{l+1}\lVert(\mathbf 1-Q_n)x\rVert^2
         & \leq\re\langle H_nx,x\rangle,
        \label{eq:nachtergaele_c2_threshold_form}
    \end{align}
    and the literal condition C3 for $n_0\leq n<N$,
    \begin{align}
        \lVert Q_nE_n\rVert
         & \leq\epsilon_l,
         & 0\leq\epsilon_l
         & <\frac{1}{\sqrt{l+1}}.
        \label{eq:nachtergaele_c3_threshold_form}
    \end{align}
    If $\gamma_{l+1}>0$, $d_{l+1}>0$, and $v\perp\operatorname{ran}(G_N)$,
    then
    \begin{align}
        \frac{\gamma_{l+1}}{d_{l+1}}
        \left(1-\epsilon_l\sqrt{l+1}\right)^2
        \sum_{n=n_0}^{N-1}\lVert E_nv\rVert^2
         & \leq
        \re\langle Hv,v\rangle
        \nonumber \\
         & \quad
        +\frac{\gamma_{l+1}}{d_{l+1}}
        \left(1-\epsilon_l\sqrt{l+1}\right)\epsilon_l\sqrt{l+1}
        \sum_{m=\max\{0,n_0-l\}}^{n_0-1}\lVert E_mv\rVert^2.
        \label{eq:nachtergaele_c1_c3_threshold_energy_estimate}
    \end{align}
    Conditions C2 and C3 have the source forms, with their onset represented by
    $n_0$. The displayed C1 inequality follows from source C1 when
    $l\leq n_0$ and the omitted local terms are nonnegative. Thus a
    source-derived application chooses $n_0$ above both $l$ and the C2--C3
    onset. If $n_0<l$, the displayed C1 inequality is instead an independent
    stronger assumption. The coefficient of the martingale mass above the
    threshold is exactly the one in Theorem~2.1(i) of
    \cite{Nachtergaele1996Spectral}. The correction is an explicit upper bound
    for the contribution of the at most $l$ martingale differences immediately
    below the threshold; it uses multiplicity $l+1$ uniformly. At $n_0=0$ the
    correction is an empty sum.
\end{theorem}

\begin{proof}
    \leanok
    \uses{thm:truncated-martingale-difference-reconstruction,
        thm:martingale-difference-pythagorean-identity,
        thm:outside-window-martingale-cancellation,
        lem:nachtergaele_weighted_inner_product,
        lem:nachtergaele_c3_projection_composition,
        thm:moving_window_sum_le}
    Write $P=\sum_{n=n_0}^{N-1}\lVert E_nv\rVert^2$ and
    $S=\sum_{m=\max\{0,n_0-l\}}^{n_0-1}\lVert E_mv\rVert^2$. The martingale
    resolution $v=\sum_{n=0}^{N-1}E_nv$ and outside-window cancellation
    reduce the identity labelled \textup{Enpsi} at index $n$ to the active
    window $m\in[n-l,n]$; this uses only $n<N$, so it is available for every
    $n\geq n_0$. Apply the weighted inner-product estimate with positive
    parameters $c_1,c_2$, then C2 and C3, and sum over $n_0\leq n<N$. The
    martingale differences are mutually orthogonal, so each window norm
    square splits into its summands, and
    (\ref{eq:moving_window_sum_Ico_le}) bounds the resulting double sum by
    $(l+1)(S+P)$. With C1 this gives
    \begin{align}
        \left(2-c_1-\frac{\epsilon_l^2}{c_2}-(l+1)c_2\right)P
         & \leq
        \frac{d_{l+1}}{c_1\gamma_{l+1}}\re\langle Hv,v\rangle
        +(l+1)c_2S.
    \end{align}
    For $\epsilon_l>0$, take $c_1=1-\epsilon_l\sqrt{l+1}$ and
    $c_2=\epsilon_l/\sqrt{l+1}$, so that
    $\epsilon_l^2/c_2=(l+1)c_2=\epsilon_l\sqrt{l+1}$ and the left factor is
    $1-\epsilon_l\sqrt{l+1}$. Multiplying by
    $c_1\gamma_{l+1}/d_{l+1}$ gives
    (\ref{eq:nachtergaele_c1_c3_threshold_energy_estimate}). If
    $\epsilon_l=0$, then $c_2=0$ is not an admissible weighted parameter. In
    that case C3 instead gives $Q_nE_n=0$, so the second inner product in
    \textup{Enpsi} vanishes; applying only the first weighted estimate with
    $c_1=1$ yields $P\leq(d_{l+1}/\gamma_{l+1})\re\langle Hv,v\rangle$,
    which is the same formula because the correction term vanishes. This
    repair is documented in \cite{gap:cpgsv21_martingale_overlap}.
\end{proof}

% Source: References/cond-mat_9410110/main.tex, the definition of E_n at
% lines 1068--1074. The vanishing below a fixed level is what the printed
% proof at lines 1195--1259 needs at the indices excluded by the lower
% endpoints of C2 and C3.
\begin{lemma}[Martingale differences below a fixed level]
    \label{lem:martingale_difference_vanishing_below_level}
    \lean{FrustrationFree.NestedGroundProjections.martingaleDifference_apply_eq_zero_of_lt}
    \leanok
    Let $(G_n)_{n\geq0}$ be nested orthogonal projections on a complex
    Hilbert space, and put $E_n=G_n-G_{n+1}$. If $G_{n_0}v=v$ for some
    $n_0$, then $E_nv=0$ for every $n<n_0$.
\end{lemma}

\begin{proof}
    \leanok
    For $m\leq n_0$ the ranges are nested, so $G_mG_{n_0}=G_{n_0}$ and hence
    $G_mv=G_mG_{n_0}v=G_{n_0}v=v$. For $n<n_0$ this applies to both $n$ and
    $n+1$, so $E_nv=G_nv-G_{n+1}v=v-v=0$.
\end{proof}

% Source: References/cond-mat_9410110/main.tex, conditions C1--C3 at
% lines 1030--1094, Theorem 2.1(i) at lines 1119--1130, and its proof at
% lines 1195--1259. The vanishing hypothesis below the threshold is the
% minimal repair of the printed statement; the source defect it repairs is
% documented in
% docs/paper-gaps/nachtergaele96_theorem_2_1_lower_endpoint.tex.
\begin{theorem}[C1--C3 energy estimate with no mass below the threshold]
    \label{thm:nachtergaele_c1_c3_below_threshold_energy_estimate}
    \lean{FrustrationFree.NestedGroundProjections.energy_lower_bound_of_nachtergaele_c1_c3_of_martingaleDifference_below_eq_zero}
    \leanok
    \uses{thm:nachtergaele_c1_c3_threshold_energy_estimate}
    Assume the hypotheses of
    Theorem~\ref{thm:nachtergaele_c1_c3_threshold_energy_estimate}, except
    that $l\leq n_0$ is assumed and
    (\ref{eq:nachtergaele_c1_threshold_form}) is replaced by condition C1 on
    the range $l\leq n<M$ of its own window sum at the two volumes at which
    the source asserts it: the full inequality
    $0\leq\sum_{n=l}^{N-1}\re\langle H_nx,x\rangle
        \leq d_{l+1}\re\langle Hx,x\rangle$ at $M=N$, and its non-negativity
    clause $0\leq\sum_{n=l}^{n_0-1}\re\langle H_nx,x\rangle$ at $M=n_0$.
    Assume in addition that $E_nv=0$ for every $n<n_0$. Then
    \begin{align}
        \frac{\gamma_{l+1}}{d_{l+1}}
        \left(1-\epsilon_l\sqrt{l+1}\right)^2\lVert v\rVert^2
         & \leq\re\langle Hv,v\rangle.
        \label{eq:nachtergaele_c1_c3_below_threshold_energy_estimate}
    \end{align}
    This is the conclusion of Theorem~2.1(i) of
    \cite{Nachtergaele1996Spectral} with each of C1, C2, and C3 imposed on the
    range at which the source imposes it. The added
    vanishing hypothesis is the minimal repair of that theorem: the two
    below-threshold quantities in
    (\ref{eq:nachtergaele_c1_c3_threshold_energy_estimate}) enter
    (\ref{eq:nachtergaele_c1_c3_below_threshold_energy_estimate}) with
    non-negative weights, so both have to vanish, and the vanishing of the
    omitted martingale mass forces the vanishing of the window correction.
    Without the added hypothesis the statement is false, by
    Theorem~\ref{thm:nachtergaele_threshold_conditions_insufficient}. The
    source defect and the repair are recorded in
    \cite{gap:nachtergaele96_theorem_2_1_lower_endpoint}.
\end{theorem}

\begin{proof}
    \leanok
    \uses{thm:nachtergaele_c1_c3_threshold_energy_estimate,
        thm:martingale-difference-final-complement-norm-decomposition}
    Condition C1 on the threshold range $n_0\leq n<N$ follows from the
    assumed form: its upper bound by subtracting from the sum over
    $l\leq n<N$ the initial sum over $l\leq n<n_0$, which is non-negative by
    C1 at the volume $\Lambda_{n_0}$, and its non-negativity clause termwise
    from C2, each of whose left-hand sides is non-negative. The window
    correction in
    (\ref{eq:nachtergaele_c1_c3_threshold_energy_estimate}) is a sum of
    squared norms of martingale differences with index below $n_0$, hence
    zero. The Pythagorean decomposition of a vector orthogonal to
    $\operatorname{ran}(G_N)$ splits $\lVert v\rVert^2$ into the sum over
    $n<n_0$, which is zero for the same reason, and the martingale mass
    $\sum_{n=n_0}^{N-1}\lVert E_nv\rVert^2$. Hence the latter equals
    $\lVert v\rVert^2$, and
    (\ref{eq:nachtergaele_c1_c3_threshold_energy_estimate}) becomes
    (\ref{eq:nachtergaele_c1_c3_below_threshold_energy_estimate}).
\end{proof}

% Source: References/cond-mat_9410110/main.tex, conditions C1--C3 at
% lines 1030--1094 and Theorem 2.1(i) at lines 1119--1130. The witness is the
% one-site chain with a weakened leftmost interaction, restricted to the span
% of its ground state and its single excitations; see
% docs/paper-gaps/nachtergaele96_theorem_2_1_lower_endpoint.tex.
\begin{theorem}[The source conditions on their own ranges do not give the
        estimate]
    \label{thm:nachtergaele_threshold_conditions_insufficient}
    \lean{FrustrationFree.not_unrestrictedNachtergaeleEstimate}
    \leanok
    \uses{thm:nachtergaele_c1_c3_below_threshold_energy_estimate}
    In the notation of
    Theorem~\ref{thm:nachtergaele_c1_c3_below_threshold_energy_estimate}: for
    every $k$ there are, on a complex Hilbert space of dimension $k+2$, nested
    orthogonal projections $(G_n)_{n\geq0}$ with $G_0=\mathbf 1$,
    $G_{N+1}=0$, and $G_n\neq0$ for every $n\leq N$, nonzero orthogonal
    projections $Q_n$, local operators $H_n$, volume Hamiltonians
    $H_{\Lambda_M}$, and a vector $v\perp\operatorname{ran}(G_N)$ such that,
    with $l=0$, $n_0=1$, $N=k+1$, $\gamma_1=d_1=1$, and $\epsilon_0=0$, the
    required commutations hold, condition C1 holds with the single constant
    $d_1$ at every volume $\Lambda_M$ on the range $l\leq n<M$ carried by its
    own window sum, the gap inequality of condition C2 and condition C3 hold
    on the range $n_0\leq n<N$ from which they are imposed, every ground space
    is the kernel of its own Hamiltonian, that is
    $\ker H_{\Lambda_M}=\operatorname{ran}(G_M)$ for $M\leq N$ and
    $\ker H_n=\operatorname{ran}(Q_n)$ for every window inside $\Lambda_N$,
    that is for $0\leq n<N$, where the $Q_n$ are also nonzero orthogonal
    projections commuting with the martingale differences outside their own
    windows, and
    \begin{align}
        \re\langle H_{\Lambda_N}v,v\rangle
         & =\tfrac12\lVert v\rVert^2
        <\frac{\gamma_1}{d_1}
        \left(1-\epsilon_0\right)^2\lVert v\rVert^2.
        \label{eq:nachtergaele_threshold_conditions_insufficient}
    \end{align}
    Hence condition C1 at every volume on the range of its own window sum
    together with conditions C2 and C3 on $n_0\leq n<N$, the ranges at which
    \cite{Nachtergaele1996Spectral} imposes them, do not imply the conclusion
    of its Theorem~2.1(i), even when every ground space is the kernel of its
    Hamiltonian and is nontrivial, as condition C2 also requires, and even
    under the conventions the source attaches to the extreme volumes.
\end{theorem}

\begin{proof}
    \leanok
    Fix an orthonormal basis $e_0,\ldots,e_{k+1}$ and let every operator be
    diagonal in it. The first $N=k+1$ basis vectors carry the single
    excitations of a chain of $N$ sites and the last carries its ground state.
    Let $G_n$ have weight $1$ at $e_i$ when $n\leq i$ and weight $0$
    otherwise, so that $G_0=\mathbf 1$, $G_{N+1}=0$ because no basis vector
    has index at least $N+1$, every $G_n$ with $n\leq N$ is nonzero
    because it retains $e_{k+1}$, $G_N$ is the projection onto $e_{k+1}$, and
    $E_n$ is the projection onto $e_n$ for $n<N$. Let $c_0=\tfrac12$, let
    $c_i=1$ for $1\leq i\leq k$, and let $c_{k+1}=0$; let $H_{\Lambda_M}$ have
    weight $c_i$ at $e_i$ when $i<M$ and weight $0$ otherwise, let $H_n$ have
    weight $c_n$ at $e_n$ and weight $0$
    elsewhere, and let $Q_n=\mathbf 1-E_n$, which is nonzero because the space
    has dimension at least two. All operators are diagonal, so the required
    commutations hold, and $Q_nE_n=0$ gives condition C3 with $\epsilon_0=0$.
    Summing $H_n$ over the C1 range $0\leq n<M$ gives $H_{\Lambda_M}$ exactly,
    at every $M$, and the weights $c_i$ are nonnegative, so condition C1 holds
    at every volume with the single constant $d_1=1$ and with equality in its
    upper bound. Since $c_i\neq0$ for $i\leq k$, the kernel of
    $H_{\Lambda_M}$ is spanned by the $e_i$ with $i\geq M$ whenever $M\leq N$,
    which is $\operatorname{ran}(G_M)$, and the kernel of $H_n$ is spanned by
    the $e_i$ with $i\neq n$ whenever $0\leq n<N$, which is
    $\operatorname{ran}(Q_n)$ and is nonzero. For $1\leq n<N$ one has $c_n=1$,
    so
    $\lVert(\mathbf 1-Q_n)x\rVert^2=\lVert E_nx\rVert^2$ equals
    $\re\langle H_nx,x\rangle$, which is condition C2 with $\gamma_1=1$.
    Take $v=e_0$. Then $v\perp\operatorname{ran}(G_N)$, the line through
    $e_{k+1}$, and $\re\langle H_{\Lambda_N}v,v\rangle=c_0=\tfrac12$, while
    $\lVert v\rVert^2=1$, which
    is (\ref{eq:nachtergaele_threshold_conditions_insufficient}). The whole
    martingale mass of $v$ sits at the index $n=0$, which C1 covers while C2
    and C3 exclude it, and it is at that index that C2 fails.
\end{proof}

% Source: References/cond-mat_9410110/main.tex, conditions C1--C3 at
% lines 1030--1094, Theorem 2.1(i) at lines 1119--1130, and its proof at
% lines 1195--1259. The lower-endpoint restriction is documented in
% docs/paper-gaps/cpgsv21_martingale_overlap.tex.
\begin{theorem}[Full-range C1--C3 energy estimate]
    \label{thm:nachtergaele_c1_c3_full_range_energy_estimate}
    \lean{FrustrationFree.NestedGroundProjections.energy_lower_bound_of_nachtergaele_c1_c3_full_range}
    \leanok
    Let $(G_n)_{n\geq0}$ be nested orthogonal projections on a
    finite-dimensional complex Hilbert space $\mathcal H$, put
    $E_n=G_n-G_{n+1}$, and assume $G_0=\mathbf 1$. Fix $l,N\in\N$, and let
    $Q_n$ be orthogonal projections
    for $0\leq n<N$. Suppose $E_mQ_n=Q_nE_m$ whenever $m<n-l$ or $n<m$.
    Let $H_n$ be the local Hamiltonians and $H$ the full Hamiltonian.
    Assume the quadratic-form version of C1 for every $x\in\mathcal H$,
    \begin{align}
        0
         & \leq\sum_{n=0}^{N-1}\re\langle H_nx,x\rangle
        \leq d_{l+1}\re\langle Hx,x\rangle,
        \label{eq:nachtergaele_c1_form}
    \end{align}
    and the quadratic-form version of C2 for every $x\in\mathcal H$ and
    $0\leq n<N$,
    \begin{align}
        \gamma_{l+1}\lVert(\mathbf 1-Q_n)x\rVert^2
         & \leq\re\langle H_nx,x\rangle,
        \label{eq:nachtergaele_c2_form}
    \end{align}
    and the literal condition C3,
    \begin{align}
        \lVert Q_nE_n\rVert
         & \leq\epsilon_l,
         & 0\leq\epsilon_l
         & <\frac{1}{\sqrt{l+1}}.
        \label{eq:nachtergaele_c3_form}
    \end{align}
    If $\gamma_{l+1}>0$, $d_{l+1}>0$, and
    $v\perp\operatorname{ran}(G_N)$, then
    \begin{align}
        \re\langle Hv,v\rangle
         & \geq
        \frac{\gamma_{l+1}}{d_{l+1}}
        \left(1-\epsilon_l\sqrt{l+1}\right)^2\lVert v\rVert^2.
        \label{eq:nachtergaele_c1_c3_full_range_energy_estimate}
    \end{align}
    This full-range statement is stronger than the unrestricted source
    theorem in general. It assumes C1, C2, and C3 for every $0\leq n<N$,
    whereas the source imposes lower thresholds on these conditions. Those
    thresholds cannot simply be dropped: without separate control of the
    initial volumes the source hypotheses imply neither the assumptions above
    nor the conclusion, by
    Theorem~\ref{thm:nachtergaele_threshold_conditions_insufficient}. The
    coefficient is exactly the one in Theorem~2.1(i) of
    \cite{Nachtergaele1996Spectral}: $\gamma_{l+1}/d_{l+1}$ times
    $(1-\epsilon_l\sqrt{l+1})^2$. This is
    Theorem~\ref{thm:nachtergaele_c1_c3_threshold_energy_estimate} at
    threshold $n_0=0$, where nothing lies below the threshold.
\end{theorem}

\begin{proof}
    \leanok
    \uses{thm:nachtergaele_c1_c3_threshold_energy_estimate,
        thm:martingale-difference-final-complement-norm-decomposition}
    Apply
    Theorem~\ref{thm:nachtergaele_c1_c3_threshold_energy_estimate} with
    $n_0=0$ and the given window length $l$. Its window correction is a sum
    over an empty range, hence zero, and its martingale mass over
    $0\leq n<N$ is $\lVert v\rVert^2$ by the Pythagorean decomposition of a
    vector orthogonal to $\operatorname{ran}(G_N)$. What remains is
    (\ref{eq:nachtergaele_c1_c3_full_range_energy_estimate}).
\end{proof}

\begin{theorem}[Energy form to norm bound]
    \label{thm:energy_form_to_norm_bound}
    \lean{FrustrationFree.norm_lower_bound_of_energy_lower_bound}
    \leanok
    Let $H$ be a linear operator on a complex Hilbert space. If
    \begin{align}
        \kappa\lVert v\rVert^2
         & \leq\re\langle Hv,v\rangle
    \end{align}
    for every $v\in(\ker H)^\perp$, then
    $\kappa\lVert v\rVert\leq\lVert Hv\rVert$ on the same subspace.
\end{theorem}

\begin{proof}
    \leanok
    Cauchy--Schwarz gives
    \begin{align}
        \kappa\lVert v\rVert^2
         & \leq\re\langle Hv,v\rangle
        \leq\lVert Hv\rVert\lVert v\rVert.
    \end{align}
    The claim is immediate for $v=0$; otherwise cancel $\lVert v\rVert$.
\end{proof}

\begin{theorem}[Norm gap from the full-range C1--C3 estimate]
    \label{thm:nachtergaele_c1_c3_full_range_norm_gap}
    \lean{FrustrationFree.NestedGroundProjections.norm_lower_bound_of_nachtergaele_c1_c3_full_range}
    \leanok
    In the setting of
    Theorem~\ref{thm:nachtergaele_c1_c3_full_range_energy_estimate}, suppose
    that $H$ is
    positive and $\operatorname{ran}(G_N)=\ker H$. Then every
    $v\in(\ker H)^\perp$ satisfies
    \begin{align}
        \frac{\gamma_{l+1}}{d_{l+1}}
        \left(1-\epsilon_l\sqrt{l+1}\right)^2\lVert v\rVert
         & \leq\lVert Hv\rVert.
        \label{eq:nachtergaele_c1_c3_full_range_norm_gap}
    \end{align}
\end{theorem}

\begin{proof}
    \leanok
    \uses{thm:nachtergaele_c1_c3_full_range_energy_estimate,
        thm:energy_form_to_norm_bound,
        lem:quadratic_form_to_norm_bound}
    Denote the coefficient in
    (\ref{eq:nachtergaele_c1_c3_full_range_norm_gap}) by $\kappa$.
    For an
    arbitrary $x$, remove its orthogonal projection onto $\ker H$ and call the
    remainder $w$. Then $Hx=Hw$ and
    $\re\langle Hx,x\rangle=\re\langle Hw,w\rangle$. The energy estimate and
    Theorem~\ref{thm:energy_form_to_norm_bound} give
    $\kappa\lVert w\rVert\leq\lVert Hx\rVert$. Hence Cauchy--Schwarz gives
    \begin{align}
        \kappa\,\re\langle Hx,x\rangle
         & \leq \kappa\lVert Hx\rVert\lVert w\rVert
        \leq \lVert Hx\rVert^2.
    \end{align}
    Thus the hypothesis of
    Lemma~\ref{lem:quadratic_form_to_norm_bound} holds globally. The
    finite-dimensional quadratic-form criterion then yields the stated norm
    bound.
\end{proof}

\begin{theorem}[Martingale criterion for a finite family of projections]
    \label{thm:martingale_criterion}
    \lean{MPSTensor.spectralGap_of_martingale_anticommutator_rowCol_of_finiteDimensional}
    \leanok
    \uses{lem:anticommutator_projection_rowsum_reduction, lem:quadratic_form_to_norm_bound}
    Let $P_i$ be a finite family of symmetric projections on a
    finite-dimensional complex Hilbert space, let $H=\sum_iP_i$, and let
    $0<\gamma\le1$.  Suppose that a real coefficient matrix $(c_{ij})$ has row
    and column sums at most one away from the diagonal, and that, for $i\ne j$
    and every $v$,
    \begin{align}
         & \re\langle P_i v,P_jv\rangle
        +\re\langle P_jv,P_iv\rangle    \\
         & \qquad\ge
        -(1-\gamma)c_{ij}
        \bigl(\re\langle P_iv,v\rangle
        +\re\langle P_jv,v\rangle\bigr).
        \label{eq:ph_martingale_condition}
    \end{align}
    Then
    \begin{align}
        \gamma\,\re\langle Hv,v\rangle
         & \le \re\langle Hv,Hv\rangle
        \notag
    \end{align}
    for every $v$.  Consequently, for every
    $v\in(\ker H)^\perp$,
    \begin{align}
        \gamma\lVert v\rVert & \le\lVert Hv\rVert.
        \notag
    \end{align}
\end{theorem}

\begin{proof}
    \leanok
    \uses{lem:anticommutator_projection_rowsum_reduction, lem:quadratic_form_to_norm_bound}
    The row-and-column reduction gives the quadratic-form inequality.  The sum
    of symmetric projections is positive, so the abstract martingale criterion
    converts that inequality into the norm lower bound on $(\ker H)^\perp$.
\end{proof}

\begin{theorem}[Reduced ground-space overlap bounds the excitation anticommutator]%
    \label{thm:friedrichs_ground_space_excitation_anticommutator}
    \lean{Submodule.re_inner_anticommutator_starProjection_orthogonal_ge_neg_of_friedrichs_bound}
    \leanok
    \uses{def:friedrichs_overlap_bound}
    Let $U$ and $V$ be ground-space ranges in a finite-dimensional complex
    Hilbert space, let $W=U\cap V$, and let $q_U$ and $q_V$ be the orthogonal
    projectors onto $U$ and $V$. Suppose that $\eta\ge0$ and
    \begin{align}
        |\langle x,y\rangle|
         & \le \eta\,\|x\|\,\|y\|
        \label{eq:ph_reduced_friedrichs_overlap}
    \end{align}
    whenever $x\in U\cap W^\perp$ and $y\in V\cap W^\perp$. Define the
    complementary excitation projections by $h_U=\Id-q_U$ and
    $h_V=\Id-q_V$. Then, for every vector $v$,
    \begin{align}
        \re\langle h_Uv,h_Vv\rangle
        +\re\langle h_Vv,h_Uv\rangle
         & \ge
        -\eta\bigl(
        \re\langle h_Uv,v\rangle
        +\re\langle h_Vv,v\rangle
        \bigr).
        \label{eq:ph_friedrichs_excitation_anticommutator}
    \end{align}
\end{theorem}

\begin{proof}
    \leanok
    Remove the component of $v$ in the common intersection $W$; both
    excitation projections are unchanged. If $\eta\le1$, apply
    (\ref{eq:ph_reduced_friedrichs_overlap}) to the reduced ground-space
    components $q_Uv$ and $q_Vv$, and combine the resulting principal-angle
    estimate with the orthogonal decompositions
    $v=q_Uv+h_Uv=q_Vv+h_Vv$. If $\eta>1$, the same conclusion follows from
    the universal overlap bound with coefficient $1$ and the non-negativity
    of the two diagonal excitation forms.
\end{proof}

\subsection{The compatible open-chain Hamiltonian}%
\label{subsec:compatible_open_chain_hamiltonian}

For an interval, the open-chain parent Hamiltonian contains only local terms
whose supports lie entirely within the
interval~\cite[Equation~(3.12)]{Nachtergaele1996Spectral}. The corresponding
zero-energy space coincides with the open-chain ground space. Each compatible
open Hamiltonian is positive as a sum of orthogonal projections. Conditions
(C1) and (C2) hold across the full martingale range: before a complete
interaction window fits, the suffix Hamiltonian vanishes, and on larger
intervals it equals the orthogonal projection onto the local excited space.
Combined with condition (C3), these identities specialize
Theorem~\ref{thm:nachtergaele_c1_c3_full_range_norm_gap} with
$d_{l+1}=\gamma_{l+1}=1$ and overlap bound $\epsilon_l$.

\begin{definition}[Compatible open-chain parent Hamiltonian]
    \label{def:open_parent_hamiltonian_es}
    \lean{MPSTensor.NonwrappingStart}
    \lean{MPSTensor.openParentHamiltonianES}
    \leanok
    \uses{rem:euclidean_local_projector_ingredients}
    For $L,N\in\N$, let
    \begin{align}
        I_{L,N} & =\{i\in\{0,\ldots,N-1\}\mid i+L\leq N\}.
    \end{align}
    The compatible length-$N$ open-chain parent Hamiltonian of interaction
    length $L$ is
    \begin{align}
        H^{\mathrm{open}}_{N,L}(A)
         & =\sum_{i\in I_{L,N}}h_{i,\mathrm{ES}}(A,L).
        \label{eq:open_parent_hamiltonian_es}
    \end{align}
    Thus no summand crosses the two endpoints of the interval. For
    $L=l+1$, the starting sites are $0,\ldots,N-l-1$, in agreement with
    \cite[Equation~(3.12)]{Nachtergaele1996Spectral}.
\end{definition}

\begin{theorem}[Full-window compatible interaction]
    \label{thm:open_parent_hamiltonian_full_window}
    \lean{MPSTensor.openParentHamiltonianES_self_eq_parentInteractionES}
    \leanok
    \uses{def:open_parent_hamiltonian_es, rem:euclidean_local_projector_ingredients}
    If $L>0$, then
    \begin{align}
        H^{\mathrm{open}}_{L,L}(A)
         & =P_L^{\mathrm{ES}}(A).
        \label{eq:open_parent_hamiltonian_full_window}
    \end{align}
\end{theorem}

\begin{proof}
    \leanok
    \uses{def:open_parent_hamiltonian_es, rem:cyclic_window_operators, def:extract_window, rem:euclidean_local_projector_ingredients}
    The interaction has one window:
    \begin{align}
        I_{L,L} & =\{0\}.
    \end{align}
    The outside region is empty, so it has a unique configuration $\tau$.
    Cyclic restriction to the sole window is the identity:
    \begin{align}
        R_{0,\tau}v & =v.
    \end{align}
    Likewise,
    \begin{align}
        \sigma|_{[0,L-1]} & =\sigma.
    \end{align}
    Therefore
    \begin{align}
        H^{\mathrm{open}}_{L,L}(A)
         & =\sum_{i\in I_{L,L}}h_{i,\mathrm{ES}}(A,L)
        =h_{0,\mathrm{ES}}(A,L)
        =P_L^{\mathrm{ES}}(A).
    \end{align}
\end{proof}

\begin{theorem}[Unit gap at the full-window volume]
    \label{thm:open_parent_hamiltonian_full_window_unit_gap}
    \lean{MPSTensor.openParentHamiltonianES_self_norm_eq_of_mem_orthogonal_ker,
        MPSTensor.openParentHamiltonianES_self_unit_gap}
    \leanok
    \uses{thm:open_parent_hamiltonian_full_window}
    If $L>0$ and
    $v\in(\ker H^{\mathrm{open}}_{L,L}(A))^\perp$, then
    \begin{align}
        \lVert H^{\mathrm{open}}_{L,L}(A)v\rVert
         & =\lVert v\rVert.
        \label{eq:open_parent_hamiltonian_full_window_norm}
    \end{align}
    In particular,
    \begin{align}
        \lVert v\rVert
         & \leq\lVert H^{\mathrm{open}}_{L,L}(A)v\rVert.
    \end{align}
    Thus $\gamma_{l+1}=1$ at the full-window volume $N=L=l+1$ in
    \cite[Theorem~2.1(i)]{Nachtergaele1996Spectral}. This neither asserts that
    the excited subspace is nonzero nor supplies a gap for any $N>L$.
\end{theorem}

\begin{proof}
    \leanok
    \uses{thm:open_parent_hamiltonian_full_window, lem:parent_interaction_es_kernel}
    The full-window identity and the local kernel formula give
    \begin{align}
        \ker H^{\mathrm{open}}_{L,L}(A)
         & =\ker P_L^{\mathrm{ES}}(A)=G_L^{\mathrm{ES}}(A).
    \end{align}
    Hence $v$ belongs to $(G_L^{\mathrm{ES}}(A))^\perp$, where the orthogonal
    projector $P_L^{\mathrm{ES}}(A)$ acts as the identity. Thus
    \begin{align}
        \left\lVert H^{\mathrm{open}}_{L,L}(A)v\right\rVert
         & =\left\lVert P_L^{\mathrm{ES}}(A)v\right\rVert
        =\lVert v\rVert.
    \end{align}
    Consequently,
    \begin{align}
        \lVert v\rVert
         & \leq\left\lVert H^{\mathrm{open}}_{L,L}(A)v\right\rVert.
    \end{align}
\end{proof}

\begin{lemma}[Positivity of the compatible open-chain Hamiltonian]
    \label{lem:open_parent_hamiltonian_es_positive}
    \lean{MPSTensor.openParentHamiltonianES_isPositive}
    \leanok
    \uses{def:open_parent_hamiltonian_es, lem:transported_es_positive}
    For every $A$, $L$, and $N$, the operator
    $H^{\mathrm{open}}_{N,L}(A)$ is positive.
\end{lemma}

\begin{proof}
    \leanok
    Each $h_{i,\mathrm{ES}}(A,L)$ is an orthogonal projection and hence is
    positive. Positivity is preserved by finite sums.
\end{proof}

\begin{definition}[Compatible suffix Hamiltonian]
    \label{def:open_suffix_parent_hamiltonian_es}
    \lean{MPSTensor.openSuffixParentHamiltonianES}
    \leanok
    \uses{def:open_parent_hamiltonian_es}
    For $L,l,n,N\in\N$, the Hamiltonian on the last $l$ sites of the prefix
    $[0,n)$, regarded as an operator on the $N$-site chain, is
    \begin{align}
        H^{\mathrm{suffix}}_{N,L;l,n}(A)
         & =\sum_{\substack{i\in I_{L,N} \\ n-l\leq i,\ i+L\leq n}}
        h_{i,\mathrm{ES}}(A,L).
        \label{eq:open_suffix_parent_hamiltonian_es}
    \end{align}
    Its local windows lie in the nonwrapping interval $[n-l,n)$.
\end{definition}

% Source: References/cond-mat_9410110/main.tex, condition C1,
% lines 1030--1041.  The vanishing terms supply the indices below the
% source lower endpoint when L=l+1.
\begin{theorem}[Vanishing suffix before a complete window]
    \label{thm:open_suffix_parent_hamiltonian_zero_before_window}
    \lean{MPSTensor.openSuffixParentHamiltonianES_eq_zero_of_lt}
    \leanok
    \uses{def:open_suffix_parent_hamiltonian_es}
    If $n<L$, then
    \begin{align}
        H^{\mathrm{suffix}}_{N,L;L,n}(A) & =0.
        \label{eq:open_suffix_parent_hamiltonian_zero_before_window}
    \end{align}
\end{theorem}

\begin{proof}
    \leanok
    \uses{def:open_suffix_parent_hamiltonian_es}
    A summand would require a starting site $i$ satisfying both $n-L\leq i$
    and $i+L\leq n$.  The latter inequality is impossible when $n<L$ and
    $i\geq0$, so the defining sum is empty.
\end{proof}

% Source: References/cond-mat_9410110/main.tex, condition C1,
% lines 1030--1041.
\begin{theorem}[Shifted full-range suffix sum]
    \label{thm:open_suffix_parent_hamiltonian_range_shift}
    \lean{MPSTensor.sum_range_openSuffixParentHamiltonianES_succ_eq_Icc}
    \leanok
    \uses{def:open_suffix_parent_hamiltonian_es,
        thm:open_suffix_parent_hamiltonian_zero_before_window}
    If $L>0$, then
    \begin{align}
        \sum_{n=0}^{N-1}H^{\mathrm{suffix}}_{N,L;L,n+1}(A)
         & =\sum_{n=L}^{N}H^{\mathrm{suffix}}_{N,L;L,n}(A).
        \label{eq:open_suffix_parent_hamiltonian_range_shift}
    \end{align}
\end{theorem}

\begin{proof}
    \leanok
    \uses{thm:open_suffix_parent_hamiltonian_zero_before_window}
    The substitution $m=n+1$ identifies the left-hand side with the sum over
    $1\leq m\leq N$.  Every term with $m<L$ vanishes by
    Theorem~\ref{thm:open_suffix_parent_hamiltonian_zero_before_window}, leaving
    precisely the stated interval.
\end{proof}

\begin{theorem}[Finite-overlap condition C1 for the compatible open chain]
    \label{thm:open_parent_hamiltonian_es_c1}
    \lean{MPSTensor.openParentHamiltonianES_C1}
    \leanok
    \uses{def:open_parent_hamiltonian_es,
        def:open_suffix_parent_hamiltonian_es}
    If $L\leq l$, then
    \begin{align}
        0
        \leq \sum_{n=l}^{N}H^{\mathrm{suffix}}_{N,L;l,n}(A)
        \leq (l-L+1)H^{\mathrm{open}}_{N,L}(A).
        \label{eq:open_parent_hamiltonian_es_c1}
    \end{align}
    Thus condition C1 of
    \cite{Nachtergaele1996Spectral} holds with $d_l=l-L+1$.
    This is the source value $pl-r+1$ in the present convention $p=1$ and
    interaction range $r=L$.
\end{theorem}

\begin{proof}
    \leanok
    \uses{lem:transported_es_positive}
    Fix a starting site $i\in I_{L,N}$. The suffix Hamiltonian indexed by $n$
    contains its local term exactly when
    \begin{align}
        i+L\leq n\leq i+l.
    \end{align}
    This interval contains at most $l-L+1$ integers. Reversing the two finite
    sums therefore writes the left-hand side as
    \begin{align}
        \sum_{i\in I_{L,N}}c_i h_{i,\mathrm{ES}}(A,L),
    \end{align}
    where $0\leq c_i\leq l-L+1$. Positivity of each local term turns the
    coefficient bound into the upper operator inequality. Positivity of their
    sum gives the lower inequality.
\end{proof}

% Source: References/cond-mat_9410110/main.tex, condition C1,
% lines 1030--1041; proof of Theorem 2.1(i), lines 1250--1255.
\begin{theorem}[Full-range compatible open-chain counting bound]
    \label{thm:open_parent_hamiltonian_c1_full_range}
    \lean{MPSTensor.openParentHamiltonianES_C1_full_range_quadratic_form}
    \leanok
    \uses{thm:open_parent_hamiltonian_es_c1,
        thm:open_suffix_parent_hamiltonian_range_shift}
    Set $L=l+1$.  For every vector $v$ on the $N$-site chain,
    \begin{align}
        0
         & \leq \sum_{n=0}^{N-1}
        \re\left\langle H^{\mathrm{suffix}}_{N,L;L,n+1}(A)v,v\right\rangle
        \leq
        \re\left\langle H^{\mathrm{open}}_{N,L}(A)v,v\right\rangle.
        \label{eq:open_parent_hamiltonian_c1_full_range}
    \end{align}
    Thus the full martingale range satisfies C1 with $d_{l+1}=1$.
\end{theorem}

\begin{proof}
    \leanok
    \uses{thm:open_parent_hamiltonian_es_c1,
        thm:open_suffix_parent_hamiltonian_range_shift,
        lem:fixed-volume-hamiltonians-positive}
    The shift identity (\ref{eq:open_suffix_parent_hamiltonian_range_shift}) rewrites the
    sum over the interval $L\leq n\leq N$.  Apply
    Theorem~\ref{thm:open_parent_hamiltonian_es_c1} with suffix length $L$.
    Each interaction then occurs once, since $L-L+1=1$.  Positivity of the
    suffix Hamiltonians gives the lower bound.
\end{proof}

% Source anchors: References/cond-mat_9410110/main.tex, condition C2,
% lines 1043--1058; Theorem 2.1(i), lines 1119--1130; and the use of C2 with
% gamma_{l+1} in the printed proof, lines 1240--1248.
\begin{theorem}[Embedded fixed-window condition C2]
    \label{thm:open_suffix_parent_hamiltonian_es_c2}
    \lean{MPSTensor.openSuffixParentHamiltonianES_eq_localTermES,
        MPSTensor.openSuffixParentHamiltonianES_isSymmetricProjection,
        MPSTensor.openSuffixParentHamiltonianES_eq_orthogonal_starProjection_ker,
        MPSTensor.openSuffixParentHamiltonianES_eq_id_sub_starProjection_ker,
        MPSTensor.openSuffixParentHamiltonianES_C2,
        MPSTensor.openSuffixParentHamiltonianES_C2_quadratic_form,
        MPSTensor.openSuffixParentHamiltonianES_norm_eq_of_mem_orthogonal_ker,
        MPSTensor.openSuffixParentHamiltonianES_unit_gap}
    \leanok
    \uses{def:open_suffix_parent_hamiltonian_es,
        thm:open_parent_hamiltonian_es_c1,
        lem:transported_es_local_projections}
    Let $L>0$ and $L\leq m\leq N$. The suffix of length $L$ ending at $m$
    contains one interaction:
    \begin{align}
        H^{\mathrm{suffix}}_{N,L;L,m}(A)
         & =h_{m-L,\mathrm{ES}}(A,L).
        \label{eq:open_suffix_parent_hamiltonian_unique_term}
    \end{align}
    Write $G^{\ker}_{N,L;m}$ for the orthogonal projection onto the kernel of
    this suffix Hamiltonian. Then
    \begin{align}
        H^{\mathrm{suffix}}_{N,L;L,m}(A)
         & =\operatorname{Proj}_{(\ker H^{\mathrm{suffix}}_{N,L;L,m}(A))^\perp}
        =I-G^{\ker}_{N,L;m}.
        \label{eq:open_suffix_parent_hamiltonian_id_sub_kernel}
    \end{align}
    Consequently condition C2 holds in Loewner order with $\gamma_L=1$:
    \begin{align}
        H^{\mathrm{suffix}}_{N,L;L,m}(A)
         & \geq I-G^{\ker}_{N,L;m}.
        \label{eq:open_suffix_parent_hamiltonian_c2}
    \end{align}
    Equivalently, for every $v$,
    \begin{align}
        \re\left\langle (I-G^{\ker}_{N,L;m})v,v\right\rangle
         & \leq
        \re\left\langle H^{\mathrm{suffix}}_{N,L;L,m}(A)v,v\right\rangle.
        \label{eq:open_suffix_parent_hamiltonian_c2_form}
    \end{align}
    If $v\in(\ker H^{\mathrm{suffix}}_{N,L;L,m}(A))^\perp$, then
    \begin{align}
        \left\lVert H^{\mathrm{suffix}}_{N,L;L,m}(A)v\right\rVert
         & =\lVert v\rVert.
        \label{eq:open_suffix_parent_hamiltonian_norm}
    \end{align}
    In particular,
    \begin{align}
        \lVert v\rVert
         & \leq\left\lVert H^{\mathrm{suffix}}_{N,L;L,m}(A)v\right\rVert.
        \label{eq:open_suffix_parent_hamiltonian_unit_gap}
    \end{align}
    In the notation of \cite{Nachtergaele1996Spectral}, set $L=l+1$ and
    $m=n+1$. The interval is then
    $\Lambda_{n+1}\setminus\Lambda_{n-l}=[n-l,n+1)$, so
    \begin{align}
        d_{l+1}
         & =1.
    \end{align}
    Condition C2 has the constant
    \begin{align}
        \gamma_{l+1}
         & =1.
    \end{align}
    The later condition C3 uses the distinct overlap parameter
    $\epsilon_l$.
\end{theorem}

\begin{proof}
    \leanok
    \uses{def:open_suffix_parent_hamiltonian_es,
        thm:open_parent_hamiltonian_es_c1,
        lem:transported_es_local_projections}
    % The interval calculation specializes the compatible open-chain sum to
    % L=l+1. Source lines 1240--1248 then use C2 with gamma_{l+1}.
    An admissible start $i$ satisfies
    \begin{align}
        m-L
         & \leq i.
    \end{align}
    It also satisfies
    \begin{align}
        i+L
         & \leq m.
    \end{align}
    Hence $i=m-L$, which proves
    (\ref{eq:open_suffix_parent_hamiltonian_unique_term}). The sole
    local term is an orthogonal projection. For any orthogonal projection $H$,
    \begin{align}
        \operatorname{ran}H
         & =(\ker H)^\perp.
    \end{align}
    Therefore
    \begin{align}
        H
         & =\operatorname{Proj}_{(\ker H)^\perp}
        =I-\operatorname{Proj}_{\ker H}.
    \end{align}
    This gives (\ref{eq:open_suffix_parent_hamiltonian_id_sub_kernel}).
    The Loewner and quadratic-form inequalities are equalities. On
    $(\ker H)^\perp=\operatorname{ran}H$, idempotence gives $Hv=v$, and hence
    exact norm preservation and the unit lower bound. Condition C1 at suffix
    length $L$ gives $d_L=L-L+1=1$. With $L=l+1$, the two constants are therefore
    $d_{l+1}=1$ and $\gamma_{l+1}=1$.
\end{proof}

% Source: References/cond-mat_9410110/main.tex, condition C2,
% lines 1043--1058; proof of Theorem 2.1(i), lines 1240--1248.
\begin{theorem}[Full-range fixed-window local gap]
    \label{thm:open_suffix_parent_hamiltonian_c2_full_range}
    \lean{MPSTensor.openSuffixParentHamiltonianES_C2_full_range_norm_sq}
    \leanok
    \uses{thm:open_suffix_parent_hamiltonian_zero_before_window,
        thm:open_suffix_parent_hamiltonian_es_c2}
    Set $L=l+1$, let $0\leq n<N$, and let $Q_n$ be the orthogonal projection
    onto the kernel of
    $H^{\mathrm{suffix}}_{N,L;L,n+1}(A)$.  Then every vector $v$ satisfies
    \begin{align}
        \lVert (I-Q_n)v\rVert^2
         & \leq
        \re\left\langle H^{\mathrm{suffix}}_{N,L;L,n+1}(A)v,v\right\rangle.
        \label{eq:open_suffix_parent_hamiltonian_c2_full_range}
    \end{align}
    Hence C2 holds on the full martingale range with $\gamma_{l+1}=1$.
\end{theorem}

\begin{proof}
    \leanok
    \uses{thm:open_suffix_parent_hamiltonian_zero_before_window,
        thm:open_suffix_parent_hamiltonian_es_c2}
    If $n<l$, the suffix Hamiltonian is zero and $Q_n=I$, so both sides of
    (\ref{eq:open_suffix_parent_hamiltonian_c2_full_range}) vanish.  If
    $n\geq l$, the suffix contains the single length-$(l+1)$ interaction and
    equals $I-Q_n$ by Theorem~\ref{thm:open_suffix_parent_hamiltonian_es_c2}.
    Since $I-Q_n$ is an orthogonal projection,
    \begin{align}
        \lVert(I-Q_n)v\rVert^2
         & =\re\langle (I-Q_n)v,v\rangle.
    \end{align}
\end{proof}

\begin{theorem}[Zero energy forces every compatible local constraint]
    \label{thm:open_parent_hamiltonian_zero_local_terms}
    \lean{MPSTensor.localTermES_eq_zero_of_openParentHamiltonianES_eq_zero}
    \leanok
    \uses{def:open_parent_hamiltonian_es}
    If $H^{\mathrm{open}}_{N,L}(A)v=0$, then, for every $i\in I_{L,N}$,
    \begin{align}
        h_{i,\mathrm{ES}}(A,L)v & =0.
    \end{align}
\end{theorem}

\begin{proof}
    \leanok
    \uses{rem:projection_geometry_rowsum_identities, lem:transported_es_local_projections}
    Taking the inner product of the sum with $v$ gives
    \begin{align}
        \sum_{i\in I_{L,N}}\re
        \langle h_{i,\mathrm{ES}}(A,L)v,v\rangle
         & =\re\langle H^{\mathrm{open}}_{N,L}(A)v,v\rangle=0.
    \end{align}
    Each summand is the squared norm of the image of $v$ under an orthogonal
    projection. Hence every local term annihilates $v$.
\end{proof}

\begin{theorem}[Restriction of an open-boundary MPS vector]
    \label{thm:contiguous_restriction_ground_space_map}
    \lean{MPSTensor.contiguousRestrictₗ_groundSpaceMap_mem_groundSpace}
    \leanok
    \uses{def:ground_space_parent, rem:contiguous_window_operators}
    Assume $s+L\leq N$. For every boundary matrix $X$ and complementary
    configuration $\tau$,
    \begin{align}
        \Res^\tau_{s,L}\Gamma_N^A(X) & \in G_L(A).
    \end{align}
\end{theorem}

\begin{proof}
    \leanok
    \uses{lem:ground_space_map_apply, lem:eval_word_append}
    For the fixed words
    \begin{align}
        w_L & =(\tau_0,\ldots,\tau_{s-1}),
    \end{align}
    and
    \begin{align}
        w_R & =(\tau_{s+L},\ldots,\tau_{N-1}),
    \end{align}
    the word of a configuration obtained by inserting the window word $\omega$
    is
    \begin{align}
        A^{w_L\omega w_R} & =A^{w_L}A^\omega A^{w_R}.
    \end{align}
    Cyclicity of the trace gives
    \begin{align}
        \operatorname{tr}\!\left(A^{w_L}A^\omega A^{w_R}X\right)
         & =\operatorname{tr}\!\left(A^\omega A^{w_R}XA^{w_L}\right),
    \end{align}
    and hence
    \begin{align}
        \Res^\tau_{s,L}\Gamma_N^A(X)
         & =\Gamma_L^A\!\left(A^{w_R}XA^{w_L}\right)\in G_L(A).
    \end{align}
\end{proof}

\begin{theorem}[Open-boundary vectors have zero compatible energy]
    \label{thm:open_ground_space_le_hamiltonian_kernel}
    \lean{MPSTensor.groundSpaceES_le_ker_openParentHamiltonianES}
    \leanok
    \uses{def:open_parent_hamiltonian_es, def:ground_space_es}
    For every $A$, $L$, and $N$,
    \begin{align}
        G_N^{\mathrm{ES}}(A)
         & \subseteq \ker H^{\mathrm{open}}_{N,L}(A).
    \end{align}
\end{theorem}

\begin{proof}
    \leanok
    \uses{lem:mem_ground_space_es, thm:contiguous_restriction_ground_space_map, rem:cyclic_window_operators, rem:euclidean_local_projector_ingredients, lem:local_term_es_kernel_restrictions}
    Let $v=\iota_N(\Gamma_N^A(X))\in G_N^{\mathrm{ES}}(A)$. For every
    non-wrapping window $[s,s+L-1]$ and complementary configuration $\tau$,
    Theorem~\ref{thm:contiguous_restriction_ground_space_map} gives
    \begin{align}
        \Res^\tau_{s,L}\Gamma_N^A(X) & \in G_L(A).
    \end{align}
    The non-wrapping identification and
    Lemma~\ref{lem:mem_ground_space_es} then give
    \begin{align}
        R_{s,\tau}v
         & =\iota_L\!\left(\Res^\tau_{s,L}\Gamma_N^A(X)\right)
        \in G_L^{\mathrm{ES}}(A).
    \end{align}
    The local kernel criterion therefore gives
    \begin{align}
        h_{s,\mathrm{ES}}(A,L)v & =0.
    \end{align}
    Thus every compatible local term annihilates $v$.
\end{proof}

\begin{theorem}[Compatible zero energy has open-boundary form]
    \label{thm:open_hamiltonian_kernel_le_ground_space}
    \lean{MPSTensor.ker_openParentHamiltonianES_le_groundSpaceES_of_isNBlkInjective}
    \leanok
    \uses{def:open_parent_hamiltonian_es, def:ground_space_es, def:l_blk_injective}
    Assume $D\geq1$, $L_0>0$, and $N\geq L_0+1$. If $A$ is
    $L_0$-block-injective, then
    \begin{align}
        \ker H^{\mathrm{open}}_{N,L_0+1}(A)
         & \subseteq G_N^{\mathrm{ES}}(A).
    \end{align}
\end{theorem}

\begin{proof}
    \leanok
    \uses{thm:open_parent_hamiltonian_zero_local_terms, lem:mem_ground_space_es, rem:contiguous_window_operators, rem:cyclic_window_operators, lem:local_term_es_kernel_restrictions, thm:normal_open_chain_range_reduction}
    In the following display, $s$ ranges over all starts satisfying
    $s+L_0+1\leq N$, and $\tau$ ranges over all complementary configurations.
    The three ingredients give
    \begin{align}
        H^{\mathrm{open}}_{N,L_0+1}(A)v=0
         & \Longrightarrow
        \bigl(\forall s,\ h_{s,\mathrm{ES}}(A,L_0+1)v=0\bigr)
        \Longrightarrow
        \bigl(\forall s,\tau,\
        \Res^\tau_{s,L_0+1}v\in G_{L_0+1}^{\mathrm{ES}}(A)\bigr)
        \Longrightarrow v\in G_N^{\mathrm{ES}}(A).
    \end{align}
    The middle implication is the local kernel criterion, and the last is the
    open-chain intersection theorem.
\end{proof}

\begin{theorem}[Ground-state property of the compatible open chain]
    \label{thm:open_parent_hamiltonian_kernel_ground_space}
    \lean{MPSTensor.ker_openParentHamiltonianES_eq_groundSpaceES_of_isNBlkInjective}
    \leanok
    \uses{def:open_parent_hamiltonian_es, def:ground_space_es, def:l_blk_injective}
    Assume $D\geq1$, $L_0>0$, and $N\geq L_0+1$. If $A$ is
    $L_0$-block-injective, then
    \begin{align}
        \ker H^{\mathrm{open}}_{N,L_0+1}(A)
         & =G_N^{\mathrm{ES}}(A).
        \label{eq:open_parent_hamiltonian_kernel_ground_space}
    \end{align}
    This is the canonical MPS instance of the ground-state property in
    \cite[Equation~(3.13)]{Nachtergaele1996Spectral}. The right-hand side is
    the open boundary-condition space and may therefore be degenerate.
\end{theorem}

\begin{proof}
    \leanok
    \uses{thm:open_ground_space_le_hamiltonian_kernel, thm:open_hamiltonian_kernel_le_ground_space}
    Combine the mutual inclusions
    Theorem~\ref{thm:open_ground_space_le_hamiltonian_kernel} and
    Theorem~\ref{thm:open_hamiltonian_kernel_le_ground_space}.
\end{proof}

% Source: References/cond-mat_9410110/main.tex, equation En, lines 1060--1073.
\begin{definition}[Block-injective MPS realization of the open-chain martingale difference]
    \label{def:open_chain_martingale_difference_es}
    \lean{MPSTensor.openChainMartingaleDifferenceES}
    \leanok
    \uses{def:open_chain_ground_subspaces_es, thm:open_parent_hamiltonian_kernel_ground_space}
    Assume $D\geq1$, $l>0$, and that $A$ is $l$-block-injective.  Fix
    $K>0$, set $n=K+l$ and $n_l=l+1$, and work in the Hilbert space on
    $\Lambda_{n+1}$.  The physical open-chain martingale difference is
    \begin{align}
        E_n
         & :=G_{\Lambda_n}-G_{\Lambda_{n+1}} \\
         & =P_{\mathcal U_{K+l}(A)}
        -P_{\mathcal G_{K+l+1}^{\mathrm{ES}}(A)}.
        \label{eq:open_chain_martingale_difference_es}
    \end{align}
    The first projector imposes the ground condition on the left $n$ sites
    while acting on all $n+1$ sites.  Block injectivity identifies these MPS
    spaces with the corresponding open-parent Hamiltonian kernels.  The
    condition $K>0$ excludes the undersized volume, where an interaction of
    length $l+1$ has no nonwrapping window and its physical ground projector is
    the identity.
\end{definition}

\begin{lemma}[Whole ground space lies in the C3 tail space]
    \label{lem:whole_ground_space_le_c3_tail_space}
    \lean{MPSTensor.groundSpaceES_le_openChainTailGroundSpaceES}
    \leanok
    \uses{lem:ground_space_in_tail_ground}
    For every $K,l\in\N$, set $n=K+l$.  Then
    \begin{align}
        \mathcal G_{n+1}^{\mathrm{ES}}(A)
         & \subseteq \mathcal V_{K,l+1}(A).
        \label{eq:whole_ground_space_le_c3_tail_space}
    \end{align}
\end{lemma}

\begin{proof}
    \leanok
    \uses{lem:ground_space_in_tail_ground}
    Every fixed-prefix suffix restriction of a whole-chain ground-space vector
    is a ground-space vector on the suffix.  This is exactly membership in
    $\mathcal V_{K,l+1}(A)$.
\end{proof}

\begin{lemma}[Reverse C3 projector composition]
    \label{lem:reverse_c3_projector_composition}
    \lean{MPSTensor.openChainTailGroundProjection_comp_groundSpaceProjection}
    \leanok
    \uses{lem:whole_ground_space_le_c3_tail_space}
    For every $K,l\in\N$, the tail projector fixes the whole-ground projector:
    \begin{align}
        G_{\Lambda_{n+1}\setminus\Lambda_{n-l}}G_{\Lambda_{n+1}}
         & =G_{\Lambda_{n+1}}.
        \label{eq:reverse_c3_projector_composition}
    \end{align}
\end{lemma}

\begin{proof}
    \leanok
    \uses{lem:whole_ground_space_le_c3_tail_space}
    The inner projection has range in
    $\mathcal G_{n+1}^{\mathrm{ES}}(A)$, which lies in the tail ground space by
    Lemma~\ref{lem:whole_ground_space_le_c3_tail_space}.  Orthogonal projection
    onto that tail space therefore leaves every vector in the inner range
    unchanged.
\end{proof}

% Source: References/cond-mat_9410110/main.tex, condition C3, lines 1083--1094;
% identity following C3, lines 1178--1188.
\begin{theorem}[Block-injective MPS instance of the exact C3 identity]
    \label{thm:exact_c3_martingale_difference_identity}
    \lean{MPSTensor.openChainTailGroundProjection_comp_martingaleDifference}
    \leanok
    \uses{def:open_chain_martingale_difference_es, lem:reverse_c3_projector_composition}
    Assume $D\geq1$, $l>0$, and that $A$ is $l$-block-injective.  Fix $K>0$.
    With $n=K+l$ and $n_l=l+1$, the MPS parent chain satisfies Nachtergaele's
    identity following condition C3:
    \begin{align}
        G_{\Lambda_{n+1}\setminus\Lambda_{n-l}}E_n
         & =G_{\Lambda_{n+1}\setminus\Lambda_{n-l}}G_{\Lambda_n}
        -G_{\Lambda_{n+1}}.
        \label{eq:exact_c3_martingale_difference_identity}
    \end{align}
\end{theorem}

\begin{proof}
    \leanok
    \uses{def:open_chain_martingale_difference_es, lem:reverse_c3_projector_composition}
    Distribute the tail projector over
    $E_n=G_{\Lambda_n}-G_{\Lambda_{n+1}}$.  The second product is
    $G_{\Lambda_{n+1}}$ by
    Lemma~\ref{lem:reverse_c3_projector_composition}, which gives the displayed
    identity.
\end{proof}

\begin{theorem}[Nachtergaele open-chain projector defect]
    \label{thm:nachtergaele_open_chain_projector_defect}
    \lean{MPSTensor.re_inner_openChain_anticommutator_ge_neg_of_groundProjection_defect}
    \leanok
    \uses{def:open_chain_ground_subspaces_es, def:ground_space_es, def:l_blk_injective}
    Assume $D\ge1$ and let $A$ be $L_0$-block-injective with $L_0>0$. Set
    $n=K+L_0$ and $l=L_0$, and write
    \begin{align}
        S & =\mathcal U_{K+L_0}(A),          &
        T & =\mathcal V_{K,L_0+1}(A). \notag
    \end{align}
    Thus $S$ and $T$ are the left and tail ground spaces in the Hilbert space
    on $\Lambda_{n+1}$. Let $q_S=P_{S^\perp}$ and $q_T=P_{T^\perp}$ be the
    complementary excitation projections. If
    \begin{align}
        \|G_{\Lambda_{n+1}\setminus\Lambda_{n-l}}G_{\Lambda_n}
        -G_{\Lambda_{n+1}}\| & \le\eta, \notag
    \end{align}
    then every vector $v$ in the Hilbert space on $\Lambda_{n+1}$ satisfies
    \begin{align}
        -\eta\bigl(
        \re\langle q_Tv,v\rangle+\re\langle q_Sv,v\rangle\bigr)
         & \le
        \re\langle q_Tv,q_Sv\rangle+\re\langle q_Sv,q_Tv\rangle. \notag
    \end{align}
\end{theorem}

\begin{proof}
    \leanok
    \uses{thm:open_chain_left_tail_ground_intersection, thm:projector_defect_to_friedrichs, thm:friedrichs_ground_space_excitation_anticommutator}
    The open-chain intersection theorem identifies $S\cap T$ with the ground
    space on $\Lambda_{n+1}$. The projector-defect adapter gives a Friedrichs
    overlap bound, and the reduced two-projection theorem gives the stated
    anticommutator inequality.
\end{proof}

\begin{theorem}[Uniform tail virtual-map bound]
    \label{thm:primitive_tail_virtual_map_norm_uniform}
    \lean{MPSTensor.IsPrimitiveMPS.tailVirtualMapES_norm_four_pow_uniform}
    \lean{MPSTensor.IsPrimitiveMPS.tailVirtualMapES_norm_uniform}
    \leanok
    \uses{def:primitive_mps, thm:spectator_boundary_hilbert_factorization}
    Let $A$ be a primitive MPS tensor with nonzero virtual dimension, and let
    $J_K=J^{\mathrm{tail}}_K$ be its length-$K$ tail virtual map.  Then there is
    a constant $C>0$ such that, for every $K\in\N$,
    \begin{align}
        \lVert J_K\rVert & \leq C.
    \end{align}
    No positive-definiteness hypothesis on the fixed point is required.
\end{theorem}

\begin{proof}
    \leanok
    \uses{thm:geometric_bound_of_spectral_radius_lt_one, thm:sum_entry_norm_sq_le_card_mul_l2_operator_norm_sq}
    The operator norm is controlled using the adjoint identity
    \begin{align}
        J_K^\dagger J_K U(X)
         & =U\!\left(X\,\mathcal E_A^K(\mathbf 1)\right).
    \end{align}
    The Hilbert--Schmidt product estimate then gives the fourth-power bound
    \begin{align}
        \lVert J_K\rVert^4
         & \leq \sum_{i,j}\left\lvert
        \mathcal E_A^K(\mathbf 1)_{ij}\right\rvert^2.
    \end{align}
    For $K\geq1$, decompose the transfer power into its fixed-point projection
    and the $K$-th power of the complementary map.  The complementary powers
    are geometrically bounded by
    Theorem~\ref{thm:geometric_bound_of_spectral_radius_lt_one}, while the
    entrywise mass is controlled by the $L^2$ operator norm with the virtual
    dimension factor from
    Theorem~\ref{thm:sum_entry_norm_sq_le_card_mul_l2_operator_norm_sq}.
    The case $K=0$ is bounded directly.  This first gives a uniform bound for
    $\lVert J_K\rVert^4$ and hence a strictly positive uniform bound for
    $\lVert J_K\rVert$.  No FNW or Nachtergaele contraction estimate enters.
\end{proof}

\begin{theorem}[Left-canonical normalization of fixed-length words]
    \label{thm:word_left_canonical_normalization}
    \lean{MPSTensor.sum_evalWord_conjTranspose_mul_evalWord}
    \leanok
    \uses{def:eval_word}
    If $A$ is left-canonical, then for every $Q\in\N$,
    \begin{align}
        \sum_{j\in\Omega_Q}(A^j)^\dagger A^j & =\mathbf 1,
    \end{align}
    where $\Omega_Q=\Fin(d^Q)$ is the set of physical words of length $Q$.
\end{theorem}

\begin{proof}
    \leanok
    \uses{def:eval_word}
    Induct on $Q$.  Split a word of length $Q+1$ into its first letter and
    its remaining suffix, apply the induction hypothesis to the suffix sum,
    and then use left-canonicity for the outer one-site sum.
\end{proof}

\begin{theorem}[Whole-increment left virtual norm]
    \label{thm:whole_increment_left_virtual_norm}
    \lean{MPSTensor.leftVirtualMapES_norm_map_of_leftCanonical}
    \lean{MPSTensor.leftVirtualMapES_norm_le_one_of_leftCanonical}
    \leanok
    \uses{thm:word_left_canonical_normalization}
    If $A$ is left-canonical, then the length-$Q$ left virtual map
    $J_Q^{\mathrm{left}}$ preserves the Hilbert--Schmidt norm.  In particular,
    \begin{align}
        \lVert J_Q^{\mathrm{left}}X\rVert & =\lVert X\rVert, \\
        \lVert J_Q^{\mathrm{left}}\rVert  & \leq1.
    \end{align}
\end{theorem}

\begin{proof}
    \leanok
    \uses{thm:word_left_canonical_normalization}
    Expand the squared norm as a sum over $j\in\Omega_Q$.  The identity from
    Theorem~\ref{thm:word_left_canonical_normalization} reduces the sum to the
    Hilbert--Schmidt norm of $X$.  The operator-norm bound follows immediately.
\end{proof}

\begin{definition}[Physical reassociation]
    \label{def:whole_increment_physical_reassociation}
    \lean{MPSTensor.physicalReassocES}
    \leanok
    For $K,L,Q\in\N$, let
    \begin{align}
        \mathscr A_{K,L,Q}\colon
        \ell^2(\Omega_{K+(L+Q)}) & \longrightarrow
        \ell^2(\Omega_{(K+L)+Q})
    \end{align}
    be the unitary coordinate map induced by associativity of addition.
\end{definition}

\begin{theorem}[Reassociation of the ground-space boundary map]
    \label{thm:whole_increment_ground_space_reassociation}
    \lean{MPSTensor.physicalReassocES_comp_groundSpaceMapES}
    \leanok
    \uses{def:whole_increment_physical_reassociation, def:ground_space_map_es}
    For every $K,L,Q\in\N$,
    \begin{align}
        \mathscr A_{K,L,Q}\Gamma_{K+(L+Q)}^{\mathrm{ES}}
         & =\Gamma_{K+L+Q}^{\mathrm{ES}}.
    \end{align}
\end{theorem}

\begin{proof}
    \leanok
    \uses{def:whole_increment_physical_reassociation, def:ground_space_map_es}
    Evaluate both sides at a physical word.  Reassociation changes only the
    parenthesization of its coordinate set, so both coefficients are the trace
    pairing with the same word product.
\end{proof}

\begin{definition}[Reassociated tail boundary map]
    \label{def:reassociated_tail_boundary_map}
    \lean{MPSTensor.reassocTailBoundaryMapES}
    \leanok
    \uses{def:whole_increment_physical_reassociation}
    Define the tail boundary map in the common $(K+L+Q)$-site coordinates by
    \begin{align}
        \widetilde\Gamma^{\mathrm{tail}}_{K;L,Q}
         & :=\mathscr A_{K,L,Q}\Gamma^{\mathrm{tail}}_{K,L+Q}.
    \end{align}
\end{definition}

\begin{theorem}[One-site reassociation]
    \label{thm:reassociated_tail_boundary_one_site}
    \lean{MPSTensor.reassocTailBoundaryMapES_one}
    \leanok
    \uses{def:reassociated_tail_boundary_map, def:spectator_boundary_maps_es}
    A one-site suffix increment needs no reassociation:
    \begin{align}
        \widetilde\Gamma^{\mathrm{tail}}_{K;L,1}
         & =\Gamma^{\mathrm{tail}}_{K,L+1}.
        \label{eq:reassociated_tail_boundary_one_site}
    \end{align}
\end{theorem}

\begin{proof}
    \leanok
    \uses{def:reassociated_tail_boundary_map,
        def:whole_increment_physical_reassociation}
    Reassociating a window whose last block carries a single site is the
    identity, so the reassociated tail map is the ordinary tail map with $L+1$
    spectator sites.
\end{proof}

\begin{theorem}[Reassociated tail boundary map on three intervals]
    \label{thm:reassociated_tail_boundary_three_block}
    \lean{MPSTensor.reassocTailBoundaryMapES_apply_threeBlock}
    \leanok
    \uses{def:reassociated_tail_boundary_map, def:ground_space_map}
    Let $K,L,Q\in\N$ and let $Y\colon(\Fin K\to\Fin d)\to\MN{D}$ be a
    boundary family. For all $u\colon\Fin K\to\Fin d$,
    $w\colon\Fin L\to\Fin d$, and $v\colon\Fin Q\to\Fin d$,
    \begin{align}
        \widetilde\Gamma^{\mathrm{tail}}_{K;L,Q}(Y)\bigl((u,w),v\bigr)
         & =\Gamma_{L+Q}(Y_u)(w,v).
    \end{align}
\end{theorem}

\begin{proof}
    \leanok
    \uses{def:reassociated_tail_boundary_map,
        def:whole_increment_physical_reassociation,
        thm:tail_boundary_map_append}
    The reassociation map only regroups the physical coordinates:
    for every vector $f$ on $K+(L+Q)$ sites,
    \begin{align}
        (\mathscr A_{K,L,Q}f)\bigl((u,w),v\bigr)
         & =f\bigl(u,(w,v)\bigr).
    \end{align}
    Applying this to $f=\Gamma^{\mathrm{tail}}_{K,L+Q}(Y)$ and then
    Theorem~\ref{thm:tail_boundary_map_append} with suffix $(w,v)$ of
    length $L+Q$ gives
    \begin{align}
        \widetilde\Gamma^{\mathrm{tail}}_{K;L,Q}(Y)\bigl((u,w),v\bigr)
         & =\Gamma^{\mathrm{tail}}_{K,L+Q}(Y)\bigl(u,(w,v)\bigr)
        =\Gamma_{L+Q}(Y_u)(w,v).
    \end{align}
\end{proof}

\begin{theorem}[Factorization of the reassociated tail map]
    \label{thm:reassociated_tail_boundary_factorization}
    \lean{MPSTensor.reassocTailBoundaryMapES_comp_tailVirtualMapES}
    \leanok
    \uses{def:reassociated_tail_boundary_map, thm:spectator_boundary_hilbert_factorization, thm:whole_increment_ground_space_reassociation}
    The reassociated tail map satisfies
    \begin{align}
        \widetilde\Gamma^{\mathrm{tail}}_{K;L,Q}
        J_K^{\mathrm{tail}}
         & =\Gamma_{K+L+Q}^{\mathrm{ES}}.
    \end{align}
\end{theorem}

\begin{proof}
    \leanok
    \uses{def:reassociated_tail_boundary_map, thm:spectator_boundary_hilbert_factorization, thm:whole_increment_ground_space_reassociation}
    Substitute the definition of the reassociated tail map, use the ordinary
    tail factorization, and then reassociate the full boundary map.
\end{proof}

\begin{theorem}[Whole-increment inverse-Gram projector residual]
    \label{thm:whole_increment_projector_residual}
    \lean{MPSTensor.whole_increment_injectiveRangeProjector_residual_eq}
    \leanok
    \uses{def:reassociated_tail_boundary_map, thm:reassociated_tail_boundary_factorization, thm:spectator_boundary_hilbert_factorization}
    Let $T=\widetilde\Gamma^{\mathrm{tail}}_{K;L,Q}$,
    $S=\Gamma^{\mathrm{left}}_{K+L,Q}$, and
    $C=\Gamma_{K+L+Q}^{\mathrm{ES}}$.  Assume separately that $T$, $S$, and
    $C$ are injective.  Since
    $TJ_K^{\mathrm{tail}}=C=SJ_Q^{\mathrm{left}}$, the difference
    $P_TP_S-P_C$ equals the exact inverse-Gram residual obtained from these two
    factorizations.
\end{theorem}

\begin{proof}
    \leanok
    \uses{thm:reassociated_tail_boundary_factorization, thm:spectator_boundary_hilbert_factorization}
    Apply the common-factorization range-projector identity to $T$, $S$, and
    $C$, using the two displayed factorizations.  No canonicality or norm
    estimate is used.
\end{proof}

\begin{theorem}[Whole-increment mixed Gram after reassociation]
    \label{thm:whole_increment_mixed_gram}
    \lean{MPSTensor.inner_reassocTailBoundaryMapES_adjoint_leftBoundaryMapES}
    \leanok
    \uses{def:reassociated_tail_boundary_map, thm:spectator_boundary_mixed_gram}
    With the notation of
    Theorem~\ref{thm:spectator_boundary_mixed_gram},
    \begin{align}
        \left\langle X,
        (\widetilde\Gamma^{\mathrm{tail}}_{K;L,Q})^\dagger
        \Gamma^{\mathrm{left}}_{K+L,Q}Z\right\rangle
         & =\sum_{u\in\Omega_K}\sum_{j\in\Omega_Q}
        \left\langle U(A^jX_u),\mathcal K_LU(Z_jA^u)\right\rangle.
    \end{align}
\end{theorem}

\begin{proof}
    \leanok
    \uses{def:reassociated_tail_boundary_map, def:whole_increment_physical_reassociation, thm:spectator_boundary_mixed_gram}
    Expand the reassociated tail map.  Its adjoint introduces
    $\mathscr A_{K,L,Q}^{-1}$, so the coordinate formula of
    Theorem~\ref{thm:spectator_boundary_mixed_gram} applies.
\end{proof}

\begin{theorem}[Centered whole-increment mixed Gram]
    \label{thm:whole_increment_mixed_gram_centered}
    \lean{MPSTensor.inner_reassocTailBoundaryMapES_adjoint_leftBoundaryMapES_centered}
    \leanok
    \uses{thm:whole_increment_mixed_gram, thm:fixed_point_gram_metric_exact}
    Let $\rho\in\MN{D}$ have nonzero trace and set
    $\mathcal K_\infty=\mathscr R(P_\rho)$.  Then
    \begin{align}
         & \left\langle X,
        (\widetilde\Gamma^{\mathrm{tail}}_{K;L,Q})^\dagger
        \Gamma^{\mathrm{left}}_{K+L,Q}Z\right\rangle
        -\sum_{u,j}\left\langle U(A^jX_u),
        \mathcal K_\infty U(Z_jA^u)\right\rangle \\
         & \qquad=
        \sum_{u,j}\left\langle U(A^jX_u),
        (\mathcal K_L-\mathcal K_\infty)U(Z_jA^u)\right\rangle.
    \end{align}
\end{theorem}

\begin{proof}
    \leanok
    \uses{thm:whole_increment_mixed_gram, thm:fixed_point_gram_metric_exact}
    Substitute the reassociated mixed-Gram formula and the action of the
    limiting metric, then distribute subtraction over the finite sums.
\end{proof}

\begin{definition}[Whole-increment left overlap factor]
    \label{def:whole_increment_left_overlap_factor}
    \lean{MPSTensor.wholeIncrementLeftOverlapFactorES}
    \leanok
    For $X=(X_u)_{u\in\Omega_K}$, define
    \begin{align}
        (F_{K,Q}X)_{u,j} & :=A^jX_u.
    \end{align}
\end{definition}

\begin{definition}[Whole-increment right overlap factor]
    \label{def:whole_increment_right_overlap_factor}
    \lean{MPSTensor.wholeIncrementRightOverlapFactorES}
    \leanok
    For $Z=(Z_j)_{j\in\Omega_Q}$, define
    \begin{align}
        (G_{K,Q}Z)_{u,j} & :=Z_jA^u.
    \end{align}
\end{definition}

\begin{theorem}[Fiber formula for the whole-increment left factor]
    \label{thm:whole_increment_left_overlap_factor_apply}
    \lean{MPSTensor.boundaryFamilyEquiv_wholeIncrementLeftOverlapFactorES_apply}
    \leanok
    \uses{def:whole_increment_left_overlap_factor}
    The matrix fiber of $F_{K,Q}X$ at $(u,j)$ is $A^jX_u$.
\end{theorem}

\begin{proof}
    \leanok
    \uses{def:whole_increment_left_overlap_factor}
    Unfold the direct-sum factor and the Frobenius coordinate identification.
\end{proof}

\begin{theorem}[Fiber formula for the whole-increment right factor]
    \label{thm:whole_increment_right_overlap_factor_apply}
    \lean{MPSTensor.boundaryFamilyEquiv_wholeIncrementRightOverlapFactorES_apply}
    \leanok
    \uses{def:whole_increment_right_overlap_factor}
    The matrix fiber of $G_{K,Q}Z$ at $(u,j)$ is $Z_jA^u$.
\end{theorem}

\begin{proof}
    \leanok
    \uses{def:whole_increment_right_overlap_factor}
    Unfold the direct-sum factor and the Frobenius coordinate identification.
\end{proof}

\begin{theorem}[Norm bounds for the whole-increment overlap factors]
    \label{thm:whole_increment_overlap_factor_norms}
    \lean{MPSTensor.wholeIncrementLeftOverlapFactorES_norm_le}
    \lean{MPSTensor.wholeIncrementRightOverlapFactorES_norm_le}
    \leanok
    \uses{def:whole_increment_left_overlap_factor, def:whole_increment_right_overlap_factor}
    The overlap factors satisfy
    \begin{align}
        \lVert F_{K,Q}\rVert & \leq\lVert J_Q^{\mathrm{left}}\rVert, \\
        \lVert G_{K,Q}\rVert & \leq\lVert J_K^{\mathrm{tail}}\rVert.
    \end{align}
\end{theorem}

\begin{proof}
    \leanok
    \uses{def:whole_increment_left_overlap_factor, def:whole_increment_right_overlap_factor}
    Expand the squared direct-sum norms fiberwise and apply the operator-norm
    bound to each fiber.  Summing gives the two inequalities without a factor
    depending on the number of spectator words.
\end{proof}

\begin{definition}[Whole-increment centered residual]
    \label{def:whole_increment_centered_residual}
    \lean{MPSTensor.wholeIncrementCenteredResidualES}
    \leanok
    \uses{def:reassociated_tail_boundary_map, def:whole_increment_left_overlap_factor, def:whole_increment_right_overlap_factor}
    Assume $\operatorname{tr}(\rho)\neq 0$ and let
    $\mathcal K_\infty=\mathscr R(P_\rho)$.  Define
    \begin{align}
        R_{K,L,Q}
         & :=(\widetilde\Gamma^{\mathrm{tail}}_{K;L,Q})^\dagger
        \Gamma^{\mathrm{left}}_{K+L,Q}
        -F_{K,Q}^\dagger
        \mathcal K_\infty^{\oplus(\Omega_K\times\Omega_Q)}G_{K,Q}.
    \end{align}
\end{definition}

\begin{theorem}[Whole-increment centered factorization]
    \label{thm:whole_increment_centered_factorization}
    \lean{MPSTensor.wholeIncrementCenteredResidualES_eq_overlapFactors}
    \leanok
    \uses{def:whole_increment_centered_residual, thm:whole_increment_left_overlap_factor_apply, thm:whole_increment_right_overlap_factor_apply, thm:whole_increment_mixed_gram_centered}
    The centered residual factors as
    \begin{align}
        R_{K,L,Q}
         & =F_{K,Q}^\dagger
        (\mathcal K_L-\mathcal K_\infty)^{\oplus
            (\Omega_K\times\Omega_Q)}G_{K,Q}.
    \end{align}
\end{theorem}

\begin{proof}
    \leanok
    \uses{def:whole_increment_centered_residual, thm:whole_increment_left_overlap_factor_apply, thm:whole_increment_right_overlap_factor_apply, thm:whole_increment_mixed_gram_centered}
    Test both sides against arbitrary boundary families, use the two fiber
    formulas, and apply the centered mixed-Gram identity.
\end{proof}

\begin{theorem}[Whole-increment centered residual bound]
    \label{thm:whole_increment_centered_residual_bound}
    \lean{MPSTensor.wholeIncrementCenteredResidualES_norm_le}
    \leanok
    \uses{thm:whole_increment_centered_factorization, thm:whole_increment_overlap_factor_norms, thm:boundary_fiberwise_map_norm}
    The factorization gives
    \begin{align}
        \lVert R_{K,L,Q}\rVert
         & \leq\lVert J_Q^{\mathrm{left}}\rVert
        \lVert\mathcal K_L-\mathcal K_\infty\rVert
        \lVert J_K^{\mathrm{tail}}\rVert.
    \end{align}
\end{theorem}

\begin{proof}
    \leanok
    \uses{thm:whole_increment_centered_factorization, thm:whole_increment_overlap_factor_norms, thm:boundary_fiberwise_map_norm}
    Apply submultiplicativity to the three factors and use the operator norm of
    a fiberwise direct sum.
\end{proof}

\begin{theorem}[Left-canonical whole-increment centered residual bound]
    \label{thm:whole_increment_centered_residual_bound_left_canonical}
    \lean{MPSTensor.wholeIncrementCenteredResidualES_norm_le_of_leftCanonical}
    \leanok
    \uses{thm:whole_increment_centered_residual_bound, thm:whole_increment_left_virtual_norm}
    If $A$ is left-canonical, then
    \begin{align}
        \lVert R_{K,L,Q}\rVert
         & \leq\lVert\mathcal K_L-\mathcal K_\infty\rVert
        \lVert J_K^{\mathrm{tail}}\rVert.
    \end{align}
\end{theorem}

\begin{proof}
    \leanok
    \uses{thm:whole_increment_centered_residual_bound, thm:whole_increment_left_virtual_norm}
    Substitute $\lVert J_Q^{\mathrm{left}}\rVert\leq1$ from
    Theorem~\ref{thm:whole_increment_left_virtual_norm} into
    Theorem~\ref{thm:whole_increment_centered_residual_bound}.
\end{proof}

\begin{theorem}[Primitive whole-increment centered residual bound]
    \label{thm:primitive_whole_increment_centered_residual_bound}
    \lean{MPSTensor.IsPrimitiveMPS.wholeIncrementCenteredResidualES_norm_le}
    \leanok
    \uses{def:primitive_mps, thm:whole_increment_centered_residual_bound_left_canonical}
    Let $A$ be a primitive MPS tensor with positive-definite fixed point
    $\rho$.  Then $A$ satisfies the left-canonical centered residual estimate
    of Theorem~\ref{thm:whole_increment_centered_residual_bound_left_canonical}.
\end{theorem}

\begin{proof}
    \leanok
    \uses{def:primitive_mps, thm:whole_increment_centered_residual_bound_left_canonical}
    Use the left-canonical normalization contained in the primitive-MPS
    hypotheses together with positive definiteness of $\rho$.
\end{proof}

\begin{theorem}[Limiting overlap product]
    \label{thm:whole_increment_limiting_overlap_product}
    \lean{MPSTensor.wholeIncrement_limitingOverlapProduct_eq}
    \leanok
    \uses{def:whole_increment_left_overlap_factor, def:whole_increment_right_overlap_factor, thm:whole_increment_left_overlap_factor_apply, thm:whole_increment_right_overlap_factor_apply, thm:limiting_metric_inverse_virtual_adjoints}
    Assume $D>0$.  For a positive-definite matrix $\rho$, the product
    $F_{K,Q}^\dagger\mathcal K_\infty^{\oplus}G_{K,Q}$ equals the corresponding
    virtual-word expression formed from $J_K^{\mathrm{tail}}$,
    $(J_Q^{\mathrm{left}})^\dagger$, $\mathcal K_\infty$, and
    $\mathcal K_\infty^{-1}$.
\end{theorem}

\begin{proof}
    \leanok
    \uses{thm:whole_increment_left_overlap_factor_apply, thm:whole_increment_right_overlap_factor_apply, thm:limiting_metric_inverse_virtual_adjoints}
    Compare matrix fibers, use the explicit limiting-metric action and its
    inverse, and evaluate the two virtual adjoints.
\end{proof}

\begin{theorem}[Gram of the reassociated tail map]
    \label{thm:reassociated_tail_boundary_gram}
    \lean{MPSTensor.reassocTailBoundaryMapES_adjoint_comp_self_eq_fiberwise_groundSpaceGram}
    \leanok
    \uses{def:reassociated_tail_boundary_map, thm:spectator_boundary_direct_sum_gram}
    Reassociation preserves the tail Gram:
    \begin{align}
        (\widetilde\Gamma^{\mathrm{tail}}_{K;L,Q})^\dagger
        \widetilde\Gamma^{\mathrm{tail}}_{K;L,Q}
         & =\mathcal K_{L+Q}^{\oplus\Omega_K}.
    \end{align}
\end{theorem}

\begin{proof}
    \leanok
    \uses{def:reassociated_tail_boundary_map, thm:spectator_boundary_direct_sum_gram}
    The reassociation is unitary, so its adjoint cancels it.  The ordinary tail
    Gram is the fiberwise direct sum of $\mathcal K_{L+Q}$.
\end{proof}

\begin{theorem}[Inverse Gram of the reassociated tail map]
    \label{thm:reassociated_tail_inverse_gram}
    \lean{MPSTensor.inverseGram_reassocTailBoundaryMapES_eq_fiberwise_inverseGram}
    \leanok
    \uses{thm:reassociated_tail_boundary_gram, thm:inverse_gram_identities}
    If $\Gamma_{L+Q}^{\mathrm{ES}}$ is injective, then the inverse Gram of the
    reassociated tail map is the fiberwise direct sum of
    $(\mathcal K_{L+Q})^{-1}$.
\end{theorem}

\begin{proof}
    \leanok
    \uses{thm:reassociated_tail_boundary_gram, thm:inverse_gram_identities}
    Multiply the proposed fiberwise inverse on both sides by the Gram from
    Theorem~\ref{thm:reassociated_tail_boundary_gram} and use the inverse-Gram
    identities.
\end{proof}

\begin{theorem}[Whole-increment finite-Gram cancellations]
    \label{thm:whole_increment_finite_gram_cancellations}
    \lean{MPSTensor.wholeIncrementLeftFiniteGramCancellation_eq_groundSpaceMapES_adjoint}
    \lean{MPSTensor.wholeIncrementTailFiniteGramCancellation_eq_groundSpaceMapES_adjoint}
    \leanok
    \uses{thm:left_finite_gram_cancellation, thm:reassociated_tail_boundary_factorization, thm:inverse_gram_identities}
    If $\Gamma_{K+L}^{\mathrm{ES}}$ and
    $\Gamma_{L+Q}^{\mathrm{ES}}$ are injective, respectively, then
    \begin{align}
        (J_Q^{\mathrm{left}})^\dagger
        \mathcal K_{K+L}^{\oplus\Omega_Q}
        (\mathcal K_{K+L}^{-1})^{\oplus\Omega_Q}
        (\Gamma_{K+L,Q}^{\mathrm{left}})^\dagger
         & =(\Gamma_{K+L+Q}^{\mathrm{ES}})^\dagger, \\
        (J_K^{\mathrm{tail}})^\dagger
        \mathcal K_{L+Q}^{\oplus\Omega_K}
        (\mathcal K_{L+Q}^{-1})^{\oplus\Omega_K}
        (\widetilde\Gamma_{K;L,Q}^{\mathrm{tail}})^\dagger
         & =(\Gamma_{K+L+Q}^{\mathrm{ES}})^\dagger.
    \end{align}
\end{theorem}

\begin{proof}
    \leanok
    \uses{thm:left_finite_gram_cancellation, thm:reassociated_tail_boundary_factorization, thm:inverse_gram_identities}
    The left identity is the finite-Gram cancellation of
    Theorem~\ref{thm:left_finite_gram_cancellation} at length $K+L$.  For the
    tail identity, cancel the fiberwise Gram with its inverse and take the
    adjoint of the reassociated tail factorization.
\end{proof}

\begin{definition}[Whole-increment tail finite-Gram correction]
    \label{def:whole_increment_tail_gram_correction}
    \lean{MPSTensor.wholeIncrementTailFiniteGramCorrectionES}
    \leanok
    Assume $D>0$, let $\rho$ be positive definite, and assume that
    $\Gamma_{L+Q}^{\mathrm{ES}}$ and $\Gamma_{K+L+Q}^{\mathrm{ES}}$ are
    injective.  The tail correction is the physical sandwich containing
    $\mathcal K_{L+Q}-\mathcal K_\infty$ between the reassociated tail
    pseudoinverse, $J_K^{\mathrm{tail}}$, and the full boundary pseudoinverse.
\end{definition}

\begin{definition}[Whole-increment full inverse-Gram correction]
    \label{def:whole_increment_full_inverse_gram_correction}
    \lean{MPSTensor.wholeIncrementFullInverseGramCorrectionES}
    \leanok
    Assume $D>0$, let $\rho$ be positive definite, and assume that
    $\Gamma_{L+Q}^{\mathrm{ES}}$ and $\Gamma_{K+L+Q}^{\mathrm{ES}}$ are
    injective.  The full correction is the physical sandwich containing
    $(\mathcal K_{K+L+Q})^{-1}-\mathcal K_\infty^{-1}$.
\end{definition}

\begin{definition}[Whole-increment left finite-Gram correction]
    \label{def:whole_increment_left_gram_correction}
    \lean{MPSTensor.wholeIncrementLeftFiniteGramCorrectionES}
    \leanok
    Assume $D>0$, let $\rho$ be positive definite, and assume that
    $\Gamma_{L+Q}^{\mathrm{ES}}$ and $\Gamma_{K+L}^{\mathrm{ES}}$ are
    injective.  The left correction is the physical sandwich containing
    $\mathcal K_{K+L}-\mathcal K_\infty$ and the left boundary
    pseudoinverse.
\end{definition}

\begin{theorem}[Reassociated tail pseudoinverse bound]
    \label{thm:reassociated_tail_pseudoinverse_bound}
    \lean{MPSTensor.norm_reassocTailBoundaryMapES_comp_inverseGram_le_sqrt}
    \leanok
    \uses{thm:reassociated_tail_inverse_gram, thm:boundary_fiberwise_map_norm}
    If $\Gamma_{L+Q}^{\mathrm{ES}}$ is injective, then
    \begin{align}
        \left\lVert\widetilde\Gamma^{\mathrm{tail}}_{K;L,Q}
        \bigl((\widetilde\Gamma^{\mathrm{tail}}_{K;L,Q})^\dagger
        \widetilde\Gamma^{\mathrm{tail}}_{K;L,Q}\bigr)^{-1}\right\rVert
         & \leq\sqrt{\lVert\mathcal K_{L+Q}^{-1}\rVert}.
    \end{align}
\end{theorem}

\begin{proof}
    \leanok
    \uses{thm:reassociated_tail_inverse_gram, thm:boundary_fiberwise_map_norm}
    Use the square-root pseudoinverse identity, substitute the fiberwise inverse
    Gram, and bound the norm of the fiberwise map.
\end{proof}

\begin{theorem}[Whole-increment virtual correction telescope]
    \label{thm:whole_increment_virtual_correction_telescope}
    \lean{MPSTensor.wholeIncrement_virtualResidual_eq_centered_sub_gramErrors}
    \leanok
    \uses{def:whole_increment_centered_residual, thm:whole_increment_limiting_overlap_product}
    Assume $D>0$, let $\rho$ be positive definite, and assume that
    $\Gamma_{K+L+Q}^{\mathrm{ES}}$ is injective.  Write
    $K_t=\mathcal K_{L+Q}^{\oplus\Omega_K}$,
    $K_t^\infty=\mathcal K_\infty^{\oplus\Omega_K}$,
    $K_\ell=\mathcal K_{K+L}^{\oplus\Omega_Q}$,
    $K_\ell^\infty=\mathcal K_\infty^{\oplus\Omega_Q}$,
    $I_f=\mathcal K_{K+L+Q}^{-1}$,
    $I_\infty=\mathcal K_\infty^{-1}$,
    $J_t=J_K^{\mathrm{tail}}$, and
    $J_\ell^\dagger=(J_Q^{\mathrm{left}})^\dagger$.  Then the exact virtual
    residual satisfies
    \begin{align}
         & (\widetilde\Gamma_{K;L,Q}^{\mathrm{tail}})^\dagger
        \Gamma_{K+L,Q}^{\mathrm{left}}
        -K_tJ_tI_fJ_\ell^\dagger K_\ell                       \\
         & \quad=R_{K,L,Q}
        -(K_t-K_t^\infty)J_tI_fJ_\ell^\dagger K_\ell
        -K_t^\infty J_t(I_f-I_\infty)J_\ell^\dagger K_\ell    \\
         & \qquad\phantom{=R_{K,L,Q}}
        -K_t^\infty J_tI_\infty J_\ell^\dagger
        (K_\ell-K_\ell^\infty).
    \end{align}
\end{theorem}

\begin{proof}
    \leanok
    \uses{def:whole_increment_centered_residual, thm:whole_increment_limiting_overlap_product}
    Substitute the centered residual and the limiting overlap product, distribute
    composition over subtraction, and telescope the product algebraically.
\end{proof}

\begin{definition}[Whole-increment centered projector sandwich]
    \label{def:whole_increment_centered_projector_sandwich}
    \lean{MPSTensor.wholeIncrementCenteredProjectorResidualES}
    \leanok
    \uses{def:whole_increment_centered_residual}
    Assume $D>0$, let $\rho$ be positive definite, and assume that
    $\Gamma_{L+Q}^{\mathrm{ES}}$ and $\Gamma_{K+L}^{\mathrm{ES}}$ are
    injective.  The centered physical sandwich is obtained by placing
    $R_{K,L,Q}$ between the reassociated tail pseudoinverse and the left
    boundary pseudoinverse.
\end{definition}

\begin{theorem}[Whole-increment physical correction telescope]
    \label{thm:whole_increment_finite_gram_corrections}
    \lean{MPSTensor.wholeIncrement_injectiveRangeProjector_residual_eq_centered_sub_corrections}
    \leanok
    \uses{thm:whole_increment_projector_residual, thm:reassociated_tail_inverse_gram, thm:reassociated_tail_boundary_gram, thm:spectator_boundary_direct_sum_gram, thm:whole_increment_virtual_correction_telescope, thm:whole_increment_finite_gram_cancellations, def:whole_increment_tail_gram_correction, def:whole_increment_full_inverse_gram_correction, def:whole_increment_left_gram_correction, def:whole_increment_centered_projector_sandwich}
    Assume $D>0$, let $\rho$ be positive definite, and suppose that
    $\Gamma_{L+Q}^{\mathrm{ES}}$, $\Gamma_{K+L}^{\mathrm{ES}}$, and
    $\Gamma_{K+L+Q}^{\mathrm{ES}}$ are injective.  Then the whole-increment
    projector defect is the centered physical sandwich minus the tail Gram
    correction, the full inverse-Gram correction, and the left Gram correction.
\end{theorem}

\begin{proof}
    \leanok
    \uses{thm:whole_increment_projector_residual, thm:reassociated_tail_inverse_gram, thm:reassociated_tail_boundary_gram, thm:spectator_boundary_direct_sum_gram, thm:whole_increment_virtual_correction_telescope, thm:whole_increment_finite_gram_cancellations, def:whole_increment_tail_gram_correction, def:whole_increment_full_inverse_gram_correction, def:whole_increment_left_gram_correction, def:whole_increment_centered_projector_sandwich}
    Substitute the common-factorization residual, the two fiberwise Gram
    identities and inverse-Gram identities, and the virtual telescope.  The
    left finite Gram then cancels its inverse, leaving the four physical
    sandwiches in the statement.
\end{proof}

\begin{theorem}[Whole-increment correction bounds]
    \label{thm:whole_increment_correction_bounds}
    \lean{MPSTensor.wholeIncrementTailFiniteGramCorrectionES_norm_le}
    \lean{MPSTensor.wholeIncrementFullInverseGramCorrectionES_norm_le}
    \lean{MPSTensor.wholeIncrementLeftFiniteGramCorrectionES_norm_le}
    \leanok
    \uses{def:whole_increment_tail_gram_correction, def:whole_increment_full_inverse_gram_correction, def:whole_increment_left_gram_correction, thm:reassociated_tail_pseudoinverse_bound, thm:reassociated_tail_inverse_gram, thm:inverse_gram_identities, thm:boundary_fiberwise_map_norm, thm:inverse_gram_limiting_resolvent, thm:left_spectator_pseudoinverse_bound}
    Assume $D>0$ and let $\rho$ be positive definite.  Put
    $I_t=\mathcal K_{L+Q}^{-1}$,
    $I_f=\mathcal K_{K+L+Q}^{-1}$,
    $I_\ell=\mathcal K_{K+L}^{-1}$, and
    $I_\infty=\mathcal K_\infty^{-1}$.  Under the injectivity assumptions
    required to define the indicated inverse Grams, the tail, full, and left
    corrections $C_t,C_f,C_\ell$ satisfy
    \begin{align}
        \lVert C_t\rVert
         & \leq \sqrt{\lVert I_t\rVert}
        \lVert\mathcal K_{L+Q}-\mathcal K_\infty\rVert
        \lVert J_K^{\mathrm{tail}}\rVert
        \sqrt{\lVert I_f\rVert}.
    \end{align}
    The full correction satisfies
    \begin{align}
        \lVert C_f\rVert
         & \leq \sqrt{\lVert I_t\rVert}
        \lVert\mathcal K_\infty\rVert
        \lVert J_K^{\mathrm{tail}}\rVert
        \lVert I_\infty\rVert
        \lVert\mathcal K_{K+L+Q}-\mathcal K_\infty\rVert
        \sqrt{\lVert I_f\rVert}.
    \end{align}
    The left correction satisfies
    \begin{align}
        \lVert C_\ell\rVert
         & \leq \sqrt{\lVert I_t\rVert}
        \lVert\mathcal K_\infty\rVert
        \lVert J_K^{\mathrm{tail}}\rVert
        \lVert I_\infty\rVert
        \lVert J_Q^{\mathrm{left}}\rVert
        \lVert\mathcal K_{K+L}-\mathcal K_\infty\rVert
        \sqrt{\lVert I_\ell\rVert}.
    \end{align}
\end{theorem}

\begin{proof}
    \leanok
    \uses{def:whole_increment_tail_gram_correction, def:whole_increment_full_inverse_gram_correction, def:whole_increment_left_gram_correction, thm:reassociated_tail_pseudoinverse_bound, thm:reassociated_tail_inverse_gram, thm:inverse_gram_identities, thm:boundary_fiberwise_map_norm, thm:inverse_gram_limiting_resolvent, thm:left_spectator_pseudoinverse_bound}
    Unfold each correction and apply submultiplicativity.  Use the reassociated
    tail pseudoinverse estimate and the fiberwise-map norm bound in all three
    cases.  For the full correction, insert the inverse resolvent identity; for
    the left correction, use the left boundary pseudoinverse estimate.
\end{proof}

\begin{definition}[Common overlap factors and centered residual]
    \label{def:c3_centered_overlap_factors_residual}
    \lean{MPSTensor.c3LeftOverlapFactorES}
    \lean{MPSTensor.c3RightOverlapFactorES}
    \lean{MPSTensor.c3CenteredVirtualResidualES}
    \leanok
    \uses{def:gram_reshuffle, def:whole_increment_left_overlap_factor, def:whole_increment_right_overlap_factor, thm:spectator_boundary_hilbert_factorization}
    Let $A$ be an MPS tensor, let $\Omega_K$ be the set of length-$K$ words
    in the physical alphabet, and let
    $\mathcal B_K=\bigoplus_{u\in\Omega_K}\MN{D}$ with its
    Hilbert--Schmidt direct-sum norm.  Define
    \begin{align}
        F_K & :\mathcal B_K\longrightarrow
        \bigoplus_{(u,j)\in\Omega_K\times\Omega_1}\MN{D}
    \end{align}
    by applying the one-site left virtual map independently to each prefix
    fiber.  Define
    \begin{align}
        G_K & :\mathcal B_1\longrightarrow
        \bigoplus_{(u,j)\in\Omega_K\times\Omega_1}\MN{D}
    \end{align}
    by applying the length-$K$ tail virtual map independently to each
    final-site fiber.  These definitions make sense without any nonemptiness
    assumption on the virtual dimension or on either configuration set.

    When $D\geq1$ and $\rho\in\MN{D}$ is positive definite, set
    $\mathcal K_\infty=\mathscr R(P_\rho)$ and define the centered virtual
    residual by
    \begin{align}
        R_{K,l}={} & (\Gamma^{\mathrm{tail}}_{K,l+1})^\dagger
        \Gamma^{\mathrm{left}}_{K+l,1}
        -\mathcal K_\infty^{\oplus\Omega_K}J^{\mathrm{tail}}_K
        \mathcal K_\infty^{-1}(J^{\mathrm{left}}_1)^\dagger
        \mathcal K_\infty^{\oplus\Omega_1}.
    \end{align}
\end{definition}

\begin{theorem}[Centered overlap factorization and geometric bound]
    \label{thm:c3_centered_overlap_factorization_geometric}
    \lean{MPSTensor.boundaryFamilyEquiv_c3LeftOverlapFactorES_apply}
    \lean{MPSTensor.boundaryFamilyEquiv_c3RightOverlapFactorES_apply}
    \lean{MPSTensor.c3LeftOverlapFactorES_norm_le}
    \lean{MPSTensor.c3RightOverlapFactorES_norm_le}
    \lean{MPSTensor.c3_centeredVirtualResidualES_eq_overlapFactors}
    \lean{MPSTensor.c3_centeredVirtualResidualES_norm_le}
    \lean{MPSTensor.IsPrimitiveMPS.c3_centeredVirtualResidualES_norm_sq_le_geometric}
    \leanok
    \uses{def:c3_centered_overlap_factors_residual, def:ground_space_gram, def:primitive_mps, thm:whole_increment_left_overlap_factor_apply, thm:whole_increment_right_overlap_factor_apply, thm:whole_increment_overlap_factor_norms}
    For every $K\in\N$, the common overlap factors have fiber formulas
    \begin{align}
        (F_KX)_{u,j} & =A^jX_u, &
        (G_KZ)_{u,j} & =Z_jA^u,
    \end{align}
    and norm bounds
    \begin{align}
        \lVert F_K\rVert & \leq\lVert J^{\mathrm{left}}_1\rVert, &
        \lVert G_K\rVert & \leq\lVert J^{\mathrm{tail}}_K\rVert.
    \end{align}
    The formulas and bounds remain valid for empty direct sums and incur no
    factor depending on the number of configurations.

    Assume $D\geq1$ and that $\rho\in\MN{D}$ is positive definite.  The
    centered residual of Definition~\ref{def:c3_centered_overlap_factors_residual}
    factors exactly at the generally nonidentity limiting metric as
    \begin{align}
        R_{K,l}
         & =F_K^\dagger
        (\mathcal K_l-\mathcal K_\infty)^{\oplus
            (\Omega_K\times\Omega_1)}G_K,
    \end{align}
    and hence
    \begin{align}
        \lVert R_{K,l}\rVert
         & \leq \lVert J^{\mathrm{left}}_1\rVert
        \lVert\mathcal K_l-\mathcal K_\infty\rVert
        \lVert J^{\mathrm{tail}}_K\rVert.
    \end{align}
    If, in addition, $A$ is primitive with fixed-point witness $\rho$, then
    there are constants $C>0$ and $0<r<1$ such that, for every $K\in\N$ and
    every $l\geq1$,
    \begin{align}
        \lVert R_{K,l}\rVert^2
         & \leq D^3\bigl(Cr^l\bigr)^2.
    \end{align}
    All formulas in this theorem are reconstructed algebra.  They are not
    sourced formulas and assert neither a full projector-defect estimate nor
    the FNW--Nachtergaele rational coefficient.
\end{theorem}

\begin{proof}
    \leanok
    \uses{thm:c3_projector_gram_error_telescope, thm:whole_increment_left_overlap_factor_apply, thm:whole_increment_right_overlap_factor_apply, thm:whole_increment_overlap_factor_norms, thm:boundary_fiberwise_map_norm, thm:primitive_ground_space_gram_geometric_fixed_point, thm:primitive_tail_virtual_map_norm_uniform}
    Evaluate the two factors on each common direct-sum fiber.  The direct-sum
    norm identity gives the two bounds without a factor counting spectator
    configurations.  Substitute these fiber formulas into the centered
    residual identity from
    Theorem~\ref{thm:c3_projector_gram_error_telescope} to obtain the exact
    factorization.  Submultiplicativity and
    Theorem~\ref{thm:boundary_fiberwise_map_norm} give the product bound.
    Under primitivity, combine the uniform bound for
    $J^{\mathrm{tail}}_K$ with the squared geometric Gram estimate and absorb
    the fixed norm of $J^{\mathrm{left}}_1$ into the constant.
\end{proof}

\begin{theorem}[Fiberwise subtraction]
    \label{thm:boundary_fiberwise_map_sub}
    \lean{MPSTensor.boundaryFiberwiseMap_sub}
    \leanok
    Let $S$ be a finite set and let $G,H\colon V\to V$ be continuous linear
    endomorphisms of the virtual Euclidean space.  Then
    \begin{align}
        (G-H)^{\oplus S} & =G^{\oplus S}-H^{\oplus S}.
    \end{align}
\end{theorem}

\begin{proof}
    \leanok
    Evaluate both sides on each fiber indexed by $s\in S$.
\end{proof}

\begin{theorem}[Invertibility of the finite ground-space Gram]
    \label{thm:ground_space_gram_is_unit_of_injective}
    \lean{MPSTensor.groundSpaceGram_isUnit_of_groundSpaceMapES_injective}
    \leanok
    \uses{def:ground_space_gram}
    If the ground-space map $\Gamma_n^{\mathrm{ES}}$ is injective, then its Gram
    endomorphism $\mathcal K_n=(\Gamma_n^{\mathrm{ES}})^\dagger
        \Gamma_n^{\mathrm{ES}}$ is invertible.
\end{theorem}

\begin{proof}
    \leanok
    \uses{def:ground_space_gram}
    The equality
    $\langle X,\mathcal K_nX\rangle=\lVert\Gamma_n^{\mathrm{ES}}X\rVert^2$
    shows that $\mathcal K_n$ is injective.  An injective endomorphism of a
    finite-dimensional space is bijective and hence invertible.
\end{proof}

\begin{theorem}[Finite-to-limiting inverse resolvent identity]
    \label{thm:inverse_gram_limiting_resolvent}
    \lean{MPSTensor.inverseGram_sub_limitingInverse_eq_resolvent}
    \leanok
    \uses{def:ground_space_gram, def:inverse_gram, def:gram_reshuffle}
    Assume $D\geq1$ and let $\rho\in\MN{D}$ be positive definite.  Set
    $\mathcal K_\infty=\mathscr R(P_\rho)$ and
    $S_\infty=\mathcal K_\infty^{-1}$.  If
    $\Gamma_n^{\mathrm{ES}}$ is injective and
    $S_n=\mathcal K_n^{-1}$, then
    \begin{align}
        S_n-S_\infty
         & =S_\infty(\mathcal K_\infty-\mathcal K_n)S_n.
    \end{align}
\end{theorem}

\begin{proof}
    \leanok
    \uses{thm:ground_space_gram_is_unit_of_injective, thm:fixed_point_gram_metric_is_unit, thm:inverse_gram_identities}
    Both Gram endomorphisms are invertible.  Apply the identity
    $A^{-1}-B^{-1}=B^{-1}(B-A)A^{-1}$ with
    $A=\mathcal K_n$ and $B=\mathcal K_\infty$.
\end{proof}

\begin{theorem}[Left finite-Gram cancellation]
    \label{thm:left_finite_gram_cancellation}
    \lean{MPSTensor.leftFiniteGramCancellation_eq_groundSpaceMapES_adjoint}
    \leanok
    \uses{def:left_boundary_map, def:ground_space_gram, def:inverse_gram}
    Suppose that $\Gamma_K^{\mathrm{ES}}$ is injective and write
    $S_K=\mathcal K_K^{-1}$.  Then
    \begin{align}
        (J_L^{\mathrm{left}})^\dagger
        \mathcal K_K^{\oplus\Omega_L}S_K^{\oplus\Omega_L}
        (\Gamma_{K,L}^{\mathrm{left}})^\dagger
         & =(\Gamma_{K+L}^{\mathrm{ES}})^\dagger.
    \end{align}
\end{theorem}

\begin{proof}
    \leanok
    \uses{thm:inverse_gram_identities, thm:spectator_boundary_hilbert_factorization}
    The direct-sum Gram and inverse Gram cancel fiberwise.  The remaining
    composition is the adjoint of
    $\Gamma_{K,L}^{\mathrm{left}}J_L^{\mathrm{left}}
        =\Gamma_{K+L}^{\mathrm{ES}}$.
\end{proof}

\begin{theorem}[Tail spectator pseudoinverse bound]
    \label{thm:tail_spectator_pseudoinverse_bound}
    \lean{MPSTensor.norm_tailBoundaryMapES_comp_inverseGram_le_sqrt}
    \leanok
    \uses{def:tail_boundary_map, def:inverse_gram}
    If $\Gamma_L^{\mathrm{ES}}$ is injective, then
    \begin{align}
        \left\lVert\Gamma_{K,L}^{\mathrm{tail}}
        S_{\Gamma_{K,L}^{\mathrm{tail}}}\right\rVert
         & \leq\sqrt{\lVert S_L\rVert}.
    \end{align}
\end{theorem}

\begin{proof}
    \leanok
    \uses{thm:inverse_gram_identities, thm:spectator_boundary_direct_sum_gram, thm:boundary_fiberwise_map_norm}
    The square-root pseudoinverse identity gives the left-hand side as the
    square root of the norm of the tail boundary inverse Gram.  This inverse
    Gram is the fiberwise lift of $S_L$, whose norm is at most
    $\lVert S_L\rVert$.
\end{proof}

\begin{theorem}[Left spectator pseudoinverse bound]
    \label{thm:left_spectator_pseudoinverse_bound}
    \lean{MPSTensor.norm_inverseGram_comp_leftBoundaryMapES_adjoint_le_sqrt}
    \leanok
    \uses{def:left_boundary_map, def:inverse_gram}
    If $\Gamma_K^{\mathrm{ES}}$ is injective, then
    \begin{align}
        \left\lVert S_{\Gamma_{K,L}^{\mathrm{left}}}
        (\Gamma_{K,L}^{\mathrm{left}})^\dagger\right\rVert
         & \leq\sqrt{\lVert S_K\rVert}.
    \end{align}
\end{theorem}

\begin{proof}
    \leanok
    \uses{thm:inverse_gram_identities, thm:spectator_boundary_direct_sum_gram, thm:boundary_fiberwise_map_norm}
    The adjoint-side square-root identity reduces the left-hand side to the
    norm of the left boundary inverse Gram.  The direct-sum inverse-Gram formula
    and the fiberwise norm estimate give the result.
\end{proof}

\begin{definition}[Finite-Gram corrections in the C3 residual]
    \label{def:c3_correction_sandwiches}
    \lean{MPSTensor.c3TailFiniteGramCorrectionES}
    \lean{MPSTensor.c3FullInverseGramCorrectionES}
    \lean{MPSTensor.c3LeftFiniteGramCorrectionES}
    \leanok
    \uses{def:ground_space_gram, def:inverse_gram, def:gram_reshuffle, def:tail_boundary_map, def:left_boundary_map, thm:c3_projector_gram_error_telescope}
    Assume $D\geq1$ and let $\rho\in\MN{D}$ be positive definite.  Set
    $\mathcal K_\infty=\mathscr R(P_\rho)$ and
    $S_\infty=\mathcal K_\infty^{-1}$.  Suppose that
    $\Gamma_{l+1}^{\mathrm{ES}}$, $\Gamma_{K+l}^{\mathrm{ES}}$, and
    $\Gamma_{K+l+1}^{\mathrm{ES}}$ are injective, and write
    \begin{align}
        T & =\Gamma^{\mathrm{tail}}_{K,l+1}, \\
        L & =\Gamma^{\mathrm{left}}_{K+l,1}, \\
        C & =\Gamma_{K+l+1}^{\mathrm{ES}}.
    \end{align}
    The three correction sandwiches in
    Theorem~\ref{thm:c3_projector_gram_error_telescope} are
    \begin{align}
        Q_1={} & T S_{l+1}^{\oplus\Omega_K}
        (\mathcal K_{l+1}-\mathcal K_\infty)^{\oplus\Omega_K}
        J_K^{\mathrm{tail}}S_{K+l+1}(J_1^{\mathrm{left}})^\dagger
        \mathcal K_{K+l}^{\oplus\Omega_1}
        S_{K+l}^{\oplus\Omega_1}L^\dagger,  \\
        Q_2={} & T S_{l+1}^{\oplus\Omega_K}
        \mathcal K_\infty^{\oplus\Omega_K}J_K^{\mathrm{tail}}
        (S_{K+l+1}-S_\infty)(J_1^{\mathrm{left}})^\dagger
        \mathcal K_{K+l}^{\oplus\Omega_1}
        S_{K+l}^{\oplus\Omega_1}L^\dagger,  \\
        Q_3={} & T S_{l+1}^{\oplus\Omega_K}
        \mathcal K_\infty^{\oplus\Omega_K}J_K^{\mathrm{tail}}S_\infty
        (J_1^{\mathrm{left}})^\dagger
        (\mathcal K_{K+l}-\mathcal K_\infty)^{\oplus\Omega_1}
        S_{K+l}^{\oplus\Omega_1}L^\dagger.
    \end{align}
\end{definition}

\begin{theorem}[Exact factorizations of the C3 corrections]
    \label{thm:c3_correction_factorizations}
    \lean{MPSTensor.c3TailFiniteGramCorrectionES_eq_factored}
    \lean{MPSTensor.c3FullInverseGramCorrectionES_eq_factored}
    \lean{MPSTensor.c3LeftFiniteGramCorrectionES_eq_factored}
    \leanok
    \uses{def:c3_correction_sandwiches}
    Under the hypotheses of
    Definition~\ref{def:c3_correction_sandwiches}, the three corrections have
    the exact factorizations
    \begin{align}
        Q_1={} & T S_{l+1}^{\oplus\Omega_K}
        (\mathcal K_{l+1}-\mathcal K_\infty)^{\oplus\Omega_K}
        J_K^{\mathrm{tail}}S_{K+l+1}C^\dagger, \\
        Q_2={} & T S_{l+1}^{\oplus\Omega_K}
        \mathcal K_\infty^{\oplus\Omega_K}J_K^{\mathrm{tail}}S_\infty
        (\mathcal K_\infty-\mathcal K_{K+l+1})
        S_{K+l+1}C^\dagger,                    \\
        Q_3={} & T S_{l+1}^{\oplus\Omega_K}
        \mathcal K_\infty^{\oplus\Omega_K}J_K^{\mathrm{tail}}S_\infty
        (J_1^{\mathrm{left}})^\dagger
        (\mathcal K_{K+l}-\mathcal K_\infty)^{\oplus\Omega_1}
        S_{\Gamma_{K+l,1}^{\mathrm{left}}}L^\dagger.
    \end{align}
\end{theorem}

\begin{proof}
    \leanok
    \uses{thm:boundary_fiberwise_map_sub, thm:inverse_gram_limiting_resolvent, thm:left_finite_gram_cancellation, thm:spectator_boundary_direct_sum_gram}
    The direct-sum inverse-Gram formula identifies the initial factor with the
    pseudoinverse factor of $T$.  Theorem~\ref{thm:left_finite_gram_cancellation}
    replaces the final product in $Q_1$ and $Q_2$ by $C^\dagger$.
    For $Q_2$, apply the resolvent identity
    \begin{align}
        S_{K+l+1}-S_\infty
         & =S_\infty(\mathcal K_\infty-\mathcal K_{K+l+1})S_{K+l+1}.
    \end{align}
    Fiberwise subtraction gives the displayed Gram differences.  The
    factorization of $Q_3$ leaves the left pseudoinverse adjacent to
    $L^\dagger$.
\end{proof}

% These raw bounds are reconstructed from the inverse-Gram projector formulas.
% They do not assert the rational coefficient in Nachtergaele's equation (6.1).
\begin{theorem}[Raw norm bounds for the C3 corrections]
    \label{thm:c3_correction_bounds}
    \lean{MPSTensor.c3TailFiniteGramCorrectionES_norm_le}
    \lean{MPSTensor.c3FullInverseGramCorrectionES_norm_le}
    \lean{MPSTensor.c3LeftFiniteGramCorrectionES_norm_le}
    \leanok
    \uses{def:c3_correction_sandwiches}
    Under the hypotheses of
    Definition~\ref{def:c3_correction_sandwiches},
    \begin{align}
        \lVert Q_1\rVert
         & \leq \sqrt{\lVert S_{l+1}\rVert}
        \lVert\mathcal K_{l+1}-\mathcal K_\infty\rVert
        \lVert J_K^{\mathrm{tail}}\rVert
        \sqrt{\lVert S_{K+l+1}\rVert},      \\
        \lVert Q_2\rVert
         & \leq \sqrt{\lVert S_{l+1}\rVert}
        \lVert\mathcal K_\infty\rVert
        \lVert J_K^{\mathrm{tail}}\rVert
        \lVert S_\infty\rVert
        \lVert\mathcal K_{K+l+1}-\mathcal K_\infty\rVert
        \sqrt{\lVert S_{K+l+1}\rVert},      \\
        \lVert Q_3\rVert
         & \leq \sqrt{\lVert S_{l+1}\rVert}
        \lVert\mathcal K_\infty\rVert
        \lVert J_K^{\mathrm{tail}}\rVert
        \lVert S_\infty\rVert
        \lVert J_1^{\mathrm{left}}\rVert
        \lVert\mathcal K_{K+l}-\mathcal K_\infty\rVert
        \sqrt{\lVert S_{K+l}\rVert}.
    \end{align}
\end{theorem}

\begin{proof}
    \leanok
    \uses{thm:c3_correction_factorizations, thm:tail_spectator_pseudoinverse_bound, thm:left_spectator_pseudoinverse_bound, thm:boundary_fiberwise_map_norm, thm:inverse_gram_identities}
    Apply submultiplicativity to the factorizations in
    Theorem~\ref{thm:c3_correction_factorizations}.  The two spectator
    pseudoinverse estimates control the exterior factors, and
    Theorem~\ref{thm:boundary_fiberwise_map_norm} controls every fiberwise Gram
    factor without a term depending on $|\Omega_K|$.  The square-root
    pseudoinverse identity controls $S_{K+l+1}C^\dagger$.
\end{proof}

\begin{theorem}[Uniform decay of the individual-block projector error]
    \label{thm:primitive_whole_increment_defect_decay}
    \lean{MPSTensor.IsPrimitiveMPS.eventually_wholeIncrement_groundProjection_defect_le}
    \leanok
    \uses{def:primitive_mps, def:reassociated_tail_boundary_map,
        def:spectator_boundary_maps_es, def:ground_space_es}
    Let $A$ be a normalized primitive tensor with nonzero bond dimension and a
    positive-definite fixed point.
    For every $\varepsilon>0$, all sufficiently large $L$ and every $K\geq0$, $Q>0$ satisfy
    \begin{align}
        \left\lVert P_{\mathcal H_K\otimes G_{L+Q}(A)}
        P_{G_{K+L}(A)\otimes\mathcal H_Q}-P_{G_{K+L+Q}(A)}\right\rVert
         & \leq\varepsilon
    \end{align}
    This is the uniform decay of the individual-block
    errors used in \cite[Section~6]{Nachtergaele1996Spectral}.
\end{theorem}
\begin{proof}
    \leanok
    \uses{thm:primitive_fnw_projector_defect_geometric}
    The geometric FNW estimate bounds the norm by
    $c\lambda^L(1+c\lambda^L)/(1-c\lambda^L)$ once $L$ exceeds the
    injectivity length and $c\lambda^L<1$. Since $0<\lambda<1$, this coefficient
    tends to zero. Neither the coefficient nor the lower bound on $L$ depends
    on $K$ or $Q$.
\end{proof}

\begin{theorem}[Uniform whole-increment defect at the coefficient $7/16$]
    \label{thm:uniform_whole_increment_defect_seven_sixteenths}
    \lean{MPSTensor.IsPrimitiveMPS.exists_uniform_wholeIncrement_defect_le_seven_sixteenths}
    \leanok
    \uses{def:primitive_mps, def:l_blk_injective,
        def:reassociated_tail_boundary_map, def:spectator_boundary_maps_es,
        def:ground_space_es}
    Let $A$ be primitive with positive-definite fixed point $\rho$ and nonzero
    bond dimension.  There is an integer $p>0$, measured in sites of $A$, such
    that $A$ is $p$-block injective and, for every prefix length $K$,
    \begin{align}
        \left\lVert
        P^{\mathrm{tail}}_{K;p,p}P^{\mathrm{left}}_{K+p,p}
        -P_{\mathcal G_{K+2p}^{\mathrm{ES}}(A)}
        \right\rVert
         & \leq \frac{7}{16}.
    \end{align}
\end{theorem}

\begin{proof}
    \leanok
    \uses{thm:primitive_whole_increment_defect_decay, thm:normal_from_primitive_posdef,
        thm:wordspan_blk_inj_succ}
    Primitivity supplies a positive injective length. Choose $p$ above it
    and large enough for the uniform decay estimate with tolerance $7/16$.
    Block injectivity propagates to $p$. Set the suffix length equal to $p$
    in that estimate, leaving the prefix length unrestricted.
\end{proof}

\begin{theorem}[Three-block whole-increment defect at the coefficient $7/16$]
    \label{thm:three_block_whole_increment_defect_seven_sixteenths}
    \lean{MPSTensor.IsPrimitiveMPS.exists_threeBlock_wholeIncrement_defect_le_seven_sixteenths}
    \leanok
    \uses{thm:uniform_whole_increment_defect_seven_sixteenths}
    Under the same hypotheses, there is an integer $p>0$, measured in sites of
    $A$, such that $A$ is $p$-block injective and
    \begin{align}
        \left\lVert
        P^{\mathrm{tail}}_{p;p,p}P^{\mathrm{left}}_{2p,p}
        -P_{\mathcal G_{3p}^{\mathrm{ES}}(A)}
        \right\rVert
         & \leq \frac{7}{16}.
    \end{align}
\end{theorem}

\begin{proof}
    \leanok
    \uses{thm:uniform_whole_increment_defect_seven_sixteenths}
    Apply the prefix-uniform theorem at $K=p$.  The three consecutive regions
    then each contain exactly $p$ sites of the input tensor.
\end{proof}

\begin{theorem}[Uniform open-chain projector-defect threshold]
    \label{thm:nachtergaele_c3_blocked_length}
    \lean{MPSTensor.IsPrimitiveMPS.exists_openChain_groundProjection_defect_lt_c3_threshold}
    \lean{MPSTensor.IsPrimitiveMPS.exists_re_inner_openChain_anticommutator_ge_c3_threshold}
    \leanok
    \uses{def:primitive_mps, def:l_blk_injective, def:open_chain_ground_subspaces_es, def:ground_space_es}
    Let $A$ be a primitive MPS tensor with nonzero bond dimension and a
    positive-definite fixed point.  There are an overlap length $l>1$, measured
    in sites of $A$, and a number $\epsilon_l\geq0$ such that $A$ is
    $l$-block injective, the projector defect is bounded
    uniformly in the left length $K$ by $\epsilon_l$, and
    \begin{align}
        \epsilon_l & <\frac{1}{\sqrt{l+1}}.
        \notag
    \end{align}
    Nachtergaele's projector-defect identity following condition C3
    \cite[Eq.~(2.4)]{Nachtergaele1996Spectral}, together with the two-projection
    estimate, gives the corresponding anticommutator bound for every $K$.  If
    $A$ is already a blocked representative, one site of $A$ is one
    blocked-chain filtration step.
\end{theorem}
\begin{proof}
    \leanok
    \uses{thm:primitive_fnw_openchain_defect_geometric, thm:wordspan_blk_inj_succ, thm:nachtergaele_open_chain_projector_defect}
    The open-chain geometric estimate supplies $c>0$, $0<\lambda<1$, and a
    positive block-injectivity length.
    Since
    \begin{align}
        \sqrt{n+1}\,\lambda^n & \longrightarrow0,
    \end{align}
    continuity of inversion at $1$ gives
    \begin{align}
        \frac{c\lambda^n(1+c\lambda^n)}{1-c\lambda^n}\sqrt{n+1}
         & \longrightarrow0.
    \end{align}
    Choose $l$ above that length and above $1$ so that $c\lambda^l<1$ and the
    last display is less than $1$.  Block injectivity propagates to $l$, and the geometric
    estimate gives the uniform defect bound with
    $\epsilon_l=c\lambda^l(1+c\lambda^l)/(1-c\lambda^l)$.  The open-chain
    projector-defect theorem gives the anticommutator statement.
\end{proof}

% Source: References/cond-mat_9410110/main.tex, condition C3, lines 1083--1094.
\begin{theorem}[Literal open-chain martingale condition C3]
    \label{thm:literal_open_chain_martingale_c3}
    \lean{MPSTensor.IsPrimitiveMPS.exists_openChain_martingaleDifference_norm_lt_c3_threshold}
    \leanok
    \uses{thm:nachtergaele_c3_blocked_length, thm:exact_c3_martingale_difference_identity}
    Let $A$ be a primitive MPS tensor with nonzero bond dimension and a
    positive-definite fixed point.  There are $l>1$ and
    $\epsilon_l\geq0$, with $A$ $l$-block injective, such that for every
    $K>0$, after setting $n=K+l$ and $n_l=l+1$,
    \begin{align}
        \left\lVert
        G_{\Lambda_{n+1}\setminus\Lambda_{n-l}}E_n
        \right\rVert
         & \leq\epsilon_l.
        \label{eq:literal_open_chain_martingale_c3}
    \end{align}
    The coefficient satisfies
    \begin{align}
        \epsilon_l
         & <\frac{1}{\sqrt{l+1}}.
        \label{eq:literal_open_chain_martingale_c3_threshold}
    \end{align}
    The length $l$ and coefficient $\epsilon_l$ are exactly those chosen
    by Theorem~\ref{thm:nachtergaele_c3_blocked_length}.
\end{theorem}

\begin{proof}
    \leanok
    \uses{thm:nachtergaele_c3_blocked_length, thm:exact_c3_martingale_difference_identity}
    The uniform threshold theorem gives
    \begin{align}
        \left\lVert
        G_{\Lambda_{n+1}\setminus\Lambda_{n-l}}G_{\Lambda_n}
        -G_{\Lambda_{n+1}}
        \right\rVert
         & \leq\epsilon_l
    \end{align}
    for every $K>0$.  Rewrite its norm by
    Theorem~\ref{thm:exact_c3_martingale_difference_identity}; the existing
    strict threshold on $\epsilon_l$ is unchanged.
\end{proof}

% Source: References/cond-mat_9410110/main.tex, condition C3, lines 1083--1094;
% proof of Theorem 2.1(i), lines 1178--1259.
\begin{theorem}[Literal fixed-final-volume martingale condition C3]
    \label{thm:literal_fixed_volume_martingale_c3}
    \lean{MPSTensor.fixedAmbient_martingaleDifference_conj_rightSpectatorConfigLinearIsometryEquiv,
        MPSTensor.openIntervalGroundProjectionES_comp_martingaleDifference_conj_rightSpectator,
        MPSTensor.norm_openIntervalGroundProjectionES_comp_martingaleDifference_le_openChain,
        MPSTensor.IsPrimitiveMPS.exists_fixedAmbient_martingaleDifference_norm_lt_c3_threshold}
    \leanok
    \uses{lem:right-spectator-extension,
        thm:fixed-volume-c3-projector-conjugacy,
        thm:fixed-volume-c3-boundary-indices,
        thm:literal_open_chain_martingale_c3}
    Let $A$ be a primitive MPS tensor with nonzero bond dimension and a
    positive-definite fixed point.  There are $l>1$ and
    $\epsilon_l\geq0$, with $A$ $l$-block injective, such that for every
    final volume $N$ and every $n<N$,
    \begin{align}
        \left\lVert
        G_{\Lambda_{n+1}\setminus\Lambda_{n-l}}^{(N)}E_n^{(N)}
        \right\rVert
         & \leq\epsilon_l,
        \label{eq:literal_fixed_volume_martingale_c3} \\
        \epsilon_l
         & <\frac{1}{\sqrt{l+1}}.
        \label{eq:literal_fixed_volume_martingale_c3_threshold}
    \end{align}
    This is the same overlap length and the same coefficient as in
    Theorem~\ref{thm:literal_open_chain_martingale_c3}.
\end{theorem}

\begin{proof}
    \leanok
    \uses{lem:right-spectator-extension,
        thm:fixed-volume-c3-projector-conjugacy,
        thm:fixed-volume-c3-boundary-indices,
        thm:literal_open_chain_martingale_c3}
    If $n<l$, then $E_n^{(N)}=0$.  If $n=l$, the endpoint identity gives
    $G_{\Lambda_{l+1}}^{(N)}E_l^{(N)}=0$.  These cases use only
    $\epsilon_l\geq0$.

    Suppose that $n>l$, and set $K=n-l>0$ and $r=N-(n+1)$.  The prefix and
    whole-prefix conjugacies identify $E_n^{(N)}$ with the direct sum of the
    $(n+1)$-site martingale difference over the $r$ right spectators.  The
    interval-projector conjugacy and fiberwise composition then identify the
    product in
    (\ref{eq:literal_fixed_volume_martingale_c3}) with the direct sum of its
    open-chain counterpart.  Isometric conjugacy preserves operator norm, and
    the fiberwise extension does not increase it.  The bound therefore follows
    from Theorem~\ref{thm:literal_open_chain_martingale_c3}.
\end{proof}

% Source: References/cond-mat_9410110/main.tex, conditions C1--C3,
% lines 1030--1094; Theorem 2.1(i), lines 1119--1130; proof lines 1195--1259.
\begin{theorem}[Compatible open-chain gap from fixed-volume overlap]
    \label{thm:open_parent_hamiltonian_fixed_c3_gap}
    \lean{MPSTensor.openParentHamiltonianES_norm_gap_of_fixedAmbient_c3}
    \leanok
    \uses{thm:nachtergaele_c1_c3_full_range_norm_gap,
        thm:physical-outside-window-martingale-cancellation,
        thm:fixed-volume-c3-boundary-indices,
        thm:fixed-volume-c3-projector-conjugacy,
        thm:open_parent_hamiltonian_c1_full_range,
        thm:open_suffix_parent_hamiltonian_c2_full_range,
        lem:open_parent_hamiltonian_es_positive,
        thm:open_parent_hamiltonian_kernel_ground_space}
    Let $D\geq1$, let $l>0$, suppose that $A$ is $l$-block injective, and let
    $N\geq l+1$.  Assume $0\leq\epsilon_l<1/\sqrt{l+1}$ and, for every
    $0\leq n<N$,
    \begin{align}
        \left\lVert Q_nE_n^{(N)}\right\rVert & \leq\epsilon_l,
        \label{eq:open_parent_hamiltonian_fixed_c3}
    \end{align}
    where $Q_n$ is the ground-space projection of the interval
    $[n-l,n+1)$ and $E_n^{(N)}$ is the fixed-volume martingale difference.
    Then every $v\in(G_N^{\mathrm{ES}}(A))^\perp$ satisfies
    \begin{align}
        \left(1-\epsilon_l\sqrt{l+1}\right)^2\lVert v\rVert
         & \leq
        \left\lVert H^{\mathrm{open}}_{N,l+1}(A)v\right\rVert.
        \label{eq:open_parent_hamiltonian_fixed_c3_gap}
    \end{align}
    This intermediate MPS specialization assumes C3 on the full range
    $0\leq n<N$, rather than only above the lower threshold in the source
    statement. It applies the full-range estimate with
    $d_{l+1}=\gamma_{l+1}=1$ and retains the coefficient of Theorem~2.1(i) in
    \cite{Nachtergaele1996Spectral}. It is not a formulation of the
    unrestricted source theorem.
    Theorem~\ref{thm:primitive_open_parent_hamiltonian_gap} supplies the
    full-range C3 hypothesis from the physical open-chain estimates.
\end{theorem}

\begin{proof}
    \leanok
    \uses{thm:nachtergaele_c1_c3_full_range_norm_gap,
        thm:physical-outside-window-martingale-cancellation,
        thm:fixed-volume-c3-boundary-indices,
        thm:fixed-volume-c3-projector-conjugacy,
        thm:open_parent_hamiltonian_c1_full_range,
        thm:open_suffix_parent_hamiltonian_c2_full_range,
        lem:open_parent_hamiltonian_es_positive,
        thm:open_parent_hamiltonian_kernel_ground_space}
    Use the fixed-volume prefix ground-space projections and their differences
    $E_n^{(N)}$. Theorem~\ref{thm:fixed-volume-c3-boundary-indices} gives
    $G_0^{(N)}=I$, while
    Theorem~\ref{thm:fixed-volume-c3-projector-conjugacy} identifies
    $\operatorname{ran}G_N^{(N)}$ with the kernel of the open Hamiltonian.
    The outside-window commutation theorem gives the locality hypothesis in
    the martingale estimate. Theorems
    \ref{thm:open_parent_hamiltonian_c1_full_range} and
    \ref{thm:open_suffix_parent_hamiltonian_c2_full_range} give C1 and C2 with
    $d_{l+1}=\gamma_{l+1}=1$, while
    (\ref{eq:open_parent_hamiltonian_fixed_c3}) is C3.  Positivity of the open
    Hamiltonian and
    $\ker H^{\mathrm{open}}_{N,l+1}(A)=G_N^{\mathrm{ES}}(A)$ identify the
    final orthogonal complement.  Substitution into
    (\ref{eq:nachtergaele_c1_c3_full_range_norm_gap}) yields
    (\ref{eq:open_parent_hamiltonian_fixed_c3_gap}).
\end{proof}

% Source: References/cond-mat_9410110/main.tex, condition C3,
% lines 1083--1094; Theorem 2.1(i), lines 1119--1130; proof lines 1195--1259.
\begin{theorem}[Uniform gap for primitive compatible open chains]
    \label{thm:primitive_open_parent_hamiltonian_gap}
    \lean{MPSTensor.IsPrimitiveMPS.exists_openParentHamiltonianES_gap}
    \leanok
    \uses{thm:literal_fixed_volume_martingale_c3,
        thm:open_parent_hamiltonian_fixed_c3_gap}
    Let $A$ be a primitive MPS tensor with nonzero bond dimension and a
    positive-definite fixed point.  There exist $l>1$ and $\epsilon_l\geq0$
    such that $A$ is $l$-block injective,
    \begin{align}
        \epsilon_l & <\frac{1}{\sqrt{l+1}},
        \label{eq:primitive_open_parent_hamiltonian_threshold}
    \end{align}
    and, for every $N\geq l+1$ and every
    $v\in(G_N^{\mathrm{ES}}(A))^\perp$,
    \begin{align}
        \left(1-\epsilon_l\sqrt{l+1}\right)^2\lVert v\rVert
         & \leq
        \left\lVert H^{\mathrm{open}}_{N,l+1}(A)v\right\rVert.
        \label{eq:primitive_open_parent_hamiltonian_gap}
    \end{align}
\end{theorem}

\begin{proof}
    \leanok
    \uses{thm:literal_fixed_volume_martingale_c3,
        thm:open_parent_hamiltonian_fixed_c3_gap}
    Theorem~\ref{thm:literal_fixed_volume_martingale_c3} supplies $l>1$,
    $\epsilon_l$, $l$-block injectivity, and the C3 estimate for every final
    volume $N$ and every $n<N$.  Apply
    Theorem~\ref{thm:open_parent_hamiltonian_fixed_c3_gap} for each
    $N\geq l+1$.  The value of $\epsilon_l$ is unchanged, so the coefficient
    is exactly the specialization
    \begin{align}
        \frac{\gamma_{l+1}}{d_{l+1}}
        \left(1-\epsilon_l\sqrt{l+1}\right)^2
         & =\left(1-\epsilon_l\sqrt{l+1}\right)^2.
    \end{align}
\end{proof}

\begin{lemma}[Blocked three-site projector and adjacent-pair transport]
    \label{lem:blocked_to_cyclic_ground_projector_comparison}
    \lean{MPSTensor.blockTensor_threeSite_openChain_defect_norm_eq}
    \leanok
    \uses{thm:three_block_whole_increment_defect_seven_sixteenths, thm:friedrichs_ground_space_excitation_anticommutator, thm:three_site_anticommutator_cyclic_transport}
    Let $A$ be an MPS tensor with nonzero physical dimension, let $p$ be any
    blocking length, and set $B=\operatorname{block}_p(A)$.  Then
    \begin{align}
        \left\lVert
        P^{\mathrm{tail}}_{1,2}(B)P^{\mathrm{left}}_2(B)
        -P_{\mathcal G_3^{\mathrm{ES}}(B)}
        \right\rVert
         & =\left\lVert
        P^{\mathrm{tail}}_{p;p,p}P^{\mathrm{left}}_{2p,p}
        -P_{\mathcal G_{3p}^{\mathrm{ES}}(A)}
        \right\rVert.
    \end{align}
    If $B$ is injective, then the bound on either side
    \begin{align}
        \left\lVert
        P^{\mathrm{tail}}_{1,2}(B)P^{\mathrm{left}}_2(B)
        -P_{\mathcal G_3^{\mathrm{ES}}(B)}
        \right\rVert
         & \leq \frac{7}{16}
    \end{align}
    gives the adjacent three-site anticommutator inequality
    \begin{align}
        h_0h_1+h_1h_0 & \geq-\frac{7}{16}(h_0+h_1).
    \end{align}
    The same quadratic-form inequality holds for every forward or backward
    adjacent pair of cyclic range-two terms on every chain of length at least
    three, including pairs crossing the periodic cut.
\end{lemma}

\begin{theorem}[Uniform gap for a blocked primitive MPS parent Hamiltonian]
    \label{thm:martingale_mps}
    \lean{MPSTensor.IsPrimitiveMPS.exists_blockTensor_parentHamiltonianES_gap_eighth,
        MPSTensor.IsPrimitiveMPS.exists_blockTensor_parentHamiltonianES_gapped}
    \leanok
    \uses{def:primitive_mps, def:l_blk_injective, def:blocked_tensor, def:parent_ground_space_es, def:parent_hamiltonian_es}
    Let $A$ be a primitive MPS tensor with nonzero physical and bond dimensions,
    and suppose that its fixed point is positive definite.  There is a blocking
    length $p>0$ such that $A$ is $p$-block injective and, with
    $B=\operatorname{block}_p(A)$, for every $N\geq4$ and every
    $v\in(\ker H_{N,2}(B))^\perp$,
    \begin{align}
        \frac{1}{8}\lVert v\rVert
         & \leq \lVert H_{N,2}(B)v\rVert.
    \end{align}
\end{theorem}

\begin{proof}
    \leanok
    \uses{lem:blocked_to_cyclic_ground_projector_comparison, lem:parent_hamiltonian_cyclic_overlap_anticommutator_gap}
    Choose $p$ so that the three-block projector defect is at most $7/16$.
    Lemma~\ref{lem:blocked_to_cyclic_ground_projector_comparison} gives the same
    coefficient for every overlapping cyclic pair.  For range two each row has
    at most two such pairs, so the martingale estimate gives
    $1-2(7/16)=1/8$.
\end{proof}

\begin{theorem}[Anticommutator bound for two projections]
    \label{thm:projection_anticommutator_upper_bound}
    \lean{LinearMap.IsSymmetricProjection.re_inner_anticommutator_le}
    \leanok
    Let $P$ and $Q$ be orthogonal projections on a real or complex Hilbert
    space. For every vector $v$,
    \begin{align}
        \operatorname{Re}\langle Pv,Qv\rangle
        +\operatorname{Re}\langle Qv,Pv\rangle
         & \leq \operatorname{Re}\langle Pv,v\rangle
        +\operatorname{Re}\langle Qv,v\rangle.
    \end{align}
\end{theorem}

\begin{proof}
    \leanok
    Positivity of the inner product gives
    \begin{align}
        0
         & \leq \operatorname{Re}\langle(P-Q)v,(P-Q)v\rangle, \\
        \operatorname{Re}\langle(P-Q)v,(P-Q)v\rangle
         & =\operatorname{Re}\langle Pv,v\rangle
        -\operatorname{Re}\langle Pv,Qv\rangle
        -\operatorname{Re}\langle Qv,Pv\rangle
        +\operatorname{Re}\langle Qv,v\rangle.
    \end{align}
    Symmetry and idempotence expand the left-hand side into the second line;
    rearranging it proves the claim.
\end{proof}

\begin{lemma}[Norm gap implies the squared quadratic-form bound]
    \label{lem:norm_gap_to_squared_quadratic_form}
    \lean{LinearMap.IsPositive.quadraticForm_sq_ge_of_norm_gap}
    \leanok
    Let $H$ be a positive operator on a finite-dimensional complex Hilbert
    space and let $\gamma\geq0$. If
    \begin{align}
        \gamma\lVert u\rVert & \leq \lVert Hu\rVert
        \qquad\text{for every }u\in(\ker H)^\perp,
    \end{align}
    then every vector $v$ satisfies
    \begin{align}
        \gamma\operatorname{Re}\langle Hv,v\rangle
         & \leq \operatorname{Re}\langle Hv,Hv\rangle.
    \end{align}
\end{lemma}

\begin{proof}
    \leanok
    Put $w=v-P_{\ker H}v$. Then $w\perp\ker H$, $Hw=Hv$, and symmetry gives
    $\langle Hv,v\rangle=\langle Hv,w\rangle$. Hence
    \begin{align}
        \gamma\operatorname{Re}\langle Hv,v\rangle
         & \leq \gamma\lVert Hv\rVert\lVert w\rVert
        \leq \lVert Hv\rVert\lVert Hw\rVert
        =\operatorname{Re}\langle Hv,Hv\rangle.
    \end{align}
\end{proof}

\begin{lemma}[Right-spectator quadratic-form bound]
    \label{lem:right_spectator_quadratic_form_gap}
    \lean{ContinuousLinearMap.quadraticForm_sq_ge_rightFiberwiseMap,
        LinearMap.IsPositive.quadraticForm_sq_ge_rightFiberwiseMap_of_norm_gap}
    \leanok
    \uses{lem:norm_gap_to_squared_quadratic_form, lem:right-spectator-extension}
    Let $G$ act on a finite active configuration space. If
    $G^2\geq\gamma G$ as a quadratic form, the same inequality holds when $G$
    acts independently over any finite right spectator coordinate. In
    particular, positivity and a norm gap $\gamma$ for $G$ imply this
    fiberwise quadratic-form inequality.
\end{lemma}

\begin{proof}
    \leanok
    Write $v_s$ for the active vector at spectator value $s$. Both quadratic
    forms split into finite sums over $s$, and the active inequality gives
    \begin{align}
        \gamma\sum_s\operatorname{Re}\langle Gv_s,v_s\rangle
         & \leq \sum_s\operatorname{Re}\langle Gv_s,Gv_s\rangle.
    \end{align}
    The final statement follows from
    Lemma~\ref{lem:norm_gap_to_squared_quadratic_form}.
\end{proof}

\begin{theorem}[Cyclic window double counting]
    \label{thm:cyclic_knabe_double_counting}
    \lean{ProjectionGeometry.cyclicWindowSum,
        ProjectionGeometry.cyclicCrossTerms,
        ProjectionGeometry.reInnerCyclicCrossTerms,
        ProjectionGeometry.reInnerCyclicCrossTerms_eq_re_inner_cyclicCrossTerms,
        ZModCyclicSums.sum_comp_addRight,
        ZModCyclicSums.sum_range_Ico_triangular,
        ZModCyclicSums.sum_range_Ico_triangular_lower,
        ZModCyclicSums.sum_cyclic_window_eq,
        ProjectionGeometry.cyclicWindowSum_sq_sum_eq,
        ProjectionGeometry.re_inner_cyclicWindowSum_sum_eq,
        ProjectionGeometry.re_inner_cyclicWindowSum_sum_apply}
    \leanok
    Let $h_i$, with $i\in\mathbb Z/N\mathbb Z$, be orthogonal projections. Define
    \begin{align}
        H      & =\sum_i h_i,                                                \\
        A_s    & =\sum_{a=0}^{m-1}h_{s+a},                                   \\
        Q_q    & =\sum_i\bigl(h_i h_{i+q}+h_i h_{i-q}\bigr),                 \\
        C_q(v) & =\sum_i\bigl(\operatorname{Re}\langle h_i v,h_{i+q}v\rangle
        +\operatorname{Re}\langle h_i v,h_{i-q}v\rangle\bigr)
        =\operatorname{Re}\langle Q_qv,v\rangle.
    \end{align}
    Then, for every vector $v$,
    \begin{align}
        \sum_s A_s^2
         & =mH+\sum_{q=1}^{m-1}(m-q)Q_q,\label{eq:cyclic_knabe_window_square}   \\
        \sum_s\operatorname{Re}\langle A_sv,A_sv\rangle
         & =m\operatorname{Re}\langle Hv,v\rangle
        +\sum_{q=1}^{m-1}(m-q)C_q(v),\label{eq:cyclic_knabe_window_square_form} \\
        \sum_s\operatorname{Re}\langle A_s v,v\rangle
         & =m\operatorname{Re}\langle Hv,v\rangle.
    \end{align}
\end{theorem}

\begin{proof}
    \leanok
    For fixed offsets $a,b\in\{0,\ldots,m-1\}$, translation by $a$ gives
    \begin{align}
        \sum_s h_{s+a}h_{s+b}
         & =\sum_i h_i h_{i+b-a}.\label{eq:cyclic_knabe_reindex}
    \end{align}
    The case $a=b$ contributes $mH$, while each of the differences
    $b-a=q$ and $b-a=-q$ contributes $(m-q)Q_q$. Summing
    (\ref{eq:cyclic_knabe_reindex}) over $a$ and $b$ gives
    (\ref{eq:cyclic_knabe_window_square}). Taking the quadratic form of
    that operator identity at $v$ turns each product $h_ih_{i\pm q}$ into
    the pair $\operatorname{Re}\langle h_iv,h_{i\pm q}v\rangle$, giving
    (\ref{eq:cyclic_knabe_window_square_form}). Translation invariance also gives
    \begin{align}
        \sum_s \operatorname{Re}\langle A_s v,v\rangle
         & =m\operatorname{Re}\langle Hv,v\rangle.
        \label{eq:cyclic_knabe_window_linear}
    \end{align}
\end{proof}

\begin{theorem}[Cyclic cross-term estimates]
    \label{thm:cyclic_knabe_cross_estimates}
    \lean{ZModCyclicSums.sum_offset_partition,
        ProjectionGeometry.re_inner_cyclicCrossTerms_le_two,
        ProjectionGeometry.re_inner_cyclicCrossTerms_nonneg,
        ProjectionGeometry.re_inner_sum_two_ge_sum_cyclicCrossTerms}
    \leanok
    \uses{thm:cyclic_knabe_double_counting}
    For $Q_q$ as in Theorem~\ref{thm:cyclic_knabe_double_counting},
    \begin{align}
        \operatorname{Re}\langle Q_qv,v\rangle
         & \leq2\operatorname{Re}\langle Hv,v\rangle.
    \end{align}
    Suppose that $h_i$ and $h_{i+d}$ commute whenever
    $R\leq d\leq N-R$. Then the cross term is nonnegative for these offsets.
    If $1\leq m$, $R\leq m$, and $2m\leq N$, it follows that
    \begin{align}
        \operatorname{Re}\langle Hv,v\rangle
        +\sum_{q=1}^{m-1}\operatorname{Re}\langle Q_qv,v\rangle
         & \leq\operatorname{Re}\langle Hv,Hv\rangle.
    \end{align}
\end{theorem}

\begin{proof}
    \leanok
    \uses{thm:projection_anticommutator_upper_bound, thm:cyclic_knabe_double_counting, rem:projection_geometry_rowsum_identities}
    Translation of the reversed leg gives
    \begin{align}
        \sum_i\operatorname{Re}\langle h_{i+q}v,h_i v\rangle
         & =\sum_i\operatorname{Re}\langle h_i v,h_{i-q}v\rangle.
        \label{eq:cyclic_knabe_reversed_leg}
    \end{align}
    Summing Theorem~\ref{thm:projection_anticommutator_upper_bound} over $i$
    and using (\ref{eq:cyclic_knabe_reversed_leg}) proves the upper bound.
    If $h_i$ and $h_{i+d}$ commute, their symmetry and idempotence give
    \begin{align}
        0 & \leq\operatorname{Re}\langle h_i v,h_{i+d}v\rangle.
        \label{eq:cyclic_knabe_commuting_nonnegative}
    \end{align}
    Partitioning the nonzero cyclic offsets into $1\leq q<m$, their
    reflections $N-q$, and $m\leq k\leq N-m$ gives
    \begin{align}
        \operatorname{Re}\langle Hv,Hv\rangle
         & =\operatorname{Re}\langle Hv,v\rangle
        +\sum_{q=1}^{m-1}\operatorname{Re}\langle Q_qv,v\rangle
        +\sum_{k=m}^{N-m}\sum_i
        \operatorname{Re}\langle h_i v,h_{i+k}v\rangle.
        \label{eq:cyclic_knabe_offset_partition}
    \end{align}
    The last sum in (\ref{eq:cyclic_knabe_offset_partition}) is nonnegative by
    (\ref{eq:cyclic_knabe_commuting_nonnegative}), which proves the lower bound.
\end{proof}

\begin{theorem}[Finite-range cyclic Knabe inequality]
    \label{thm:finite_range_cyclic_knabe}
    \lean{ProjectionGeometry.quadraticForm_sum_projections_of_cyclic_knabe}
    \leanok
    \uses{thm:cyclic_knabe_double_counting}
    Let $1\leq R\leq m$, $2m\leq N$, and let $h_i$ be cyclically indexed
    orthogonal projections such that $h_i$ commutes with $h_{i+d}$ whenever
    $R\leq d\leq N-R$. Suppose that, for every $s$ and $v$,
    \begin{align}
        \gamma\operatorname{Re}\langle A_sv,v\rangle
         & \leq\operatorname{Re}\langle A_sv,A_sv\rangle.
    \end{align}
    If $(R-1)^2<m\gamma$, then, with
    \begin{align}
        \delta & =\frac{m\gamma-(R-1)^2}{m-R+1}>0,
    \end{align}
    one has
    \begin{align}
        \delta\operatorname{Re}\langle Hv,v\rangle
         & \leq\operatorname{Re}\langle Hv,Hv\rangle.
    \end{align}
    The coefficient for general $R$ is derived here and is not attributed to
    the cited papers.
    % Its source comparison is recorded in
    % docs/paper-gaps/knabe88_finite_range_coefficient.tex.
\end{theorem}

\begin{proof}
    \leanok
    \uses{thm:cyclic_knabe_double_counting, thm:cyclic_knabe_cross_estimates}
    Put $c=m-R+1$. Summing the window gaps and applying
    Theorem~\ref{thm:cyclic_knabe_double_counting} gives
    \begin{align}
        m\gamma\operatorname{Re}\langle Hv,v\rangle
         & \leq\sum_s\operatorname{Re}\langle A_sv,A_sv\rangle.
    \end{align}
    Decompose $m-q=c+(R-1-q)$ and expand
    (\ref{eq:cyclic_knabe_window_square_form}): the weights $q\leq R-2$
    carry the coefficient $c$ plus the positive part $R-1-q$, while for
    $q\geq R$ the coefficient $R-1-q$ is nonpositive and the cross term is
    nonnegative, so those summands are bounded by $c\,C_q(v)$. For the
    positive parts, the anticommutator bound and the triangular identity
    \begin{align}
        2\sum_{q=1}^{R-1}(R-1-q) & =(R-1)(R-2)
    \end{align}
    combine with the decomposition
    \begin{align}
        m-c+(R-1)(R-2) & =(R-1)^2,
        \qquad c=m-R+1,
        \label{eq:cyclic_knabe_coefficient}
    \end{align}
    so that
    \begin{align}
        \sum_s\operatorname{Re}\langle A_sv,A_sv\rangle
         & \leq\bigl(m+(R-1)(R-2)\bigr)\operatorname{Re}\langle Hv,v\rangle
        +c\sum_{q=1}^{m-1}C_q(v)
        \notag                                                              \\
         & \leq\bigl(m+(R-1)(R-2)\bigr)\operatorname{Re}\langle Hv,v\rangle
        +c\bigl(\operatorname{Re}\langle Hv,Hv\rangle
        -\operatorname{Re}\langle Hv,v\rangle\bigr).
        \notag
    \end{align}
    The cyclic cross-term estimate therefore gives
    \begin{align}
        \sum_s\operatorname{Re}\langle A_sv,A_sv\rangle
         & \leq c\operatorname{Re}\langle Hv,Hv\rangle
        +(R-1)^2\operatorname{Re}\langle Hv,v\rangle.
    \end{align}
    Rearranging and dividing by $c>0$ proves the claim.
\end{proof}

\begin{theorem}[Nearest-neighbor cyclic Knabe inequality]
    \label{thm:nearest_neighbor_cyclic_knabe}
    \lean{ProjectionGeometry.quadraticForm_sum_projections_of_cyclic_knabe_nearest_neighbor}
    \leanok
    \uses{thm:finite_range_cyclic_knabe}
    Under the hypotheses of Theorem~\ref{thm:finite_range_cyclic_knabe} with
    $R=2$, if $1<m\gamma$, then
    \begin{align}
        \frac{m\gamma-1}{m-1}\operatorname{Re}\langle Hv,v\rangle
         & \leq\operatorname{Re}\langle Hv,Hv\rangle.
    \end{align}
    Equivalently, the open-window threshold is $\gamma>1/m$. This is Knabe's
    nearest-neighbor threshold~\cite{Knabe1988Energy}, quoted in
    \cite[lines~1483--1489]{PerezGarcia2007Matrix} and
    \cite[lines~2194--2197]{Cirac2021Matrix}.
\end{theorem}

\begin{proof}
    \leanok
    \uses{thm:finite_range_cyclic_knabe}
    Substitute $R=2$ into the coefficient of
    Theorem~\ref{thm:finite_range_cyclic_knabe}.
\end{proof}

% Source scope: Papers/quant-ph_0608197/MPSarchive.tex:1460--1500 gives
% the gap definition, the MPS gap assertion, and the nearest-neighbor Knabe
% threshold. Papers/2011.12127/TN-Review-main.tex:2184--2188,2194--2197 gives
% the MPS gap assertion and the nearest-neighbor Knabe principle. The arbitrary
% finite-range coefficient below is the TNLean derivation documented in
% docs/paper-gaps/knabe88_finite_range_coefficient.tex.

\begin{definition}[Cyclic coordinates for periodic local interactions]
    \label{def:cyclic_parent_local_coordinates}
    \lean{MPSTensor.zmodLocalTermES}
    \leanok
    \uses{rem:euclidean_local_projector_ingredients}
    For an $N$-site periodic chain with $N>0$, index the range-$R$
    local interactions by $\mathbb Z/N\mathbb Z$.  If
    $s\in\mathbb Z/N\mathbb Z$ and $i\in\{0,\ldots,N-1\}$ is its standard
    representative, write
    \begin{align}
        h_s^{(N,R)} & =h_{i,\mathrm{ES}}(A,R).
    \end{align}
\end{definition}

\begin{theorem}[Total sum in cyclic coordinates]
    \label{thm:cyclic_parent_local_sum}
    \lean{MPSTensor.sum_zmodLocalTermES_eq_parentHamiltonianES}
    \leanok
    \uses{def:cyclic_parent_local_coordinates, def:parent_hamiltonian_es}
    The sum over the cyclic coordinates is the periodic parent Hamiltonian:
    \begin{align}
        \sum_{s\in\mathbb Z/N\mathbb Z}h_s^{(N,R)}
         & =H_{N,R}^{\mathrm{ES}}(A).
        \label{eq:cyclic_parent_local_sum}
    \end{align}
\end{theorem}

\begin{proof}
    \leanok
    \uses{def:cyclic_parent_local_coordinates, def:parent_hamiltonian_es}
    Reindex the standard sum over $\{0,\ldots,N-1\}$ by its canonical
    bijection with $\mathbb Z/N\mathbb Z$.  Each summand is unchanged.
\end{proof}

\begin{theorem}[Oriented separation of cyclic interactions]
    \label{thm:cyclic_parent_oriented_separation}
    \lean{MPSTensor.finEquiv_symm_add_eq_cyclicForwardSite,
        MPSTensor.cyclicWindowsDisjoint_cyclicForwardSite_of_oriented_separation,
        MPSTensor.zmodLocalTermES_commute_of_oriented_separation}
    \leanok
    \uses{def:cyclic_parent_local_coordinates, rem:cyclic_window_operators,
        lem:cyclic_window_nonoverlap_positive}
    Let $R\leq e$ and $e+R\leq N$.  Addition by $e$ in
    $\mathbb Z/N\mathbb Z$ moves the standard representative forward by $e$
    sites.  The two length-$R$ cyclic windows beginning at $s$ and $s+e$ are
    disjoint.  Consequently,
    \begin{align}
        h_s^{(N,R)}h_{s+e}^{(N,R)}
         & =h_{s+e}^{(N,R)}h_s^{(N,R)}.
        \label{eq:cyclic_parent_oriented_commutation}
    \end{align}
\end{theorem}

\begin{proof}
    \leanok
    \uses{def:cyclic_parent_local_coordinates, rem:cyclic_window_operators,
        lem:cyclic_window_nonoverlap_positive}
    Measure every site from the first window's initial site in the forward
    cyclic order.  Membership in both windows would give an offset smaller
    than $R$ and also an offset of the form $e+b$ with $b<R$.  The inequalities
    $R\leq e$ and $e+R\leq N$ exclude this possibility.  Local interactions on
    disjoint cyclic windows commute, which gives
    (\ref{eq:cyclic_parent_oriented_commutation}).
\end{proof}

\begin{definition}[Cyclic active-block coordinates]
    \label{def:cyclic_active_block_coordinates}
    \lean{MPSTensor.cyclicActiveBlockConfigEquiv,
        MPSTensor.cyclicActiveBlockConfigLinearIsometryEquiv,
        MPSTensor.cyclicActiveBlockConfigEquiv_symm_apply_window,
        MPSTensor.cyclicActiveBlockConfigLinearIsometryEquiv_apply_apply,
        MPSTensor.cyclicActiveBlockConfigLinearIsometryEquiv_symm_apply_apply,
        MPSTensor.cyclicActiveBlockConfigEquiv_symm_apply_spectator}
    \leanok
    \uses{def:cyclic_window_complement_indexing}
    Let $W\leq N$ and fix a site $i$.  Splitting the cyclically ordered sites
    into the first $W$ sites and their complement gives
    \begin{align}
        [d]^N & \simeq [d]^W\times[d]^{N-W},                 \\
        U_i\colon\ell^2([d]^N)
              & \longrightarrow\ell^2([d]^W\times[d]^{N-W}).
        \label{eq:cyclic_active_block_isometry}
    \end{align}
    If $(\sigma,\tau)$ is an active configuration and a complementary
    configuration, then their inverse image has value $\sigma_r$ at the site
    $i+r$ for $0\leq r<W$, and value $\tau_r$ at the site $i+W+r$ for
    $0\leq r<N-W$.  The isometry and its inverse act by composition with this
    configuration bijection.
\end{definition}

\begin{theorem}[A local interaction inside a cyclic active block]
    \label{thm:cyclic_active_block_local_conjugacy}
    \lean{MPSTensor.localTermES_conj_cyclicActiveBlockConfigLinearIsometryEquiv}
    \leanok
    \uses{def:cyclic_active_block_coordinates,
        rem:euclidean_local_projector_ingredients, lem:right-spectator-extension}
    Suppose that $W\leq N$ and $q+R\leq W$.  Under the isometry $U_i$ in
    (\ref{eq:cyclic_active_block_isometry}), the range-$R$ interaction
    beginning at the cyclic site $i+q$ becomes the corresponding nonwrapping
    interaction on the active $W$-site interval, acting independently on the
    complementary configuration:
    \begin{align}
        U_i h_{i+q}^{(N,R)}U_i^{-1}
         & =\bigl(h_q^{(W,R)}\bigr)^{\oplus[d]^{N-W}}.
        \label{eq:cyclic_active_block_local_conjugacy}
    \end{align}
\end{theorem}

\begin{proof}
    \leanok
    \uses{def:cyclic_active_block_coordinates,
        rem:euclidean_local_projector_ingredients, lem:right-spectator-extension}
    The condition $q+R\leq W$ places the whole interaction window inside the
    active block.  Joining an active configuration to a complementary one
    therefore preserves both the extracted length-$R$ word and every
    configuration obtained by replacing that word.  The local projector acts
    only on the active coordinates, which proves
    (\ref{eq:cyclic_active_block_local_conjugacy}) on each complementary
    configuration.
\end{proof}

\begin{theorem}[A cyclic Knabe block is an open parent Hamiltonian]
    \label{thm:cyclic_knabe_block_open_parent_hamiltonian}
    \lean{MPSTensor.cyclicWindowSum_zmodLocalTermES_conj_cyclicActiveBlock}
    \leanok
    \uses{def:cyclic_parent_local_coordinates,
        def:cyclic_active_block_coordinates,
        thm:cyclic_active_block_local_conjugacy,
        def:open_parent_hamiltonian_es}
    Let $1\leq R\leq m$ and $2m\leq N$, and put
    \begin{align}
        W & =m+R-1.
        \label{eq:cyclic_knabe_active_volume}
    \end{align}
    For every $s\in\mathbb Z/N\mathbb Z$, let
    $A_s=\sum_{q=0}^{m-1}h_{s+q}^{(N,R)}$.  The $m$ interaction windows
    occupy exactly the $W$ cyclic sites beginning at $s$.  Under the active
    block isometry $U_s$,
    \begin{align}
        U_sA_sU_s^{-1}
         & =\left(H^{\mathrm{open}}_{W,R}(A)\right)^{\oplus[d]^{N-W}}.
        \label{eq:cyclic_knabe_block_open_hamiltonian}
    \end{align}
\end{theorem}

\begin{proof}
    \leanok
    \uses{thm:cyclic_active_block_local_conjugacy,
        def:open_parent_hamiltonian_es}
    The starts $q=0,\ldots,m-1$ are precisely the nonwrapping range-$R$
    starts in an interval of length $W=m+R-1$.  Apply
    Theorem~\ref{thm:cyclic_active_block_local_conjugacy} to every summand and
    reindex the sum by these nonwrapping starts.  Fiberwise extension commutes
    with finite sums, giving
    (\ref{eq:cyclic_knabe_block_open_hamiltonian}).
\end{proof}

\begin{theorem}[Periodic finite-range gap from one open active volume]
    \label{thm:parent_hamiltonian_gap_from_open_knabe_block}
    \lean{MPSTensor.parentHamiltonianES_gap_of_openParentHamiltonianES_gap}
    \leanok
    \uses{def:open_parent_hamiltonian_es, def:parent_hamiltonian_es,
        thm:cyclic_parent_local_sum, thm:cyclic_parent_oriented_separation,
        thm:cyclic_knabe_block_open_parent_hamiltonian,
        lem:right_spectator_quadratic_form_gap,
        thm:finite_range_cyclic_knabe, lem:quadratic_form_to_norm_bound}
    Let $1\leq R\leq m$, put $W=m+R-1$, and suppose that
    $H^{\mathrm{open}}_{W,R}(A)$ has norm gap $\gamma$.  Thus, for every
    $u\in(\ker H^{\mathrm{open}}_{W,R}(A))^\perp$,
    \begin{align}
        \gamma\lVert u\rVert
         & \leq\left\lVert H^{\mathrm{open}}_{W,R}(A)u\right\rVert.
        \label{eq:open_knabe_block_norm_gap}
    \end{align}
    If $(R-1)^2<m\gamma$, define
    \begin{align}
        \delta
         & =\frac{m\gamma-(R-1)^2}{m-R+1}.
        \label{eq:parent_hamiltonian_finite_range_knabe_delta}
    \end{align}
    Then $\delta>0$, and for every $N\geq2m$ and every
    $v\in(\ker H_{N,R}^{\mathrm{ES}}(A))^\perp$,
    \begin{align}
        \delta\lVert v\rVert
         & \leq\left\lVert H_{N,R}^{\mathrm{ES}}(A)v\right\rVert.
        \label{eq:parent_hamiltonian_finite_range_knabe_gap}
    \end{align}
    The coefficient in
    (\ref{eq:parent_hamiltonian_finite_range_knabe_delta}) is the TNLean
    finite-range derivation documented in
    \path{docs/paper-gaps/knabe88_finite_range_coefficient.tex}.  It is not a
    theorem of the
    MPS sources cited for the gap definition and the nearest-neighbor Knabe
    principle.
\end{theorem}

\begin{proof}
    \leanok
    \uses{thm:cyclic_parent_local_sum,
        thm:cyclic_parent_oriented_separation,
        thm:cyclic_knabe_block_open_parent_hamiltonian,
        lem:right_spectator_quadratic_form_gap,
        thm:finite_range_cyclic_knabe, lem:quadratic_form_to_norm_bound}
    The norm gap (\ref{eq:open_knabe_block_norm_gap}) and positivity give
    $(H^{\mathrm{open}}_{W,R})^2\geq\gamma H^{\mathrm{open}}_{W,R}$ as a
    quadratic form.  Lemma~\ref{lem:right_spectator_quadratic_form_gap} gives
    the same inequality after adjoining the $N-W$ complementary sites, and
    Theorem~\ref{thm:cyclic_knabe_block_open_parent_hamiltonian} gives it for
    every cyclic sum $A_s$.  The oriented-separation theorem verifies the
    commutation hypothesis in
    Theorem~\ref{thm:finite_range_cyclic_knabe}.  That theorem gives
    \begin{align}
        \delta H_{N,R}^{\mathrm{ES}}(A)
         & \leq\left(H_{N,R}^{\mathrm{ES}}(A)\right)^2
    \end{align}
    as a quadratic form, using the total-sum identity
    (\ref{eq:cyclic_parent_local_sum}).  The inequality
    $(R-1)^2<m\gamma$ makes $\delta$ positive, and the finite-dimensional
    spectral criterion yields
    (\ref{eq:parent_hamiltonian_finite_range_knabe_gap}).
\end{proof}

\begin{theorem}[Uniform finite-range gap for primitive MPS parent Hamiltonians]
    \label{thm:primitive_parent_hamiltonian_finite_range_knabe_gap}
    \lean{MPSTensor.IsPrimitiveMPS.exists_parentHamiltonianES_gap_of_finiteRangeKnabe}
    \leanok
    \uses{thm:primitive_open_parent_hamiltonian_gap,
        thm:open_parent_hamiltonian_kernel_ground_space,
        thm:parent_hamiltonian_gap_from_open_knabe_block}
    Let $A$ be a primitive MPS tensor with nonzero bond dimension and a
    positive-definite fixed point.  There exist $l>1$, $\epsilon_l\geq0$, an
    integer $m\geq l+1$, and $\delta>0$ such that $A$ is $l$-block injective
    and
    \begin{align}
        \epsilon_l
         & <\frac{1}{\sqrt{l+1}}.
        \label{eq:primitive_finite_range_knabe_threshold}
    \end{align}
    Set
    \begin{align}
        R & =l+1,                                   \\
        W & =m+R-1=m+l,                             \\
        \gamma_{\mathrm{open}}
          & =\left(1-\epsilon_l\sqrt{l+1}\right)^2.
        \label{eq:primitive_finite_range_knabe_parameters}
    \end{align}
    The integer $m$ may be chosen so that
    \begin{align}
        l^2 & <m\gamma_{\mathrm{open}},
        \label{eq:primitive_finite_range_knabe_numerator}
    \end{align}
    and
    \begin{align}
        \delta
         & =\frac{m\gamma_{\mathrm{open}}-l^2}{m-l}.
        \label{eq:primitive_finite_range_knabe_delta}
    \end{align}
    For every $N\geq2m$ and every
    $v\in(\ker H_{N,l+1}^{\mathrm{ES}}(A))^\perp$,
    \begin{align}
        \delta\lVert v\rVert
         & \leq\left\lVert H_{N,l+1}^{\mathrm{ES}}(A)v\right\rVert.
        \label{eq:primitive_finite_range_knabe_gap}
    \end{align}
    The existence of a uniform MPS parent-Hamiltonian gap and the
    nearest-neighbor Knabe principle are stated in
    \cite[lines~1460--1500]{PerezGarcia2007Matrix} and
    \cite[lines~2184--2188,2194--2197]{Cirac2021Matrix}.  Those passages do
    not state the arbitrary finite-range coefficient
    (\ref{eq:primitive_finite_range_knabe_delta}).
\end{theorem}

\begin{proof}
    \leanok
    \uses{thm:primitive_open_parent_hamiltonian_gap,
        thm:open_parent_hamiltonian_kernel_ground_space,
        thm:parent_hamiltonian_gap_from_open_knabe_block}
    Theorem~\ref{thm:primitive_open_parent_hamiltonian_gap} gives $l$,
    $\epsilon_l$, and the norm gap $\gamma_{\mathrm{open}}$ on every open
    volume at least $l+1$.  Since $\gamma_{\mathrm{open}}>0$, choose
    $m\geq l+1$ with $l^2<m\gamma_{\mathrm{open}}$.  At the active volume
    $W=m+l$, block injectivity identifies the open ground space with the kernel
    of the open Hamiltonian, so
    (\ref{eq:primitive_open_parent_hamiltonian_gap}) is exactly the hypothesis
    (\ref{eq:open_knabe_block_norm_gap}) with $R=l+1$.  Apply
    Theorem~\ref{thm:parent_hamiltonian_gap_from_open_knabe_block}.  Substituting
    $R-1=l$ and $m-R+1=m-l$ gives
    (\ref{eq:primitive_finite_range_knabe_delta}) and
    (\ref{eq:primitive_finite_range_knabe_gap}).
\end{proof}

\section{Overlap with a sum of subspaces}

\begin{theorem}[Lower Gram estimate]
    \label{thm:subspace_overlap_lower_gram}
    \lean{Submodule.norm_sum_sq_ge_of_overlapMatrix}
    \leanok
    Let $v_1,\ldots,v_m$ lie in a complex inner product space, and let
    $B\in M_m(\R)$ have zero diagonal. Suppose that
    $|\langle v_i,v_j\rangle|\leq B_{ij}\lVert v_i\rVert\lVert v_j\rVert$
    whenever $i\neq j$. With the Euclidean operator norm on $B$,
    \begin{align}
        (1-\lVert B\rVert)\sum_i\lVert v_i\rVert^2
         & \leq\left\lVert\sum_i v_i\right\rVert^2.
    \end{align}
    This is the lower Gram estimate in the proof of
    \cite[Lemma~9]{Nachtergaele1996Spectral}.
\end{theorem}

\begin{proof}
    \leanok
    Put $w_i=\lVert v_i\rVert$. Expanding the squared norm gives
    \begin{align}
        \sum_i\lVert v_i\rVert^2
        -\left\lVert\sum_i v_i\right\rVert^2
         & \leq \sum_{i,j}B_{ij}w_iw_j
        =\langle w,Bw\rangle
        \leq\lVert B\rVert\lVert w\rVert^2.
    \end{align}
    Rearranging proves the assertion.
\end{proof}

\begin{theorem}[Overlap with a sum of subspaces]
    \label{thm:subspace_overlap_sum}
    \lean{Submodule.norm_inner_le_iSup_of_overlapMatrix}
    \leanok
    \uses{thm:subspace_overlap_lower_gram}
    Let $U,V_1,\ldots,V_m$ be subspaces of a complex inner product space.
    Let $a\in\R^m$ and $B\in M_m(\R)$, where $B$ has zero
    diagonal and $\lVert B\rVert<1$. Suppose that
    $|\langle u,v\rangle|\leq B_{ij}\lVert u\rVert\lVert v\rVert$
    for $i\neq j$, $u\in V_i$, and $v\in V_j$. Suppose also that
    $|\langle x,u\rangle|\leq a_i\lVert x\rVert\lVert u\rVert$
    for every $i$, $x\in U$, and $u\in V_i$.
    Then every $x\in U$ and $y\in V_1+\cdots+V_m$ satisfy
    \begin{align}
        |\langle x,y\rangle|
         & \leq\frac{\lVert a\rVert}{\sqrt{1-\lVert B\rVert}}
        \lVert x\rVert\lVert y\rVert.
    \end{align}
    This is \cite[Lemma~9]{Nachtergaele1996Spectral}, expressed as a bound
    on inner products.
\end{theorem}

\begin{proof}
    \leanok
    \uses{thm:subspace_overlap_lower_gram}
    Write $y=\sum_i v_i$ with $v_i\in V_i$. The triangle inequality and
    Cauchy--Schwarz give
    \begin{align}
        |\langle x,y\rangle|
         & \leq\lVert x\rVert\sum_i a_i\lVert v_i\rVert
        \leq\lVert x\rVert\lVert a\rVert
        \left(\sum_i\lVert v_i\rVert^2\right)^{1/2}.
    \end{align}
    Apply Theorem~\ref{thm:subspace_overlap_lower_gram} to the final factor.
\end{proof}

\begin{theorem}[Uniform overlap bound]
    \label{thm:subspace_overlap_sum_uniform}
    \lean{Submodule.iSup_overlap_bound_of_uniform}
    \leanok
    Let $U,V_1,\ldots,V_m$ be subspaces of a complex inner product space,
    and let $\varepsilon\geq0$ satisfy $\varepsilon m\leq1$.
    Suppose that every pair of vectors in distinct $V_i$ satisfies
    $|\langle u,v\rangle|\leq\varepsilon\lVert u\rVert\lVert v\rVert$,
    and that the same inequality holds for vectors in $U$ and $V_i$, for
    every $i$. Set
    \begin{align}
        \eta & =\frac{\varepsilon\sqrt m}
        {\sqrt{1-\varepsilon(m-1)}}.
    \end{align}
    Then $\eta\leq1$, and every $x\in U$ and $y\in V_1+\cdots+V_m$ satisfy
    $|\langle x,y\rangle|\leq\eta\lVert x\rVert\lVert y\rVert$.
    This is the uniform bound in \cite[Lemma~9]{Nachtergaele1996Spectral}.
\end{theorem}

\begin{proof}
    \leanok
    \uses{thm:subspace_overlap_sum}
    For $y=\sum_i v_i$, the off-diagonal terms satisfy
    \begin{align}
        \sum_{i\neq j}\varepsilon\lVert v_i\rVert\lVert v_j\rVert
         & =\varepsilon\left[\left(\sum_i\lVert v_i\rVert\right)^2
            -\sum_i\lVert v_i\rVert^2\right]
        \leq\varepsilon(m-1)\sum_i\lVert v_i\rVert^2.
    \end{align}
    Take $a_i=\varepsilon$ and $B_{ij}=\varepsilon(1-\delta_{ij})$. For
    $w\in\R^m$ write $|w|$ for the vector with entries $|w_i|$. Since $B$ is
    symmetric with nonnegative entries, the same computation applied to $|w|$
    gives
    \begin{align}
        |\langle w,Bw\rangle|
         & \leq\langle|w|,B|w|\rangle
        =\varepsilon\left[\left(\sum_i|w_i|\right)^2-\sum_i w_i^2\right]
        \leq\varepsilon(m-1)\lVert w\rVert^2,
    \end{align}
    and hence
    $\lVert B\rVert=\sup_{\lVert w\rVert=1}|\langle w,Bw\rangle|
        \leq\varepsilon(m-1)\leq1-\varepsilon$, which is less
    than one when $\varepsilon>0$ and vanishes when $\varepsilon=0$. With
    $\lVert a\rVert=\varepsilon\sqrt m$,
    Theorem~\ref{thm:subspace_overlap_sum} gives
    \begin{align}
        |\langle x,y\rangle|
         & \leq\frac{\varepsilon\sqrt m}{\sqrt{1-\varepsilon(m-1)}}
        \lVert x\rVert\lVert y\rVert.
    \end{align}
    Finally,
    $\varepsilon^2m\leq1-\varepsilon(m-1)$ follows from
    $\varepsilon m\leq1$, and gives $\eta\leq1$.
\end{proof}

\section{Comparison of sums of orthogonal projectors}

\begin{theorem}[Off-diagonal pairing estimate]
    \label{thm:ph_off_diagonal_pairing}
    \lean{Submodule.sum_offDiagonal_norm_inner_le_of_overlapMatrix}
    \leanok
    Let $(v_i)$ and $(w_i)$ be finite families in a complex inner product space.
    Let $B$ be a real matrix with zero diagonal and $\beta=\lVert B\rVert<1$.
    Suppose that, for $i\ne j$, each of
    $|\langle v_i,v_j\rangle|$, $|\langle w_i,w_j\rangle|$, and
    $|\langle v_i,w_j\rangle|$ is bounded by $B_{ij}$ times the product
    of the corresponding vector norms. Put $\delta=\beta/(1-\beta)$. Then
    \begin{align}
        \sum_{i\ne j}|\langle v_i,w_j\rangle|
         & \leq\delta\left\lVert\sum_i v_i\right\rVert
        \left\lVert\sum_i w_i\right\rVert.
    \end{align}
    This is the pairing estimate in the proof of
    \cite[Lemma~6.1]{Nachtergaele1996Spectral}.
\end{theorem}
\begin{proof}
    \leanok
    \uses{thm:subspace_overlap_lower_gram}
    Write $a_i=\lVert v_i\rVert$ and $b_i=\lVert w_i\rVert$.
    The left-hand side is at most $\langle a,Bb\rangle\leq\beta\lVert a\rVert\lVert b\rVert$.
    The lower Gram estimate gives
    $\lVert a\rVert\leq\lVert\sum_i v_i\rVert/\sqrt{1-\beta}$,
    and likewise for $b$.
\end{proof}

\begin{theorem}[Sum of block projectors]
    \label{thm:ph_sum_block_projectors}
    \lean{Submodule.norm_sum_starProjection_sub_iSup_le, Submodule.norm_sum_starProjection_le_of_overlapMatrix}
    \leanok
    Let $(V_i)$ be a finite family of subspaces of a finite-dimensional
    complex inner product space. Let $B$ have zero diagonal and
    $\beta=\lVert B\rVert<1/2$, and suppose that
    $|\langle u,v\rangle|\leq B_{ij}\lVert u\rVert\lVert v\rVert$
    for $i\ne j$, $u\in V_i$, and $v\in V_j$.
    Write $P_i$ for the orthogonal projector onto $V_i$, $P$ for the
    projector onto $\sum_i V_i$, $X=\sum_i P_i$, and
    $\delta=\beta/(1-\beta)$. Then
    \begin{align}
        \lVert X-P\rVert & \leq\frac{\delta}{1-\delta}, &
        \lVert X\rVert   & \leq\frac{1}{1-\delta}.
    \end{align}
    This is the finite-dimensional projector estimate underlying
    \cite[Lemma~6.1(i)]{Nachtergaele1996Spectral}.
\end{theorem}
\begin{proof}
    \leanok
    \uses{thm:ph_off_diagonal_pairing}
    For $u=\sum_j v_j$ with $v_j\in V_j$, self-adjointness gives
    \begin{align}
        \langle z,Xu-u\rangle & =\sum_{i\ne j}\langle P_i z,v_j\rangle.
    \end{align}
    The pairing estimate bounds its modulus by
    $\delta\lVert Xz\rVert\lVert u\rVert$.
    Since $XP=X$, it follows that $\lVert X-P\rVert\leq\delta\lVert X\rVert$.
    Combine this with $\lVert X\rVert\leq\lVert X-P\rVert+1$ and $\delta<1$.
\end{proof}

\begin{theorem}[Off-diagonal products of projectors]
    \label{thm:ph_off_diagonal_projectors}
    \lean{Submodule.norm_sum_starProjection_comp_sub_sum_le}
    \leanok
    Let $(U_i)$ and $(V_i)$ be finite families of subspaces of a
    finite-dimensional complex inner product space. Let $B$ have zero
    diagonal and $\beta=\lVert B\rVert<1$. Suppose that, for $i\ne j$,
    $|\langle u,v\rangle|\leq B_{ij}\lVert u\rVert\lVert v\rVert$
    whenever $(u,v)$ lies in $U_i\times U_j$, $V_i\times V_j$, or
    $U_i\times V_j$.
    Write $P_i,Q_i$ for their orthogonal projectors and put
    $X_U=\sum_i P_i$, $X_V=\sum_i Q_i$, $\delta=\beta/(1-\beta)$.
    Then
    \begin{align}
        \left\lVert X_U X_V-\sum_i P_iQ_i\right\rVert
         & \leq\delta\lVert X_U\rVert\lVert X_V\rVert.
    \end{align}
\end{theorem}
\begin{proof}
    \leanok
    \uses{thm:ph_off_diagonal_pairing}
    Pair the operator on the left with unit vectors $z,x$.
    Its matrix coefficient is $\sum_{i\ne j}\langle P_i z,Q_j x\rangle$.
    Apply the off-diagonal pairing estimate and take the supremum over $z,x$.
\end{proof}

\begin{theorem}[Comparison of joint projectors]
    \label{thm:ph_joint_projector_comparison}
    \lean{Submodule.norm_iSup_starProjection_sub_comp_le_of_overlapMatrix}
    \leanok
    Under the hypotheses of Theorem~\ref{thm:ph_off_diagonal_projectors},
    assume $\beta<1/2$. Let $W_i\subseteq U_i$, and let $R_i$ be the
    orthogonal projector onto $W_i$.
    Let $P,Q,R$ project onto $\sum_i U_i$, $\sum_i V_i$, and $\sum_i W_i$,
    respectively. Then
    \begin{align}
        \lVert R-PQ\rVert & \leq\frac{4\delta}{(1-\delta)^2}
        +\sum_i\lVert R_i-P_iQ_i\rVert.
    \end{align}
    This is the abstract projector comparison used in
    \cite[Lemma~6.1(ii)]{Nachtergaele1996Spectral}.
\end{theorem}
\begin{proof}
    \leanok
    \uses{thm:ph_sum_block_projectors, thm:ph_off_diagonal_projectors}
    Put $X_W=\sum_i R_i$, $Z=\sum_i P_iQ_i$,
    $a=\delta/(1-\delta)$, and $q=1/(1-\delta)$.
    The triangle inequality, applied through $X_W$, $Z$, and $X_U X_V$, gives
    \begin{align}
        \lVert R-PQ\rVert
         & \leq a+\sum_i\lVert R_i-P_iQ_i\rVert+aq+(aq+a) \\
         & \leq4aq+\sum_i\lVert R_i-P_iQ_i\rVert.
    \end{align}
    Here $\lVert X_U X_V-PQ\rVert\leq aq+a$ follows by expanding
    $(X_U-P)X_V+P(X_V-Q)$, and $a\leq aq$ since $q\geq1$.
\end{proof}

\section{Overlap bounds with spectator sites}

\begin{theorem}[Overlap from restrictions to a common interval]
    \label{thm:ph_middle_interval_overlap}
    \lean{MPSTensor.norm_inner_le_of_middle_groundSpaceES}
    \leanok
    Let $A,B$ be tensors with the same physical dimension and arbitrary bond
    dimensions. Suppose $\varepsilon\geq0$ and
    $|\langle a,b\rangle|\leq\varepsilon\lVert a\rVert\lVert b\rVert$
    for every $a\in G_L(A)$ and $b\in G_L(B)$.
    Let $x,y$ be vectors on three successive intervals of lengths $K,L,Q$.
    For each prefix configuration $u$ and suffix configuration $v$, define
    $x_{u,v}(w)=x(u,w,v)$ and $y_{u,v}(w)=y(u,w,v)$.
    If $x_{u,v}\in G_L(A)$ and $y_{u,v}\in G_L(B)$ for all $u,v$, then
    \begin{align}
        |\langle x,y\rangle| & \leq\varepsilon\lVert x\rVert\lVert y\rVert.
    \end{align}
    No positivity assumption on $K,L,Q$ is required.
    This is the finite-sum estimate used to extend the overlap bounds in
    \cite[Lemma~6.1(ii)]{Nachtergaele1996Spectral}.
\end{theorem}
\begin{proof}
    \leanok
    Decomposing along the prefix and suffix configurations gives
    \begin{align}
        \langle x,y\rangle & =\sum_{u,v}\langle x_{u,v},y_{u,v}\rangle,
        \notag                                                          \\
        \lVert x\rVert^2   & =\sum_{u,v}\lVert x_{u,v}\rVert^2,
        \qquad
        \lVert y\rVert^2=\sum_{u,v}\lVert y_{u,v}\rVert^2.
        \notag
    \end{align}
    The triangle inequality, the overlap bound, and Cauchy--Schwarz for the
    sequences $(\lVert x_{u,v}\rVert)$ and $(\lVert y_{u,v}\rVert)$ give
    \begin{align}
        |\langle x,y\rangle|
         & \leq\sum_{u,v}|\langle x_{u,v},y_{u,v}\rangle|                    \\
         & \leq\varepsilon\sum_{u,v}\lVert x_{u,v}\rVert\lVert y_{u,v}\rVert
        \leq\varepsilon\lVert x\rVert\lVert y\rVert.
    \end{align}
\end{proof}

\begin{theorem}[Overlap of extended boundary ranges]
    \label{thm:ph_extended_boundary_overlap}
    \lean{MPSTensor.norm_inner_leftBoundaryMapES_le}
    \lean{MPSTensor.norm_inner_reassocTailBoundaryMapES_le}
    \lean{MPSTensor.norm_inner_left_reassocTailBoundaryMapES_le}
    \leanok
    \uses{def:spectator_boundary_maps_es, def:reassociated_tail_boundary_map}
    Assume the ground-space overlap bound of
    Theorem~\ref{thm:ph_middle_interval_overlap}.
    Let $U_A$ be the range of $\Gamma^{\mathrm{left}}_{K+L,Q}(A)$
    and $V_A$ the range of the reassociated tail boundary map
    $\widetilde\Gamma^{\mathrm{tail}}_{K;L,Q}(A)$,
    both in the physical space on $K+L+Q$ sites. Define $U_B,V_B$ in the same way.
    Then $|\langle x,y\rangle|\leq\varepsilon\lVert x\rVert\lVert y\rVert$
    for vectors in any of the pairs of ranges
    $U_A\times U_B$, $V_A\times V_B$, and $U_A\times V_B$.
\end{theorem}
\begin{proof}
    \leanok
    \uses{thm:ph_middle_interval_overlap, def:reassociated_tail_boundary_map,
        thm:reassociated_tail_boundary_three_block, thm:left_boundary_map_append,
        lem:ground_space_in_tail_ground, lem:prefix_restrict_ground_space_map}
    For a left boundary vector with boundary family $(Z_v)$, fixing $u,v$
    leaves the middle vector $\Gamma_L(Z_vA^u)$.
    For a right boundary vector with boundary family $(Y_u)$, the corresponding
    middle vector is $\Gamma_L(A^vY_u)$.
    Both belong to the appropriate length-$L$ ground space, so the
    common-interval estimate applies to each of the three pairs.
\end{proof}

\begin{theorem}[Joint projectors on overlapping MPS intervals]
    \label{thm:ph_block_interval_projectors}
    \lean{MPSTensor.norm_iSup_groundSpaceES_sub_comp_le_of_overlapMatrix}
    \leanok
    Let $(A_i)$ be a finite family of tensors with the same physical dimension,
    with bond dimensions allowed to depend on $i$. Let $B$ be a nonnegative
    real matrix with zero diagonal and $\beta=\lVert B\rVert<1/2$.
    Suppose that for $i\ne j$, $x\in G_L(A_i)$, and $y\in G_L(A_j)$,
    $|\langle x,y\rangle|\leq B_{ij}\lVert x\rVert\lVert y\rVert$.
    In the physical space on $K+L+Q$ sites, put
    \begin{align}
        U_i & =G_{K+L}(A_i)\otimes\mathcal H_Q,  &
        V_i & =\mathcal H_K\otimes G_{L+Q}(A_i), &
        W_i & =G_{K+L+Q}(A_i).
    \end{align}
    Let $P_i,Q_i,R_i$ be the orthogonal projectors onto these spaces, and let
    $P,Q,R$ project onto their respective sums. With
    $\delta=\beta/(1-\beta)$,
    \begin{align}
        \lVert R-PQ\rVert & \leq\frac{4\delta}{(1-\delta)^2}
        +\sum_i\lVert R_i-P_iQ_i\rVert.
    \end{align}
    This is the physical interval comparison in the proof of
    \cite[Lemma~6.1(ii)]{Nachtergaele1996Spectral}, before inserting the
    individual block error bounds.
\end{theorem}
\begin{proof}
    \leanok
    \uses{thm:ph_extended_boundary_overlap, thm:ph_joint_projector_comparison,
        thm:spectator_boundary_hilbert_factorization, thm:range_ground_space_map_es}
    The boundary-map factorization gives $W_i\subseteq U_i$.
    The extended-boundary overlap estimate supplies all three bounds for
    $U_i\times U_j$, $V_i\times V_j$, and $U_i\times V_j$ with the same
    matrix $B$. Apply the joint-projector comparison.
\end{proof}

\begin{theorem}[Uniform decay of the joint projector error]
    \label{thm:ph_block_interval_defect_decay}
    \lean{MPSTensor.eventually_iSup_groundSpaceES_projector_defect_le}
    \leanok
    \uses{thm:ph_block_interval_projectors}
    Let $(A_i)_{i\in I}$ be a finite family of tensors, and use the interval
    spaces and projectors of Theorem~\ref{thm:ph_block_interval_projectors}.
    Suppose that, for each $\varepsilon>0$ and all sufficiently large $L$,
    the following two bounds hold. For distinct $i,j$, $x\in G_L(A_i)$,
    and $y\in G_L(A_j)$,
    $|\langle x,y\rangle|\leq\varepsilon\lVert x\rVert\lVert y\rVert$.
    For every $i$, $K\geq0$, and $Q>0$,
    $\lVert Q_iP_i-R_i\rVert\leq\varepsilon$.
    Then, for every $\eta>0$, all sufficiently large $L$, and every
    $K\geq0$ and $Q>0$,
    \begin{align}
        \lVert R-PQ\rVert & \leq\eta.
    \end{align}
\end{theorem}
\begin{proof}
    \leanok
    \uses{thm:ph_block_interval_projectors}
    Put $r=|I|$ and let $C$ have zero diagonal and all off-diagonal entries
    equal to one. As $\varepsilon\downarrow0$, the quantities
    $\beta=\lVert\varepsilon C\rVert$, $\delta=\beta/(1-\beta)$, and
    $4\delta/(1-\delta)^2$ tend to zero. Choose $\varepsilon>0$ so that
    $\beta<1/2$ and the last quantity is less than $\eta/(r+1)$.
    For all sufficiently large $L$, the overlap bound holds with
    $B=\varepsilon C$, and every individual projector error is at most
    $\eta/(r+1)$. Since $P_i$, $Q_i$, and $R_i$ are self-adjoint,
    $(R_i-P_iQ_i)^*=R_i-Q_iP_i$, and hence
    \begin{align}
        \lVert R_i-P_iQ_i\rVert & =\lVert Q_iP_i-R_i\rVert.
    \end{align}
    The interval comparison therefore gives
    \begin{align}
        \lVert R-PQ\rVert
         & \leq\frac{\eta}{r+1}+r\frac{\eta}{r+1}=\eta.
    \end{align}
\end{proof}

\begin{theorem}[Uniform projector error decay for primitive blocks]
    \label{thm:ph_primitive_block_interval_defect_decay}
    \lean{MPSTensor.eventually_iSup_groundSpaceES_projector_defect_le_of_inequivalent}
    \leanok
    \uses{thm:ph_block_interval_projectors, def:primitive_mps}
    Let $(A_i)_{i\in I}$ be a finite family of normalized primitive tensors
    with positive bond dimensions and positive-definite fixed points.
    Suppose that, if the bond dimensions of two distinct members agree, they
    are not related by a gauge transformation and a phase. Use the interval projectors of
    Theorem~\ref{thm:ph_block_interval_projectors}.
    For every $\eta>0$, all sufficiently large $L$, and every $K\geq0$
    and $Q>0$,
    \begin{align}
        \lVert R-PQ\rVert & \leq\eta.
    \end{align}
    This is the uniform decay consequence of
    \cite[Lemma~6.1(ii)]{Nachtergaele1996Spectral} for a positive right
    outer interval.
\end{theorem}
\begin{proof}
    \leanok
    \uses{thm:ph_block_interval_defect_decay,
        thm:inequivalent_block_ground_space_overlap_decay,
        thm:primitive_whole_increment_defect_decay}
    Inequivalence gives decay of the overlaps between distinct blocks.
    Primitivity gives uniform decay of each individual projector error.
    Since the family is finite, the thresholds can be chosen simultaneously
    for all blocks and all pairs. Apply the joint projector error estimate.
\end{proof}

\begin{theorem}[Interval spaces of a block-diagonal tensor]
    \label{thm:ph_block_sum_interval_spaces}
    \lean{MPSTensor.range_leftBoundaryMapES_toTensorFromBlocks_eq_iSup,
        MPSTensor.range_tailBoundaryMapES_toTensorFromBlocks_eq_iSup,
        MPSTensor.range_reassocTailBoundaryMapES_toTensorFromBlocks_eq_iSup}
    \leanok
    \uses{def:whole_increment_physical_reassociation}
    Let $B=\bigoplus_{j=1}^r\mu_jA_j$ with $\mu_j\ne0$ for every $j$.
    For every $K,L\geq0$, the extended interval spaces satisfy
    \begin{align}
        G_K(B)\otimes\mathcal H_L
         & =\sum_{j=1}^r\bigl(G_K(A_j)\otimes\mathcal H_L\bigr),  \\
        \mathcal H_K\otimes G_L(B)
         & =\sum_{j=1}^r\bigl(\mathcal H_K\otimes G_L(A_j)\bigr).
    \end{align}
    For every $K,L,Q\geq0$, with the reassociation $\mathscr A_{K,L,Q}$ of
    Definition~\ref{def:whole_increment_physical_reassociation}, the
    right interval space of length $L+Q$ satisfies
    \begin{align}
        \mathscr A_{K,L,Q}\bigl(\mathcal H_K\otimes G_{L+Q}(B)\bigr)
         & =\sum_{j=1}^r
        \mathscr A_{K,L,Q}\bigl(\mathcal H_K\otimes G_{L+Q}(A_j)\bigr).
    \end{align}
\end{theorem}
\begin{proof}
    \leanok
    \uses{thm:block_sum_local_ground_space,
        def:whole_increment_physical_reassociation}
    Let $\Omega_L$ be the set of physical words of length $L$. For subspaces
    $S_1,\dots,S_r$ of a vector space, the finite coordinate identity
    \begin{align}
        \prod_{\tau\in\Omega_L}\Bigl(\sum_{j=1}^rS_j\Bigr)
         & =\sum_{j=1}^r\prod_{\tau\in\Omega_L}S_j
    \end{align}
    holds: the right side is contained in the left, and a family
    $(\psi_\tau)_\tau$ with $\psi_\tau=\sum_j\psi_{\tau,j}$ and
    $\psi_{\tau,j}\in S_j$ is the sum over $j$ of the families
    $(\psi_{\tau,j})_\tau$.
    Let $T_{K,L}$ be the coordinate isomorphism sending
    $(\psi_\tau)_{\tau\in\Omega_L}$ to the vector with coefficient
    $\psi_\tau(\sigma)$ at $(\sigma,\tau)\in\Omega_K\times\Omega_L$. Then
    $G_K(C)\otimes\mathcal H_L=T_{K,L}\prod_{\tau\in\Omega_L}G_K(C)$ for every
    tensor $C$. The local identity $G_K(B)=\sum_jG_K(A_j)$ and the finite
    coordinate identity give
    \begin{align}
        G_K(B)\otimes\mathcal H_L
         & =T_{K,L}\prod_{\tau\in\Omega_L}\Bigl(\sum_{j=1}^rG_K(A_j)\Bigr)
        =\sum_{j=1}^rT_{K,L}\prod_{\tau\in\Omega_L}G_K(A_j)
        =\sum_{j=1}^r\bigl(G_K(A_j)\otimes\mathcal H_L\bigr),
    \end{align}
    since linear maps preserve sums of subspaces. The same argument, with
    the free sites on the left and coordinates indexed by $\Omega_K$, gives
    the second identity. Applying the linear map $\mathscr A_{K,L,Q}$ to the
    second identity at lengths $K$ and $L+Q$ gives the third.
\end{proof}

\begin{theorem}[Projector error decay for a block-diagonal tensor]
    \label{thm:ph_block_diagonal_projector_decay}
    \lean{MPSTensor.eventually_toTensorFromBlocks_projector_defect_le}
    \leanok
    Let $B=\bigoplus_{j=1}^r\mu_jA_j$, where every $\mu_j$ is nonzero
    and the $A_j$ are pairwise gauge-phase inequivalent normalized primitive
    tensors with positive bond dimensions and positive definite invariant
    densities. Write $P_I(B)$ for the orthogonal projector onto the ground
    space on the interval $I$, extended by the identity on the remaining
    sites. For every $\eta>0$, all sufficiently large $L$, and every
    $K\geq0$ and $Q>0$,
    \begin{align}
        \bigl\lVert P_{[1,K+L+Q]}(B)
        -P_{[1,K+L]}(B)P_{[K+1,K+L+Q]}(B)\bigr\rVert
         & \leq\eta.
    \end{align}
\end{theorem}
\begin{proof}
    \leanok
    \uses{thm:ph_primitive_block_interval_defect_decay,
        thm:ph_block_interval_projectors,
        thm:block_sum_local_ground_space, thm:ph_block_sum_interval_spaces}
    For each block put
    \begin{align}
        U_j & =G_{K+L}(A_j)\otimes\mathcal H_Q,                                &
        V_j & =\mathscr A_{K,L,Q}\bigl(\mathcal H_K\otimes G_{L+Q}(A_j)\bigr), &
        W_j & =G_{K+L+Q}(A_j),
    \end{align}
    and let $P,S,R$ be the orthogonal projectors onto $\sum_jU_j$,
    $\sum_jV_j$, and $\sum_jW_j$. These are the projectors written $P,Q,R$
    in Theorem~\ref{thm:ph_block_interval_projectors}; the middle one is
    renamed here because $Q$ denotes a length. The local identity at
    length $K+L+Q$, the first identity of
    Theorem~\ref{thm:ph_block_sum_interval_spaces} at lengths $K+L$ and $Q$,
    and its reassociated identity at lengths $K,L,Q$ give
    \begin{align}
        P_{[1,K+L+Q]}(B)   & =P_{G_{K+L+Q}(B)}
        =P_{\sum_jW_j}=R,                                               \\
        P_{[1,K+L]}(B)     & =P_{G_{K+L}(B)\otimes\mathcal H_Q}
        =P_{\sum_jU_j}=P,                                               \\
        P_{[K+1,K+L+Q]}(B) & =P_{\mathscr A_{K,L,Q}(\mathcal H_K\otimes
                G_{L+Q}(B))}=P_{\sum_jV_j}=S.
    \end{align}
    Hence the left side of the asserted bound is $\lVert R-PS\rVert$, and
    Theorem~\ref{thm:ph_primitive_block_interval_defect_decay} with the
    prescribed $\eta$ bounds it by $\eta$ for all sufficiently large $L$,
    every $K\geq0$, and every $Q>0$.
\end{proof}

\section{Overlaps of distinct block ground spaces}

\begin{theorem}[Mixed boundary Gram identity]
    \label{thm:mixed_boundary_gram_identity}
    \lean{MPSTensor.inner_single_mixed_groundSpaceGram_single,
        MPSTensor.adjoint_groundSpaceMapES_comp_eq_mixedTransfer}
    \leanok
    \uses{def:ground_space_map_es, def:mixed_transfer_rect}
    Let $A$ and $B$ have bond dimensions $D_A$ and $D_B$, respectively,
    and let $\Gamma_n^{\mathrm{ES}}(A)$ and $\Gamma_n^{\mathrm{ES}}(B)$ be
    their Hilbert-space boundary maps, writing $E_{ab}$ also for the
    coordinate vector of a matrix unit. Let $F_{BA}^{\mathrm{rect}}$ be the
    rectangular mixed transfer operator on $M_{D_B\times D_A}(\C)$. For
    matrix units $E_{ab}$ and $E_{ce}$,
    \begin{align}
        \langle\Gamma_n^{\mathrm{ES}}(A)E_{ab},\Gamma_n^{\mathrm{ES}}(B)E_{ce}\rangle
         & =\bigl[(F_{BA}^{\mathrm{rect}})^n(E_{ca})\bigr]_{eb}.
    \end{align}
    Thus $\Gamma_n^{\mathrm{ES}}(A)^*\Gamma_n^{\mathrm{ES}}(B)$ is the rectangular reshuffling of
    $(F_{BA}^{\mathrm{rect}})^n$. This is the boundary-coordinate form of
    equations~\textup{ip} and~\textup{P12} in the proof of the lemma
    \textup{disjoint} in \cite{Nachtergaele1996Spectral}.
\end{theorem}

\begin{proof}
    \leanok
    Expanding the trace boundary maps gives
    \begin{align}
        \langle\Gamma_n^{\mathrm{ES}}(A)E_{ab},\Gamma_n^{\mathrm{ES}}(B)E_{ce}\rangle
         & =\sum_{|w|=n}\overline{A^w_{ba}}B^w_{ec}
        =\sum_{|w|=n}\bigl[B^wE_{ca}(A^w)^*\bigr]_{eb}.
    \end{align}
    The word expansion of $(F_{BA}^{\mathrm{rect}})^n$ gives the assertion.
\end{proof}

\begin{theorem}[Mixed boundary Gram decay]
    \label{thm:mixed_boundary_gram_decay}
    \lean{MPSTensor.adjoint_groundSpaceMapES_comp_tendsto_zero,
        MPSTensor.adjoint_groundSpaceMapES_comp_tendsto_zero_of_spectralRadius_lt_one}
    \leanok
    \uses{thm:mixed_boundary_gram_identity, def:ground_space_map_es, def:mixed_transfer_rect}
    If $(F_{BA}^{\mathrm{rect}})^n(X)\to0$ for every boundary matrix $X$, then
    $\lVert\Gamma_n^{\mathrm{ES}}(A)^*\Gamma_n^{\mathrm{ES}}(B)\rVert\to0$.
    In particular this holds when the spectral radius of
    $F_{BA}^{\mathrm{rect}}$ is less than one.
\end{theorem}

\begin{proof}
    \leanok
    \uses{thm:mixed_boundary_gram_identity}
    The preceding identity gives convergence of every entry of the mixed
    Gram operator. Its domain and codomain have fixed finite dimensions,
    so entrywise convergence implies convergence in operator norm.
    A spectral radius less than one gives decay of the powers of
    $F_{BA}^{\mathrm{rect}}$, as in equation~\textup{limP12} in the proof of the
    lemma \textup{disjoint} in \cite{Nachtergaele1996Spectral}.
\end{proof}

\begin{theorem}[Ground-space overlap from mixed Gram bounds]
    \label{thm:ground_space_overlap_mixed_gram}
    \lean{MPSTensor.norm_inner_groundSpaceES_le_mixedGram}
    \leanok
    \uses{def:ground_space_map_es}
    Suppose that $\Gamma_n^{\mathrm{ES}}(A)$ and $\Gamma_n^{\mathrm{ES}}(B)$ are injective, and put
    $\mathcal K_n(A)=\Gamma_n^{\mathrm{ES}}(A)^*\Gamma_n^{\mathrm{ES}}(A)$ and
    $\mathcal K_n(B)=\Gamma_n^{\mathrm{ES}}(B)^*\Gamma_n^{\mathrm{ES}}(B)$. For vectors
    $x\in\operatorname{ran}\Gamma_n^{\mathrm{ES}}(A)$ and
    $y\in\operatorname{ran}\Gamma_n^{\mathrm{ES}}(B)$,
    \begin{align}
        |\langle x,y\rangle|
         & \leq\sqrt{\lVert \mathcal K_n(A)^{-1}\rVert}\,
        \lVert\Gamma_n^{\mathrm{ES}}(A)^*\Gamma_n^{\mathrm{ES}}(B)\rVert\,
        \sqrt{\lVert \mathcal K_n(B)^{-1}\rVert}\,
        \lVert x\rVert\lVert y\rVert.
    \end{align}
\end{theorem}

\begin{proof}
    \leanok
    Write $x=\Gamma_n^{\mathrm{ES}}(A)u$ and $y=\Gamma_n^{\mathrm{ES}}(B)v$. The left inverse
    $\mathcal K_n(A)^{-1}\Gamma_n^{\mathrm{ES}}(A)^*$ has norm
    $\sqrt{\lVert \mathcal K_n(A)^{-1}\rVert}$, so
    $\lVert u\rVert\leq\sqrt{\lVert \mathcal K_n(A)^{-1}\rVert}\lVert x\rVert$;
    the analogous estimate holds for $v$. Apply these bounds to
    $|\langle u,\Gamma_n^{\mathrm{ES}}(A)^*\Gamma_n^{\mathrm{ES}}(B)v\rangle|$.
    This makes explicit the normalization step using
    equation~\textup{C1C2} in \cite{Nachtergaele1996Spectral}.
\end{proof}

\begin{theorem}[Uniform decay of block ground-space overlaps]
    \label{thm:block_ground_space_overlap_decay}
    \lean{MPSTensor.IsPrimitiveMPS.eventually_norm_inner_groundSpaceES_le_of_mixed_gap}
    \leanok
    \uses{thm:mixed_boundary_gram_decay, thm:ground_space_overlap_mixed_gram,
        thm:primitive_ground_space_gram_inverse_tendsto, def:mixed_transfer_rect}
    Let $A$ and $B$ be trace-preserving tensors with positive-definite
    fixed matrices $\rho$ and $\sigma$. Suppose their complementary
    transfer maps have spectral radius less than one, and that
    $\rho_{\spec}(F_{BA}^{\mathrm{rect}})<1$.
    For every $\varepsilon>0$, there is $n_0$ such that for every
    $n\geq n_0$, $x\in G_n(A)$, and $y\in G_n(B)$,
    \begin{align}
        |\langle x,y\rangle| & \leq\varepsilon\lVert x\rVert\lVert y\rVert.
    \end{align}
    This is the bound on inner products of ground-space vectors obtained
    from the spectral-radius estimate in the proof of the lemma \textup{disjoint}
    in \cite{Nachtergaele1996Spectral}.
\end{theorem}

\begin{proof}
    \leanok
    \uses{thm:mixed_boundary_gram_decay, thm:ground_space_overlap_mixed_gram,
        thm:primitive_ground_space_gram_inverse_tendsto}
    By Theorem~\ref{thm:primitive_ground_space_gram_inverse_tendsto}, the
    self-Gram operators $\mathcal K_n(A)$ and $\mathcal K_n(B)$ converge to invertible limits
    $\mathcal K_\infty(A)$ and $\mathcal K_\infty(B)$ determined by $\rho$ and $\sigma$, and
    their inverses converge to $\mathcal K_\infty(A)^{-1}$ and $\mathcal K_\infty(B)^{-1}$.
    Since the limits are invertible, so are $\mathcal K_n(A)$ and $\mathcal K_n(B)$, and hence
    $\Gamma_n^{\mathrm{ES}}(A)$ and $\Gamma_n^{\mathrm{ES}}(B)$ are injective, for all
    sufficiently large $n$. By Theorem~\ref{thm:mixed_boundary_gram_decay},
    the coefficient in Theorem~\ref{thm:ground_space_overlap_mixed_gram}
    satisfies
    \begin{align}
        c(n)
         & =\sqrt{\lVert \mathcal K_n(A)^{-1}\rVert}\,
        \lVert\Gamma_n^{\mathrm{ES}}(A)^*\Gamma_n^{\mathrm{ES}}(B)\rVert\,
        \sqrt{\lVert \mathcal K_n(B)^{-1}\rVert}
        \longrightarrow
        \sqrt{\lVert \mathcal K_\infty(A)^{-1}\rVert}\cdot0\cdot
        \sqrt{\lVert \mathcal K_\infty(B)^{-1}\rVert}=0.
    \end{align}
    Choose $n_0$ with $c(n)\leq\varepsilon$ for $n\geq n_0$; then
    Theorem~\ref{thm:ground_space_overlap_mixed_gram} gives the bound,
    uniformly in $x$ and $y$.
\end{proof}

\begin{theorem}[Asymptotic orthogonality of inequivalent blocks]
    \label{thm:inequivalent_block_ground_space_overlap_decay}
    \lean{MPSTensor.IsPrimitiveMPS.eventually_norm_inner_groundSpaceES_le_of_inequivalent}
    \leanok
    \uses{thm:block_ground_space_overlap_decay, thm:transfer_operator_gap_irr_tp,
        thm:transfer_operator_gap_rect_irr_tp}
    Let $A$ and $B$ be trace-preserving tensors with positive-definite
    fixed matrices and complementary transfer spectral radii less than one.
    Suppose that, if their bond dimensions agree, they are not gauge-phase
    equivalent. Then, for every $\varepsilon>0$ and all sufficiently large $n$,
    every $x\in G_n(A)$ and $y\in G_n(B)$ satisfy
    \begin{align}
        |\langle x,y\rangle| & \leq\varepsilon\lVert x\rVert\lVert y\rVert.
    \end{align}
    This is the normalized tensor form of equation~\textup{limitepsilonm}
    in \cite{Nachtergaele1996Spectral}.
\end{theorem}

\begin{proof}
    \leanok
    \uses{thm:block_ground_space_overlap_decay, thm:transfer_operator_gap_irr_tp,
        thm:transfer_operator_gap_rect_irr_tp}
    The positive-definite fixed matrices and complementary gaps imply
    irreducibility. For equal bond dimensions, inequivalence and
    Theorem~\ref{thm:transfer_operator_gap_irr_tp} give
    $\rho_{\spec}(F_{BA}^{\mathrm{rect}})=\rho_{\spec}(F_{BA})<1$.
    For unequal bond dimensions, a peripheral eigenvector of
    $F_{BA}^{\mathrm{rect}}$ would force equality of the bond dimensions, so
    Theorem~\ref{thm:transfer_operator_gap_rect_irr_tp} gives
    \begin{align}
        \rho_{\spec}(F_{BA}^{\mathrm{rect}}) & <1.
    \end{align}
    Apply Theorem~\ref{thm:block_ground_space_overlap_decay}.
\end{proof}

\section{The normalized primitive gauge of a normal tensor}

A normal tensor need not be trace-preserving or have spectral radius one.
Standard treatments normalize the tensor by scaling the spectral radius of the
transfer map to one and applying a gauge transformation by the square root of a
positive fixed point~\cite{Nachtergaele1996Spectral,PerezGarcia2007Matrix}. This
gauge reduction extends the finite-range spectral gap to arbitrary normal
tensors.

\begin{theorem}[Rescaling invariance of block injectivity and normality]
    \label{thm:normality_rescaling_invariance}
    \lean{MPSTensor.wordSpan_smul_eq, MPSTensor.isNBlkInjective_smul_iff,
        MPSTensor.isNormal_smul_iff}
    \leanok
    \uses{def:word_span, def:l_blk_injective, def:normal}
    Let $\zeta\neq0$ and let $A$ be a tensor.  For every $N$ the length-$N$
    word spans of $A$ and of $\zeta A$ agree, so $\zeta A$ is $N$-block
    injective exactly when $A$ is, and $\zeta A$ is normal exactly when $A$ is.
\end{theorem}

\begin{proof}
    \leanok
    \uses{lem:eval_word_smul}
    Every length-$N$ word of $\zeta A$ is the corresponding word of $A$
    multiplied by $\zeta^{N}$, and $\zeta^{N}$ is invertible, so each span
    contains the other.
\end{proof}

\begin{theorem}[A normal tensor has a nonzero matrix]
    \label{thm:normal_tensor_nonzero_matrix}
    \lean{MPSTensor.exists_apply_ne_zero_of_isNormal}
    \leanok
    \uses{def:normal}
    Let $D>0$ and let $A$ be a normal tensor.  Then $A^i\neq0$ for at least
    one $i$.
\end{theorem}

\begin{proof}
    \leanok
    \uses{def:l_blk_injective}
    If every $A^i$ vanished, then every word of positive length would vanish,
    so the span of the length-$N$ words would be $\{0\}$ for the length $N>0$
    at which $A$ is block injective.  That span is $\MN{D}$, which contains
    $\Id\neq0$ because $D>0$.
\end{proof}

\begin{lemma}[Trace pairing of a conjugated word against a boundary matrix]
    \label{lem:trace_gauge_boundary}
    \lean{MPSTensor.trace_gauge_boundary}
    \leanok
    Let $X$ be invertible and let $E$ and $Y$ be $D\times D$ matrices.  Then
    \begin{align}
        \tr\bigl(XEX^{-1}Y\bigr) & =\tr\bigl(E\,X^{-1}YX\bigr). \notag
    \end{align}
\end{lemma}

\begin{proof}
    \leanok
    Cyclicity of the trace moves the leading factor $X$ to the end of the
    product.
\end{proof}

\begin{theorem}[Gauge invariance of the local ground space]
    \label{thm:local_ground_space_gauge_invariance}
    \lean{MPSTensor.GaugeEquiv.groundSpace_le,
        MPSTensor.GaugeEquiv.groundSpace_eq}
    \leanok
    \uses{def:gauge_equiv, def:ground_space_parent}
    Let $A$ and $B$ be gauge equivalent.  For every $L$ one has
    $G_L(A)=G_L(B)$.
\end{theorem}

\begin{proof}
    \leanok
    \uses{lem:eval_word_gauge, lem:trace_gauge_boundary}
    If $B^i=XA^iX^{-1}$, then $B^\sigma=XA^\sigma X^{-1}$ for every length-$L$
    configuration $\sigma$ by Lemma~\ref{lem:eval_word_gauge}, so by
    Lemma~\ref{lem:trace_gauge_boundary} the boundary parametrization of $B$
    at $Y$ agrees with that of $A$ at $X^{-1}YX$.  This gives
    $G_L(B)\subseteq G_L(A)$, and the symmetric argument applied to
    $A^i=X^{-1}B^iX$ gives the reverse inclusion.
\end{proof}

\begin{theorem}[Rescaling invariance of the local ground space]
    \label{thm:local_ground_space_rescaling}
    \lean{MPSTensor.groundSpace_smul_eq}
    \leanok
    \uses{def:ground_space_parent}
    Let $\zeta\neq0$.  For every $L$ one has $G_L(\zeta A)=G_L(A)$.
\end{theorem}

\begin{proof}
    \leanok
    \uses{lem:eval_word_smul}
    The boundary parametrization of $\zeta A$ at length $L$ sends $X$ to the
    value of the boundary parametrization of $A$ at $\zeta^{L}X$, and
    multiplication by $\zeta^{L}$ is a bijection of $\MN{D}$.
\end{proof}

\begin{theorem}[Local ground spaces determine the parent Hamiltonian]
    \label{thm:parent_hamiltonian_from_ground_space}
    \lean{MPSTensor.parentInteraction_eq_of_groundSpace_eq,
        MPSTensor.localTerm_eq_of_parentInteraction_eq,
        MPSTensor.parentHamiltonian_eq_of_groundSpace_eq,
        MPSTensor.parentHamiltonianES_eq_of_groundSpace_eq}
    \leanok
    \uses{def:ground_space_parent, def:parent_interaction,
        def:parent_hamiltonian, def:parent_hamiltonian_es}
    Let $A$ and $B$ have the same bond dimension and satisfy $G_L(A)=G_L(B)$.
    Then their canonical parent interactions at length $L$ agree, hence so do
    all their translated local terms and their parent Hamiltonians
    $H_{N,L}(A)=H_{N,L}(B)$ and $H_{N,L}^{\mathrm{ES}}(A)
        =H_{N,L}^{\mathrm{ES}}(B)$ on every periodic chain.
\end{theorem}

\begin{proof}
    \leanok
    \uses{def:ground_space_es, def:parent_interaction, def:local_term_parent,
        def:parent_hamiltonian, def:parent_hamiltonian_es}
    The parent interaction at length $L$ is the orthogonal projector onto
    $G_L(A)^\perp$, so it depends on $G_L(A)$ alone.  Every translated local
    term is built from that single projector, and the parent Hamiltonian is
    their sum.
\end{proof}

\begin{theorem}[Normality alone gives an irreducible tensor]
    \label{thm:normal_irreducible_tensor}
    \lean{MPSTensor.isPrimitivePaper_of_isNormal,
        MPSTensor.isIrreducibleTensor_of_isPrimitivePaper}
    \leanok
    \uses{def:normal, def:paper_primitive, def:has_inv_proj,
        def:irreducible_tensor}
    Let $A$ be a normal tensor.  Then $A$ is primitive in the sense of
    Definition~\ref{def:paper_primitive}, and $A$ has no nontrivial invariant
    projection, so $A$ is irreducible in the sense of
    Definition~\ref{def:irreducible_tensor}.  No normalization is imposed on
    $A$.
\end{theorem}

\begin{proof}
    \leanok
    \uses{def:eventually_full_kraus_rank, rem:wolf_6_8_pairwise}
    Normality is eventual full Kraus rank
    (Remark~\ref{rem:wolf_6_8_pairwise}), so $S_N(A)=\MN{D}$ for some $N>0$.
    Applying the length-$N$ words to a nonzero $\varphi$ then spans $\C^D$,
    which is the spreading condition.  A subspace invariant under all matrices
    of $A$ and containing a nonzero vector therefore contains the span of the
    images of that vector, namely all of $\C^D$; so no invariant projection is
    nontrivial.
\end{proof}

\begin{theorem}[Normalized primitive gauge of a normal tensor]
    \label{thm:normalized_primitive_gauge}
    \lean{MPSTensor.exists_isPrimitiveMPS_gauge_of_isNormal}
    \leanok
    \uses{def:normal, def:primitive_mps, def:gauge_equiv, def:mpv,
        def:ground_space_parent, def:parent_interaction,
        def:chain_ground_space}
    Let $D>0$ and let $A$ be a normal tensor.  There exist a scalar
    $\zeta\neq0$, a tensor $B$, and a matrix $\rho>0$ such that $\zeta A$ and
    $B$ are gauge equivalent, the pair $(B,\rho)$ is primitive, and
    \begin{align}
        V^{(N)}(B)_{\sigma}
         & =\zeta^{N}V^{(N)}(A)_{\sigma} \notag
    \end{align}
    for every chain length $N$ and every configuration $\sigma$.  The local
    ground spaces, the parent interactions, and the periodic chain ground
    spaces of $A$ and $B$ also coincide:
    \begin{align}
        G_L(A)              & =G_L(B), \notag             \\
        h_L(A)              & =h_L(B), \notag             \\
        \mathcal G_{N,L}(A) & =\mathcal G_{N,L}(B) \notag
    \end{align}
    for all $L$ and $N$, where $h_L$ is the canonical parent interaction and
    $\mathcal G_{N,L}$ the periodic chain ground space.

    This is the normalization taken without loss of generality in
    \cite[equations~(3.1)--(3.2b), lines~1394--1435]{Nachtergaele1996Spectral}
    and at the start of the proof of
    \cite[Theorem~4, lines~765--770]{PerezGarcia2007Matrix}.  The gauge here is
    built from the Perron eigenvector of the adjoint transfer map, giving the
    trace-preserving orientation.  That orientation is what conditions~1 and~2
    of \cite[Theorem~4, lines~752--758]{PerezGarcia2007Matrix},
    $\sum_iA^i(A^i)^\dagger=\Id$ and $\sum_i(A^i)^\dagger\Lambda A^i=\Lambda$
    with $\Lambda>0$, become after the further gauge by $\Lambda^{1/2}$:
    condition~2 turns into $\sum_i(B^i)^\dagger B^i=\Id$ and condition~1 into
    the statement that $\Lambda$ is a positive definite fixed point of the
    transfer map of $B$.
\end{theorem}

\begin{proof}
    \leanok
    \uses{thm:normal_irreducible_tensor,
        thm:normalized_normal_strongly_irreducible,
        thm:primitive_from_strongly_irred,
        thm:normality_rescaling_invariance,
        thm:local_ground_space_rescaling, thm:normal_tensor_nonzero_matrix,
        thm:tp_data_irreducible, thm:gauge_invariance_normal,
        thm:gauge_invariance, thm:local_ground_space_gauge_invariance,
        lem:eval_word_smul,
        thm:parent_hamiltonian_from_ground_space,
        lem:parent_spanning_of_local_space_eq}
    Theorem~\ref{thm:normal_irreducible_tensor} makes $A$ an irreducible
    tensor, with no normalization imposed on $A$, and
    Theorem~\ref{thm:normal_tensor_nonzero_matrix} supplies the nonzero matrix
    that the Perron--Frobenius eigenvector requires.
    Theorem~\ref{thm:tp_data_irreducible} therefore returns a
    positive definite $\sigma$ and an eigenvalue $r>0$ of the adjoint transfer
    map.  Set $\zeta=r^{-1/2}$ and $B^i=\zeta\,\sigma^{1/2}A^i\sigma^{-1/2}$;
    then $\sum_i(B^i)^\dagger B^i=\Id$, and $\zeta A$ and $B$ are gauge
    equivalent.  Rescaling preserves normality by
    Theorem~\ref{thm:normality_rescaling_invariance} and gauging preserves it
    by Theorem~\ref{thm:gauge_invariance_normal}, so $B$ is normal.  Since $B$
    is trace preserving,
    Theorem~\ref{thm:normalized_normal_strongly_irreducible} makes $B$
    strongly irreducible, and
    Theorem~\ref{thm:primitive_from_strongly_irred} then returns a positive
    definite $\rho$ with $(B,\rho)$ primitive.  Gauge equivalence
    leaves the matrix product vectors unchanged by
    Theorem~\ref{thm:gauge_invariance}, and multiplying every matrix by
    $\zeta$ multiplies every length-$N$ word by $\zeta^{N}$ by
    Lemma~\ref{lem:eval_word_smul}, hence every length-$N$ coefficient of $B$
    is $\zeta^{N}$ times that of $A$.  By
    Theorem~\ref{thm:local_ground_space_rescaling} together with gauge
    invariance of the local ground space, $G_L(A)=G_L(\zeta A)=G_L(B)$.  The
    parent interaction then agrees by
    Theorem~\ref{thm:parent_hamiltonian_from_ground_space} and the periodic
    chain ground space by Lemma~\ref{lem:parent_spanning_of_local_space_eq}.
\end{proof}

\begin{theorem}[Uniform finite-range gap for normal MPS parent Hamiltonians]
    \label{thm:normal_parent_hamiltonian_finite_range_knabe_gap}
    \lean{MPSTensor.exists_parentHamiltonianES_gap_of_isNormal}
    \leanok
    \uses{def:normal, def:l_blk_injective, def:parent_hamiltonian_es}
    Let $D>0$ and let $A$ be a normal tensor.  There exist $l>1$,
    $\epsilon_l\geq0$, an integer $m\geq l+1$, and $\delta>0$ such that $A$ is
    $l$-block injective, $\epsilon_l<1/\sqrt{l+1}$,
    \begin{align}
        \delta
         & =\frac{m\left(1-\epsilon_l\sqrt{l+1}\right)^2-l^2}{m-l},
        \label{eq:normal_finite_range_knabe_delta}
    \end{align}
    and, for every $N\geq2m$ and every
    $v\in(\ker H_{N,l+1}^{\mathrm{ES}}(A))^\perp$,
    \begin{align}
        \delta\lVert v\rVert
         & \leq\left\lVert H_{N,l+1}^{\mathrm{ES}}(A)v\right\rVert.
        \label{eq:normal_finite_range_knabe_gap}
    \end{align}
    No normalization is imposed on $A$.  The uniform gap of matrix product
    state parent Hamiltonians is asserted at
    \cite[lines~2183--2187]{Cirac2021Matrix}; the arbitrary finite-range
    coefficient (\ref{eq:normal_finite_range_knabe_delta}) is the derivation
    recorded in
    \path{docs/paper-gaps/knabe88_finite_range_coefficient.tex}.
\end{theorem}

\begin{proof}
    \leanok
    \uses{thm:normalized_primitive_gauge,
        thm:primitive_parent_hamiltonian_finite_range_knabe_gap,
        thm:parent_hamiltonian_from_ground_space,
        thm:normality_rescaling_invariance,
        lem:gauge_invariance_block_injective}
    Replace $A$ by the normalized primitive representative $B$ of
    Theorem~\ref{thm:normalized_primitive_gauge} and apply
    Theorem~\ref{thm:primitive_parent_hamiltonian_finite_range_knabe_gap} to
    $B$.  Since $G_L(A)=G_L(B)$ for every $L$,
    Theorem~\ref{thm:parent_hamiltonian_from_ground_space} makes the two
    parent Hamiltonians equal, so both the kernel and the norm bound transfer
    verbatim.  Block injectivity of $B$ at length $l$ passes to $\zeta A$ by
    Lemma~\ref{lem:gauge_invariance_block_injective} applied to the inverse
    gauge, and from $\zeta A$ to $A$ by
    Theorem~\ref{thm:normality_rescaling_invariance}.
\end{proof}

\begin{lemma}[Positive gap transfer under Loewner domination]
    \label{lem:positive_gap_transfer_loewner}
    \lean{LinearMap.IsPositive.re_inner_ge_of_norm_gap,
        LinearMap.IsPositive.norm_gap_of_le_of_ker_eq}
    \leanok
    Let $P$ and $Q$ be finite-dimensional complex endomorphisms, and let
    $\gamma\geq0$.  Suppose that $P$ is positive and
    \begin{align}
        \gamma\lVert v\rVert
         & \leq\lVert Pv\rVert
    \end{align}
    for every $v\in(\ker P)^\perp$.  Then
    \begin{align}
        \gamma\lVert v\rVert^2
         & \leq\operatorname{Re}\langle Pv,v\rangle
    \end{align}
    on $(\ker P)^\perp$.  If also $P\leq Q$ and $\ker P=\ker Q$, then
    \begin{align}
        \gamma\lVert v\rVert
         & \leq\lVert Qv\rVert
    \end{align}
    for every $v\in(\ker Q)^\perp$.
\end{lemma}

% Source: Rozmán--Molnár--Schuch, arXiv:2607.19078v1, Lemma block-obc
% (statement at main.tex lines 607--651, argument at lines 795--811).  The
% source tex is not vendored under Papers/, and references.bib has no key for
% it; cite by footnote until the paper is vendored and a bib entry is added.
\begin{theorem}[Kernels of positively weighted sums of positive operators]
    \label{thm:weighted_sum_kernel}
    \lean{LinearMap.IsPositive.apply_eq_zero_of_re_inner_eq_zero,
        WeightedPositiveKernel.re_inner_weighted_sum_apply,
        WeightedPositiveKernel.apply_weighted_sum_eq_zero_iff,
        WeightedPositiveKernel.apply_sum_eq_zero_iff,
        WeightedPositiveKernel.ker_weighted_sum_eq_iInf,
        WeightedPositiveKernel.ker_sum_eq_iInf,
        WeightedPositiveKernel.ker_weighted_sum_eq_sum}
    \leanok
    Let $(P_i)_{i\in I}$ be a finite family of positive operators on an
    inner product space.  A positive operator $P$ annihilates a vector $v$
    as soon as $\operatorname{Re}\langle Pv,v\rangle=0$, and consequently,
    for real weights $w_i>0$,
    \begin{align}
        \ker\Bigl(\sum_{i\in I}w_iP_i\Bigr)
        =\ker\Bigl(\sum_{i\in I}P_i\Bigr)
        =\bigcap_{i\in I}\ker P_i,
        \label{eq:ph_weighted_sum_kernel}
    \end{align}
    equivalently $(\sum_{i\in I}w_iP_i)v=0$ exactly when $P_iv=0$ for
    every $i\in I$.  For the empty family both sums are zero and all three
    kernels are the full space.  This is the analytic kernel step of the
    blocked-sum comparison in Lemma~\texttt{block-obc} of
    Rozmán--Molnár--Schuch\footnote{arXiv:2607.19078v1.}, in which the
    blocked parent Hamiltonian is a weighted sum $\sum_i\gamma_ih_i$ of
    positive local terms with $\gamma_i>0$.
\end{theorem}

\begin{proof}
    \leanok
    If $(\sum_iw_iP_i)v=0$, the real part of the inner product with $v$
    expands as
    \begin{align}
        \operatorname{Re}\Bigl\langle\Bigl(\sum_iw_iP_i\Bigr)v,v\Bigr\rangle
        =\sum_iw_i\operatorname{Re}\langle P_iv,v\rangle
        =0.
    \end{align}
    Each summand is nonnegative, so each vanishes, and $w_i>0$ forces
    $\operatorname{Re}\langle P_iv,v\rangle=0$.  For a positive operator
    $P$, the quadratic form is nonnegative along the real affine line
    $t\mapsto v+t\cdot(Pv)$ and vanishes at $t=0$; its linear coefficient
    is $2\lVert Pv\rVert^2$, so evaluating at small negative $t$ forces
    $\lVert Pv\rVert=0$, that is $Pv=0$.  The converse inclusion holds
    term by term, and the unweighted comparison is the case $w_i=1$.
\end{proof}

\begin{theorem}[Loewner domination of positively weighted sums]
    \label{thm:weighted_sum_loewner}
    \lean{WeightedPositiveKernel.isPositive_weighted_sum,
        WeightedPositiveKernel.weighted_sum_le_sum}
    \leanok
    Let $(P_i)_{i\in I}$ be a finite family of positive operators on an
    inner product space, with real weights $(w_i)_{i\in I}$.  If
    $w_i\geq0$ for every $i$, then $\sum_{i\in I}w_iP_i$ is positive.  If
    $w_i\leq1$ for every $i$, then
    \begin{align}
        \sum_{i\in I}w_iP_i
         & \leq\sum_{i\in I}P_i
    \end{align}
    in Loewner order.  This is the comparison step of the blocked-sum
    argument in the source cited in
    Theorem~\ref{thm:weighted_sum_kernel}, where the coefficients satisfy
    $0<\gamma_i\leq1$.
\end{theorem}

\begin{proof}
    \leanok
    Positivity of the weighted sum follows from positivity of each scalar
    multiple $w_iP_i$ and of finite sums of positive operators.  For the
    domination, the difference is again a weighted sum,
    \begin{align}
        \sum_{i\in I}P_i-\sum_{i\in I}w_iP_i
        =\sum_{i\in I}(1-w_i)P_i,
    \end{align}
    now with nonnegative weights $1-w_i\geq0$, hence positive.
\end{proof}

\begin{theorem}[Gap transfer through a positively weighted sum]
    \label{thm:weighted_sum_gap_transfer}
    \lean{WeightedPositiveKernel.norm_gap_sum_of_weighted_sum}
    \leanok
    Let $(P_i)_{i\in I}$ be a finite family of positive operators on a
    finite-dimensional complex inner product space, with real weights
    $0<w_i\leq1$, and let $\gamma\geq0$.  If
    \begin{align}
        \gamma\lVert v\rVert
         & \leq\Bigl\lVert\Bigl(\sum_{i\in I}w_iP_i\Bigr)v\Bigr\rVert
    \end{align}
    for every $v\in\bigl(\ker\sum_{i\in I}w_iP_i\bigr)^\perp$, then
    \begin{align}
        \gamma\lVert v\rVert
         & \leq\Bigl\lVert\Bigl(\sum_{i\in I}P_i\Bigr)v\Bigr\rVert
    \end{align}
    for every $v\in\bigl(\ker\sum_{i\in I}P_i\bigr)^\perp$.
\end{theorem}

\begin{proof}
    \leanok
    \uses{thm:weighted_sum_kernel, thm:weighted_sum_loewner, lem:positive_gap_transfer_loewner}
    The weighted sum is positive and, by
    Theorem~\ref{thm:weighted_sum_loewner}, dominated by the unweighted
    sum; by the kernel equality~(\ref{eq:ph_weighted_sum_kernel}) of
    Theorem~\ref{thm:weighted_sum_kernel}, the two sums have the same
    kernel.  Lemma~\ref{lem:positive_gap_transfer_loewner} transfers the
    norm gap from the smaller weighted sum to the larger unweighted sum.
\end{proof}

\begin{theorem}[Uniform gap for normal MPS parent Hamiltonians]
    \label{thm:normal_mps_parent_hamiltonian_gap}
    \notready
    \discussion{5926}
    \uses{thm:primitive_parent_hamiltonian_finite_range_knabe_gap, thm:anticommutator_gap_bound, thm:parent_hamiltonian_gapped}
    Let $A$ be a normal tensor.  There exist an interaction range $L>1$ and
    a constant $\gamma>0$ such that, for every $N\ge2L$ and every
    $v\in(\ker H_{N,L}(A))^\perp$,
    \begin{align}
        \gamma\lVert v\rVert
         & \le \lVert H_{N,L}(A)v\rVert.
        \notag
    \end{align}
    Theorem~\ref{thm:primitive_parent_hamiltonian_finite_range_knabe_gap}
    proves the estimate for a primitive tensor with a positive-definite fixed
    point, at range $L=l+1$, with no divisibility restriction on the chain
    length, but only for $N\ge2m$, where the Knabe volume $m$ satisfies
    $m\ge l+1$ and is chosen large enough that $l^2<m\gamma_{\mathrm{open}}$.
    Two steps therefore remain.  The chain lengths with $2(l+1)\le N<2m$ are
    not covered by that theorem, so a finite-volume estimate on this finite
    set of lengths, together with a common positive constant, is still needed.
    The passage from a primitive tensor with a positive-definite fixed point
    to an arbitrary normal tensor is also still needed.
\end{theorem}

\begin{proof}
    \notready
    \discussion{5926}
    \uses{thm:primitive_parent_hamiltonian_finite_range_knabe_gap, thm:anticommutator_gap_bound, thm:parent_hamiltonian_gapped}
    The primitive case with a positive-definite fixed point, at the lengths
    $N\ge2m$, is
    Theorem~\ref{thm:primitive_parent_hamiltonian_finite_range_knabe_gap}.  It
    remains to bound the gap from below on the finitely many lengths
    $2(l+1)\le N<2m$ and to take the minimum of the two constants, and then to
    reduce an arbitrary normal tensor to the primitive case while preserving
    the ground-space complement and a uniform positive constant.
\end{proof}

\begin{theorem}[Gap bound from a cyclic anticommutator estimate]
    \label{thm:anticommutator_gap_bound}
    \leanok
    \uses{lem:parent_hamiltonian_cyclic_overlap_anticommutator_gap}
    Let $L>1$. Suppose that for every $N\ge2L$ and every overlapping
    off-diagonal pair of cyclic length-$L$ windows,
    \begin{align}
        h_{i,\mathrm{ES}}h_{j,\mathrm{ES}}
        +h_{j,\mathrm{ES}}h_{i,\mathrm{ES}}
         & \ge
        -\left(1-\frac{1}{4L}\right)\frac{1}{2(L-1)}
        (h_{i,\mathrm{ES}}+h_{j,\mathrm{ES}})
        \notag
    \end{align}
    in quadratic form. Then the parent Hamiltonian has the explicit gap
    constant $\gamma_L:=1/(4L)>0$.
    Concretely, if $N\ge2L$ and
    $v\in(\mathcal G_{N,L}^{\mathrm{ES}}(A))^\perp$, then
    $\gamma_L\|v\| \le \|H_{N,L}^{\mathrm{ES}}(A)v\|$.
\end{theorem}

\begin{proof}
    \leanok
    The finite-overlap martingale reduction applies with the cyclic overlap
    relation: at most $2(L-1)$ off-diagonal neighbours per row; for
    non-overlapping pairs the disjoint-window estimate
    $\re\langle h_{i,\mathrm{ES}}v,h_{j,\mathrm{ES}}v\rangle\ge 0$ holds;
    and the anticommutator inequality is retained as an explicit hypothesis.
\end{proof}

\begin{theorem}[Conditional spectral gap for MPS parent Hamiltonians]
    \label{thm:parent_hamiltonian_gapped}
    \lean{MPSTensor.parentHamiltonian_gapped_of_anticommutator}
    \leanok
    \uses{thm:anticommutator_gap_bound}
    For a tensor $A$ and a range $L>1$, assume the cyclic-window
    anticommutator estimate of Theorem~\ref{thm:anticommutator_gap_bound}.
    Then there exists $\gamma>0$
    such that for every $N\ge2L$ and every vector
    $v\in(\mathcal G_{N,L}^{\mathrm{ES}}(A))^\perp$ one has
    $\gamma\|v\| \le \|H_{N,L}^{\mathrm{ES}}(A)v\|$.
    By Lemma~\ref{lem:quadratic_form_to_norm_bound}, this implies the usual
    uniform lower bound on every nonzero eigenvalue of the parent Hamiltonian.
\end{theorem}

\begin{proof}
    \leanok
    This follows immediately from the corresponding explicit theorem by taking
    $\gamma$ to be the constant produced there.
\end{proof}
